% !TeX program = xelatex
% Build: xelatex coveo_session_recommender_math.tex (until references stabilize).
% Requires XeLaTeX, polyglossia, and standard Windows fonts: Times New Roman, Arial, and Consolas.
% Standalone file; no external .bib files, images, or included files are required.
\documentclass[12pt,a4paper]{article}
\usepackage{fontspec}
\setmainfont{Times New Roman}
\setsansfont{Arial}
\setmonofont{Consolas}[Scale=MatchLowercase]
\newfontfamily\cyrillicfont{Times New Roman}[Script=Default]
\newfontfamily\cyrillicfontsf{Arial}[Script=Default]
\newfontfamily\cyrillicfonttt{Consolas}[Script=Default,Scale=MatchLowercase]
\usepackage{polyglossia}
\setdefaultlanguage{english}
\usepackage{amsmath,amssymb,amsthm,mathtools}
\usepackage{array,booktabs,longtable,geometry,enumitem}
\usepackage{microtype,xurl,hyperref,ragged2e}
\AtBeginDocument{\RaggedRight}
\geometry{margin=2.2cm}
\hypersetup{hidelinks,pdftitle={Coveo Recommender System: Mathematical Specification},pdflang={en}}
\setlength{\parindent}{0pt}
\setlength{\parskip}{0.32\baselineskip}
\setlength{\emergencystretch}{3em}
\setlist{nosep,leftmargin=2em}
\setcounter{tocdepth}{2}
\setcounter{secnumdepth}{3}
\allowdisplaybreaks[2]
\newcolumntype{L}[1]{>{\raggedright\arraybackslash}p{#1}}
\newcommand{\code}[1]{{\ttfamily\small\nolinkurl{#1}}}
\newcommand{\R}{\mathbb R}
\newcommand{\N}{\mathbb N}
\newcommand{\E}{\mathbb E}
\newcommand{\Prob}{\mathbb P}
\newcommand{\ind}{\mathbf 1}
\newcommand{\eps}{\varepsilon}
\newcommand{\norm}[1]{\left\lVert#1\right\rVert}
\newcommand{\abs}[1]{\left|#1\right|}
\newcommand{\pos}[1]{\left(#1\right)_+}
\newcommand{\given}{\,\vert\,}
\DeclareMathOperator*{\argmin}{arg\,min}
\DeclareMathOperator*{\argmax}{arg\,max}
\DeclareMathOperator{\Var}{Var}
\DeclareMathOperator{\Cov}{Cov}
\DeclareMathOperator{\sigmoid}{sigmoid}
\DeclareMathOperator{\Top}{Top}
\DeclareMathOperator{\clip}{clip}
\newtheoremstyle{plainru}{0.45\baselineskip}{0.45\baselineskip}{\normalfont}{}{\bfseries}{.}{0.5em}{}
\theoremstyle{plainru}
\newtheorem{definition}{Definition}[section]
\newtheorem{assumption}{Assumption}[section]
\newtheorem{proposition}{Proposition}[section]
\newtheorem{invariant}{Invariant}[section]
\newtheorem{testcase}{Test}[section]
\newtheorem{remark}{Remark}[section]
\makeatletter
\renewcommand*\l@section[2]{%
  \ifnum \c@tocdepth >\z@
    \addpenalty\@secpenalty\addvspace{0.6em \@plus\p@}%
    \setlength\@tempdima{2.8em}%
    \begingroup\parindent\z@\rightskip\@pnumwidth
      \parfillskip-\@pnumwidth\leavevmode\bfseries
      \advance\leftskip\@tempdima\hskip-\leftskip
      #1\nobreak\hfil\nobreak\hb@xt@\@pnumwidth{\hss#2}\par
    \endgroup\fi}
\renewcommand*\l@subsection{\@dottedtocline{2}{2.8em}{3.4em}}
\makeatother
\begin{document}
\begin{center}
{\Large\bfseries An E-Commerce Recommender System:\\[0.3em]
Mathematical Specification Using Coveo SIGIR eCommerce Data\par}
\vspace{0.7em}
{\large Interaction Graph, Item Features,\par
and Session Search Context\par}
\vspace{0.7em}
{Version 2.10.\par}
{Demo service requirements clarified on October 8, 2026.\par}
\end{center}
\begin{abstract}
The initial version uses a transition graph, provided description
and image vectors, search, and LambdaMART ranking. A neural extension
learns item representations and session intents.
Each stage specifies its inputs, training procedure, validation,
and conditions for proceeding to the next stage.
The business objective is to increase the number of sessions with a purchase per 1000 eligible sessions.
The first Coveo evaluation measures whether the first recommendation list includes a recorded future purchase
among sessions with at least one such purchase. Telemetry completeness
remains unknown; evaluation over all sessions requires it to be confirmed.
The data support measuring agreement with historical purchases;
causal effects require a randomized experiment with recommendation impressions.
This specification contains no training results or load measurements.
\end{abstract}

\section*{Implementation sequence}
\phantomsection\label{sec:reading-route}
\label{sec:first-run}\label{sec:stage-acceptance}
\addcontentsline{toc}{section}{Implementation sequence}

\begin{longtable}{L{0.22\textwidth}L{0.69\textwidth}}
\toprule Stage & Sections and outcome \\\midrule\endhead
Baseline system & Part~\ref{part:core}, recipe~\ref{app:first-config},
and tests~\ref{app:tests-base}. The first report uses
$C_{\mathrm{smoke}}$; no configuration tuning or access to $D$. \\
Selection and final evaluation & Period and selection sequence:
Section~\ref{sec:eval-primary-comparison}; result states:
Section~\ref{sec:acceptance}. \\
Neural model & Part~\ref{part:neural}, parameters~\ref{app:neural-config},
tests~\ref{app:tests-neural}. Requires a reproducible baseline,
diagnostics of its losses, and measured cost. \\
Separate studies & Sequence~\ref{sec:sequence-control},
graph updates~\ref{sec:freshness}, SID~\ref{sec:sid-overview},
other models~\ref{sec:model-extensions}. Each requires a predefined
question, control, and budget. \\
Experiment with impressions & Protocol~\ref{sec:eval-online},
after serving validation. Measures the causal effect. \\
\bottomrule
\end{longtable}

Conditions for moving between stages are specified in the
run manifest~\ref{app:run-stages}. Parameter values are in
Appendix~\ref{app:parameters}; report contents are in
\ref{app:diagnostic-output}. If $D$ has already been used to evaluate the baseline,
the next extension requires a new final evaluation period.

The current implementation goal is local historical evaluation through $D$
under the \code{observed_purchases_v1} protocol
(Section~\ref{sec:eval-observed}) in \code{ideal_deterministic} mode.
It defines common rules for recorded history, targets, metrics, and acceptance.
Estimating purchase frequency over all sessions would require confirmed
observation completeness and a separate formulation; it is outside the current scope.
Service readiness is checked
separately and is not a prerequisite for this evaluation. Neural studies,
the service, and A/B testing remain later stages subject to all requirements
of their respective sections.

\clearpage
\begingroup
\small
\setlength{\parskip}{0.1\baselineskip}
\tableofcontents
\endgroup
\clearpage

\part{Core system: data, training, and evaluation}\label{part:core}
\section{Purpose and chosen formulation}\label{sec:scope}

Recommendations use the current session without a persistent
customer identifier. Inputs are the available history, the latest search
query, and the item feature version. The response contains an ordered list,
the model version, candidate sources, and computation status.

The first implementation is policy \(\pi_{\mathrm{base}}\): graph
transitions, similarity of the provided item vectors, search context,
and LambdaMART ranking. This policy is called GCS+Rk:
G denotes the graph, C the description and image features,
S search, and Rk the ranker.
The first run is defined in Part~I, neural representations in
Part~II, and additional studies in the appendices.

Initial settings and load targets define requirements;
performance and business impact remain to be measured.

\begin{longtable}{L{0.28\textwidth}L{0.64\textwidth}}
\toprule Object & Primary variant \\\midrule\endhead
Deployment objective & Increase the fraction of eligible sessions with a purchase, subject to
latency, error, and cost constraints. \\
Outcome unit & A session admitted under rules checked before the first
experimental recommendation impression. \\
Item retrieval in the first run & One-step transitions on the historical
graph, content and search candidates, recent items, and popularity. \\
Ranking & Trees with pairwise training; the primary label is a subsequent
purchase of the item. Clicks and cart actions are considered separately. \\
Historical evaluation & Common protocol, sessions, and requests across methods;
the first Coveo evaluation follows Section~\ref{sec:eval-observed}. \\
Evidence of impact & A randomized experiment counting all
assigned sessions, including serving failures. \\
Prototype subsample & Complete sessions from a predefined set;
the result is not claimed as an official competition result. \\
\bottomrule
\end{longtable}

\begin{invariant}[Three distinct quantities]
Model scores, agreement of recommendations with historical purchases, and changes
in conversion after deployment are different quantities. They are not equated
without additional data and assumptions.
\end{invariant}

\section{Notation and units}\label{sec:symbols}

\(\ind\{A\}\) is the indicator of condition \(A\), \(\E\) denotes
expectation, and \(\Prob\) probability. \(\Top_m\) returns at most
\(m\) elements with the highest scores; ties are resolved
by lexicographic identifier order. Time is measured in seconds.

\small
\begin{longtable}{L{0.23\textwidth}L{0.69\textwidth}}
\toprule Symbol & Meaning \\\midrule\endhead
\(s,i,j\) & Session and item identifiers. \\
\(e,t(e),a(e)\) & An event, its time, and the time it becomes available to the system. \\
\(q,s(q),t_q\) & A recommendation request, its session, and its time. \\
\(t_{0,s},H,g\) & First eligible request, outcome horizon, and inactivity
gap, equal to 30 minutes for Coveo sessions. \\
\(\mathcal H_q,\mathcal F_q\) & Available history and all information
available to the algorithm at request time. \\
\(r_G,r_N,r(q),\mathcal I_q\) & Times of the graph snapshot, publication
of the learned retriever, the snapshot used, and the eligible catalog. \\
\(\mathcal Q_s,Q_{\max}\) & Session requests and their predefined limit. \\
\(\pi,\pi_0\) & Complete recommendation policy and comparison policy. \\
\(\pi_{\mathrm{base}}\) & The first implementation, GCS+Rk. \\
\(\mathcal C_q(\pi),M\) & Candidates and maximum candidate count. \\
\(L_q(\pi),K\) & Final list and its length limit. \\
\(Y_q,y_{qi}\) & Set of future purchases after a request and binary label
of candidate \(i\). \\
\(Y_s(\pi)\) & Indicator of a purchase in a session when policy
\(\pi\) is actually used. Distinct from the set \(Y_q\). \\
\(B(\pi),\Delta B\) & Purchasing sessions per 1000 eligible sessions and their causal increase. \\
\(\operatorname{Hit}_s(\pi),\operatorname{Hit@20}_J\) & A first-list hit
on recorded future purchases and its rate among sessions with such a target. \\
\(f_q(i)\) & Ranking score, not a purchase probability. \\
\(Z_s\) & Assigned experiment arm for impressions, 0 or 1. \\
\(\xi,\widehat\xi\) & Procedure configuration and selected configuration. \\
\(A_G,A_N,B,C,E,D\) & Blocks for topology and statistics, neural retriever
training, ranker training, selection, reserve, and final evaluation;
\(A=A_G\cup A_N\). \\
\bottomrule
\end{longtable}
\normalsize

The index \(r\) on an aggregate restricts events to those available at the snapshot.
Adding a score to a monetary value requires a justified conversion
coefficient. Comparison policy \(\pi_0\) is fixed in advance;
it need not coincide with \(\pi_{\mathrm{base}}\).
Block \(A_N\) is reserved for the neural model and does not augment the graph
or item statistics in the first run.

\section{Coveo data and observation limits}\label{sec:data}

\subsection{Observed objects}

The release contains tables of behavior, search events, and item
features~\cite{coveo-data,coveo-paper}. The formulas below use
the following event representation:
\begin{equation}\label{eq:data-event}
e=(s,i,t,a,\kappa,\eta),\qquad
\kappa\in\{\mathrm{detail},\mathrm{add},\mathrm{remove},
\mathrm{purchase},\mathrm{search},\mathrm{page}\}.
\end{equation}
Here \(i=\varnothing\) is allowed for a page without an item and for search;
\(\eta\) contains additional fields. Time \(t\) is obtained by dividing
\code{server_timestamp_epoch_ms} by 1000. Ingestion times \(a\)
are absent from the original dataset: in historical replay, the primary model
assumes \(a=t\), while delays are examined in a separate test.

Search is represented as
\begin{equation}\label{eq:data-search}
v=(s,t,u,R,C),
\end{equation}
where \(u\) is \code{query_vector}, \(R\) is the list
\code{product_skus_hash}, and \(C\) is \code{clicked_skus_hash}.
The item content vector \(d_i\) comes from
\code{description_vector}, and the image vector \(v_i\) from
\code{image_vector}. Category \(c_i\) is a hashed hierarchy,
and \(b_i^{\mathrm{price}}\) is an ordinal price bucket.

\subsection{What cannot be recovered from the dataset}

\begin{longtable}{L{0.32\textwidth}L{0.60\textwidth}}
\toprule Missing object & Consequence for the formulation \\\midrule\endhead
Actual order price, quantity, and margin & Monetary
revenue and profit cannot be computed from a price bucket number. \\
Impression log for our recommendation widget & Agreement with a purchase
does not establish the recommendation's influence. \\
Recommendation assignment probabilities & The primary historical experiment
does not use causal reweighting of impressions. \\
Historical inventory and item activation dates & Catalog eligibility
is defined by an observation rule, not by stock availability. \\
Time of each click within a search row & Field \(C\) cannot be used
as a feature of the current or next request without a separate timestamped event. \\
Initial cart contents and quantities & Cart state denotes
the last observed action, not an exact balance. \\
Historical content encoder versions & Use of the provided
vectors is conditional on the released feature snapshot. \\
Confirmation of telemetry completeness & An empty set of recorded purchases
does not prove that no purchase occurred. Evaluation over positively observed
targets is defined in Section~\ref{sec:eval-observed}. \\
\bottomrule
\end{longtable}

\begin{assumption}[Availability of static features]\label{ass:static-content}
In historical evaluation, the provided vectors and categories
are assumed available for items in the historical catalog. This is an assumption
about side information, not evidence of its actual availability
in the past. A mandatory comparison without the provided metadata tests
sensitivity to this assumption (Section~\ref{sec:metadata-sensitivity}).
It does not recover the unknown historical version of the vectors.
\end{assumption}

Price buckets are encoded as categories or compared by order
and equality; their differences and ratios have no monetary meaning.
Hashed categories are compared by equality and their known hierarchy.

\subsection{Snapshot and reproducible subsample}

A data snapshot is defined by the tuple
\begin{equation}\label{eq:data-snapshot}
\mathsf D=(\mathrm{version},\mathrm{hashes},\mathrm{schema},
T_{\min},T_{\max},\mathrm{sampling},\mathrm{cleaning}).
\end{equation}
Within predefined time blocks, sessions are sorted
by a stable pseudorandom key \(h_{\mathrm{seed}}(s)\).
The initial algorithm is SHA-256 over the concatenation of the seed as an unsigned
64-bit big-endian integer, the ID byte length as an unsigned 64-bit big-endian integer,
and the exact UTF-8 bytes of the ID, without case conversion or Unicode normalization.
Sorting is lexicographic over the 32 digest bytes, with ties resolved
by those same ID bytes; the seed is 17. Size $n$ is set separately
for each period before inspecting its outcomes.
The first \(n\) identifiers are taken with all behavior and search rows,
regardless of future purchases. This produces nested training samples,
such as 200 and 500 thousand sessions.

The published counts include sessions, events, items, purchasing sessions, and search
queries, as well as the fractions of missing vectors and unknown test items.
Purchasing sessions are counted by identifier, not by \code{purchase}
rows; catalog size is reported separately from sample size.

\section{Time, cleaning, and available history}\label{sec:time}

\subsection{Event ordering}

Behavior and search are merged by \(s\) and time; simultaneous events
form a block. Row numbers define storage order, not causal order.
Under historical \(a=t\), a block contains all rows at that time;
in the live service, its completeness is defined by the closed envelope
in Section~\ref{sec:ops-transport-block}.

A single item view may generate both \code{detail} and
\code{pageview}. The primary item sequence uses
\code{product_action}; a page without an action does not create a second
item transition. All source rows are retained for audit.
Without an event identifier, exact row equality alone
does not prove a technical duplicate. The primary variant uses the declared
canonicalization \code{event_semantics=set_v1}: each combination
\((s,t,i,\kappa)\) maps to one canonical event, retaining the
multiplicity of its source rows. For search, its payload additionally
preserves the set of distinct \code{query_vector} values and the multiplicity
of each, rather than the first or last value.
\label{sec:data-search-canonicalization}
Valid vectors are compared before normalization: by dimension and exact
FP64 coordinates, with $-0$ and $+0$ treated as equal. Missing and
invalid fields form separate markers with reason codes;
source rows are retained.
Fields $R,C$ from~\eqref{eq:data-search} do not participate in this representation.
A search block is unambiguous only if it contains one valid value
and no markers. Two or more distinct valid vectors set
an ambiguity flag; a marker also disables this context
with the corresponding reason. Repeated copies of a vector do not create
ambiguity. Row order does not select one of the queries.
Below, $\mathcal H_q$, view counters,
history $\mathcal D_q$, and the sequence encoder use these
canonical events. This is a representation assumption, not a reconstruction
of the true action count; source rows and their multiplicities are not removed.

A sensitivity check on $C_{\mathrm{smoke}}$ is mandatory for
\code{event_semantics=multiset_v1}: multiplicity is included in counters,
history, and averaging within time groups. Binary purchase targets remain
sets, and edge support still counts distinct sessions.
The report includes the fraction of affected sessions and changes in features and coverage.
The choice of semantics is fixed before $C_{\mathrm{select}}$ and is shared by
the models being compared. Transport retries in the live service are deduplicated by
stable event ID before this operation; matching fields are not a substitute.

\begin{invariant}[Simultaneity]
No directed transitions are constructed between events in the same time block,
and no targets are assigned as ``after the request.'' The request runs
after the entire available block has been processed; its future targets have
strictly greater times.
\end{invariant}

\subsection{Information at decision time}

Let \(a_q\) be the computation start time. Then
\begin{equation}\label{eq:time-history}
\mathcal H_q=\{e:s(e)=s(q),\ t(e)\le t_q,\ a(e)\le a_q\}.
\end{equation}
\(\mathcal F_q\) includes \(\mathcal H_q\), available reference-data
snapshots, published model parameters, and service state.
Serving and independent application of the ranker on $C,E,D$ require
\begin{equation}\label{eq:time-adapted}
\mathcal C_q,\ f_q,\ L_q\quad\text{to be functions only of }\mathcal F_q.
\end{equation}
Future session length, final conversion, the next search, and future
popularity are prohibited. Snapshot \(r(q)<t_q\) is selected by its
publication time. In $B$, the ranker uses historically eligible features
and candidates, but training predictions are not an independent evaluation.
The neural model trains on continuations of examples in \(A_N\) and is evaluated
on subsequent independent periods.

The initial and neural variants use only \(u\), without result list \(R\),
which depends on the unknown search engine policy.
Reordering the returned \(R\) is a separate task;
\(C\) remains a subsequent response. Its results are not combined
with recommendation results.

\subsection{Observed cart state}

Define \(b_{qi}\in\{\mathrm{add},\mathrm{remove},\mathrm{unknown}\}\).
Its initial value is \(\mathrm{unknown}\). The latest unambiguous
observed \code{add} sets \(\mathrm{add}\), and \code{remove}
sets \(\mathrm{remove}\). If the latest time group contains
both actions, the value becomes \(\mathrm{unknown}\) again.
A repeated removal does not create a negative item quantity.

The model uses indicators of the three states and time since the latest
unambiguous event. Actual cart contents are not reconstructed.

\section{Requests, eligible items, and future targets}\label{sec:requests}

\subsection{Observation start and schedule}

A session is eligible if it contains a \code{detail}
with a valid item identifier and no purchase has been observed by the end
of that time block. The first such block defines \(t_{0,s}\).
The rule does not require a future purchase, multiple views,
a content vector, or a nonempty response from the selected algorithm.
Sessions that start with a purchase before the first eligible view
are outside the target population; their fraction is published.

The sole request $q_s^{(1)}$ runs at \(t_{0,s}\):
$\mathcal Q_s=\{q_s^{(1)}\}$, $Q_{\max}=1$ for all methods.
Its prefix includes the entire block of simultaneous events.
A future purchase does not select the request time. Subsequent lists
are not evaluated in this experiment. The latest search block is checked
for ambiguity as specified in Section~\ref{sec:data-search-canonicalization}.

\begin{remark}
The historical schedule \(\mathcal Q_s\) follows the recorded path.
Recommendations may change future actions and requests; log replay
does not model this change.
\end{remark}

\subsection{Recorded target window}

Let $t_{\mathrm{last},s}^{\mathrm{obs}}$ be the time of the last recorded
event in the session's merged behavior and search stream. The target window is
\begin{equation}\label{eq:request-horizon}
\tau_s=\min(t_{0,s}+H,t_{\mathrm{last},s}^{\mathrm{obs}}).
\end{equation}
The last record and a wait of $g=1800$ seconds do not confirm session
closure or telemetry completeness. The calendar buffer $H+g$ in
Section~\ref{sec:eval-maturity} protects period boundaries without proving
the absence of missing purchases. The first version uses horizon $H=86400$ seconds.

\subsection{Catalog and labels}

In the primary historical variant, the catalog is fixed using $A_G$:
\begin{align}
\mathcal I_G&=\{i:\exists e\text{ from eligible sessions in }A_G,\ i(e)=i\},\notag\\
\mathcal I_q&=\mathcal I_G\setminus\mathcal P_q.
\label{eq:request-catalog}
\end{align}
Here \(\mathcal P_q\) denotes items already purchased in the visible prefix.
The first request is admitted before any observed purchase,
so $\mathcal P_q=\varnothing$. The exclusion remains
a required safeguard when reusing the service. Recommendations
after a purchase, repeat purchases of consumables, and item quantities
require a separate task with a newly specified schedule,
catalog, and labels; the binary target of the current experiment does not model quantity.
Metadata do not confirm the historical availability of future items.
A prefix item unknown to the snapshot receives the corresponding flag;
both primary and fallback recommendations stay within \(\mathcal I_q\).

In a real store, set~\eqref{eq:request-catalog} is replaced
by the catalog available at request time, subject to stock,
region, and publication restrictions. The quality of such filtering is not measured on Coveo.
Items with an observed \code{add} are allowed in the primary variant;
recommendations restricted to items new to the cart are evaluated
separately with the same denominator.

For $s=s(q)$, the recorded target set and candidate labels are
\begin{align}
Y_q=Y_s^{\mathrm{obs}}&=\{i:\exists e\text{ in the export},\ s(e)=s,
\ \kappa(e)=\mathrm{purchase},\ t_{0,s}<t(e)\le\tau_s,\ i(e)=i\},
\label{eq:request-target}\\
y_{qi}&=\ind\{i\in Y_q\}.\label{eq:request-label}
\end{align}
The target contains every distinct purchased SKU in the window, not only the first
or last. Repeated purchase records contribute one set element;
multiple items in the same time block are retained.
Items in \(Y_q\) that are absent from \(\mathcal I_q\) are not removed
from evaluation. This separates catalog coverage error from ranking error.
An empty $Y_q$ means no target purchase was recorded, not a confirmed
non-purchase. Even when $Y_q\ne\varnothing$, completeness of the purchase list is unknown.
A zero candidate label means only that the item is absent from $Y_q$;
its use is governed by Section~\ref{sec:model-ranking}.

\section{Primary business criterion}\label{sec:business}

\subsection{Purchasing sessions and causal lift}

Policy \(\pi\) includes the model, recommendation rules, fallback behavior,
and request-count limits. Let \(Y_s(\pi)\) denote
the indicator of at least one purchase after policy assignment
and before the end of its outcome window. A policy change
may alter the session path and its end. The outcome is defined
by the same window rule for every possible path.

\begin{definition}[Policy value]
For a predefined population of eligible sessions,
\begin{equation}\label{eq:business-value}
B(\pi)=1000\,\E[Y_s(\pi)],\qquad
\Delta B(\pi,\pi_0)=1000\,\E[Y_s(\pi)-Y_s(\pi_0)].
\end{equation}
The unit of \(\Delta B\) is additional sessions with a purchase
per 1000 eligible sessions.
\end{definition}

Below, ``eligible'' means sessions admitted under the rule in
Section~\ref{sec:requests}. Let $\mathcal S_{\mathrm{all}}$
be all recorded sessions from the same source and calendar window,
and $\mathcal S_{\mathrm{elig}}$ those that pass this rule. Report separately
$r_{\mathrm{elig}}=|\mathcal S_{\mathrm{elig}}|/|\mathcal S_{\mathrm{all}}|$
and exclusion counts by reason; outcome maturity is assessed in the next,
separate step. With an empty denominator, the fraction is unavailable.
Formula~\eqref{eq:business-value} is conditional on $s\in\mathcal S_{\mathrm{elig}}$.
It does not measure conversion across all visits. Multiplying by $r_{\mathrm{elig}}$
alone does not yield the overall effect: this requires unchanged eligibility, no
effect outside the eligible population, and a common source population. A KPI for all visits
requires a separate experiment with assignment before selection and its own
denominator. When assignment occurs only after admission, the admission rate cannot
be broken down by an arm that has not yet been assigned.

In A/B testing with a stable key, the primary cohort is further restricted
to entry sessions $\mathcal S_{\mathrm{entry}}\subseteq\mathcal S_{\mathrm{elig}}$
under~\ref{sec:eval-itt}. Then~\eqref{eq:business-value} and the business threshold
explicitly use $B_{\mathrm{entry}}(\pi)=1000\E[Y_s(\pi)\mid
s\in\mathcal S_{\mathrm{entry}}]$, and the fraction of this cohort
among eligible sessions is also reported. Confirmation of $\Delta B_{\mathrm{entry}}$
does not extend to all $\mathcal S_{\mathrm{elig}}$; a product needing such
generalization requires a separate estimand and protocol. The notation
$\Delta B$ in acceptance criteria always specifies the chosen population.

Any purchase in an eligible session counts as one success regardless of the number
of items or a recommendation click. Selection by such a click would depend
on behavior after arm assignment.

For example, increasing the fraction of purchasing sessions from 0.030 to 0.032
corresponds to \(\Delta B=2\), an absolute increase of 0.2 percentage
points, and a relative increase of about 6.7\%. This is an arithmetic example,
not an effect estimate for Coveo or the system being developed.

The deployment configuration is selected by solving
\begin{equation}\label{eq:business-optimization}
\max_{\pi\in\Pi}\ \Delta B(\pi,\pi_0)
\quad\text{subject to}\quad
\begin{cases}
\Prob(\ell_q>L_{\max})\le\rho_L,\\
\Prob(\mathrm{error}_q=1)\le\rho_{\mathrm{err}},\\
\E[C_s^{\mathrm{serve}}]\le c_{\max},\\
L_q\text{ satisfies item constraints}.
\end{cases}
\end{equation}
Here \(\ell_q\) is response latency, and \(C_s^{\mathrm{serve}}\) is
the cost of serving a session in the chosen currency unit.
The product owner sets the constraints before comparing systems.
When monetary rates are unavailable, a measurable computational budget
is fixed instead of cost; no monetary savings are claimed.

\subsection{Why a purchase label does not solve the causal problem}

A historical log contains one outcome for each session,
not both \(Y_s(\pi)\) and \(Y_s(\pi_0)\). Predicting a future purchase
can select a model for testing but does not identify
the difference~\eqref{eq:business-value}.

\begin{proposition}[Effects are not identifiable from matches]\label{prop:business-nonidentification}
For $0<p_0=\E[Y_s(\pi_0)]<1$, the same distribution of observed purchases
and the same recommendation matches are compatible with zero,
positive, or negative effects of a new widget.
\end{proposition}
\begin{proof}
Fix the outcome distribution under the original policy. Three models with
$Y_s(\pi)=Y_s(\pi_0)$, $Y_s(\pi)=1$, and $Y_s(\pi)=0$ produce effects
$0$, $1000(1-p_0)>0$, and $-1000p_0<0$, respectively, but agree on the original
log because the new policy was never assigned there.
In general, binary outcomes bound the effect to the interval
$[-1000p_0,\,1000(1-p_0)]$. At $p_0=0$ or $p_0=1$, not every
sign is possible, but the effect magnitude remains unidentified.
The absence of one class in a finite log does not prove
the corresponding boundary value of $p_0$.
\end{proof}

\subsection{Monetary extension when the required data become available}

If a real system provides paid order lines, returns,
cost of goods, and additional expenses, one can define
\begin{equation}\label{eq:business-margin}
G_s(\pi)=\sum_{o\in\mathcal O_s(\pi)}
\bigl(R_o-C_o^{\mathrm{goods}}-C_o^{\mathrm{delivery}}
-C_o^{\mathrm{return}}\bigr)-C_s^{\mathrm{serve}},
\end{equation}
where \(\mathcal O_s\) denotes orders attributed to the session under a predefined
rule. All terms must use the same currency unit
and consistent treatment of taxes and returns.
Then \(1000\E[G_s(\pi)-G_s(\pi_0)]\) is the additional
contribution margin per 1000 sessions. The return horizon is set separately.
In Coveo, this functional has status \code{not_observed}.

A price bucket cannot substitute for \(R_o\). Converting \(\Delta B\) into money
requires assumptions about order counts, baskets, and margins, which may
change after deployment. The primary business criterion is \(\Delta B\).

\subsection{Levels of evidence}

\begin{longtable}{L{0.26\textwidth}L{0.32\textwidth}L{0.32\textwidth}}
\toprule Level & Supported conclusion & Required data \\\midrule\endhead
Historical suitability & The system finds future purchases in a fixed
log and stays within budget. & Coveo and a reproducible protocol. \\
Operational suitability & Requests are served within the required latency
using correct features. & Load tests and a shadow deployment. \\
Business superiority & \(\Delta B\) exceeds a predefined practical threshold. &
Randomized assignment and mature outcomes. \\
Monetary benefit & Additional margin is positive. & An experiment,
financial data, and completed accounting for returns. \\
\bottomrule
\end{longtable}

\section{Item interaction graph}\label{sec:model-graph}

The graph uses history \(\mathcal H_q\) and eligible catalog
\(\mathcal I_q\) from Sections~\ref{sec:time} and~\ref{sec:requests}.
The time unit \(t_{\mathrm{unit}}=1\,\mathrm{s}\) makes logarithm arguments
dimensionless.
Topology and statistics are built only from $A_G$ by time $r_G$.
For recorded events, history eligibility is defined
in Section~\ref{sec:eval-observed}; confirmed telemetry completeness
is not assumed for this protocol.
The baseline policy requires \(r_G<t_q\); there is no separate training
of a graph encoder. At the second stage, the neural retriever
is trained on \(A_N\) and published by \(r_N\), so its requests
from \(B,C,E,D\) additionally require \(r_G<r_N<t_q\).

\subsection{Vertices, relations, and one session's contribution}

A graph vertex is an item identifier. Three relations are defined:
\begin{equation}\label{eq:model-relations}
\mathcal R_G=
\{\mathrm{detail}\!\to\!\mathrm{detail},\quad
  \mathrm{detail}\!\to\!\mathrm{add},\quad
  \mathrm{detail}\!\to\!\mathrm{purchase}\}.
\end{equation}
All events, including purchases, must precede the snapshot.

For a historical session $s$, item events are combined into nonempty
groups sharing a timestamp; groups are ordered by time.
Let $i_{s\ell},\kappa_{s\ell},t_{s\ell}$ denote the item, action,
and time of event $\ell$, and $k_s(\ell)$ its time-group index.
No directed edges are created within a group: the order is unknown.
For relation $r=(\kappa_1,\kappa_2)$, the set of eligible pairs is
\begin{align}\label{eq:model-event-pairs}
\mathcal P^r_{sij}(r_0)=\{(\ell,m):\;&i_{s\ell}=i,
\ i_{sm}=j,\ i\ne j,\ \kappa_{s\ell}=\kappa_1,\ \kappa_{sm}=\kappa_2,\nonumber\\
&0<t_{sm}-t_{s\ell}\le\Delta_G,\quad
0<k_s(m)-k_s(\ell)\le B_G,\quad t_{sm}<r_0\}.
\end{align}
Here $r_0$ is snapshot time, $\Delta_G>0$ is the time window,
and $B_G\in\N$ is the maximum distance between time groups
of item events. Events without an item do not create such groups.
Pairs do not depend on row order within a group; repeated item occurrences
are handled by session features, without self-loops $i\to i$.

A session's contribution is bounded by one and defined by a maximum:
\begin{equation}\label{eq:model-session-edge}
w^r_{sij}=
\begin{cases}
\displaystyle\max_{(\ell,m)\in\mathcal P^r_{sij}(r_0)}
\exp\!\left(-\frac{t_{sm}-t_{s\ell}}{\tau_G}\right),
&\mathcal P^r_{sij}(r_0)\ne\varnothing,\\[0.4em]
0,&\text{otherwise},
\end{cases}
\end{equation}
where $\tau_G>0$ has time units. Then
\begin{equation}\label{eq:model-edge-weight}
W^r_{ij}=\sum_s w^r_{sij},\qquad
n^r_{ij}=\sum_s\ind\{w^r_{sij}>0\},\qquad
0\le W^r_{ij}\le n^r_{ij}.
\end{equation}
The sums use snapshot sessions truncated at the snapshot boundary
under the data preparation protocol.

\subsection{Sparsification and graph-row regularization}

Among edges with support $n^r_{ij}\ge n_G\ge1$, retain
at most $d_G$ neighbors with the greatest $W^r_{ij}$ for each pair $(r,i)$.
Ties are resolved by a fixed identifier order.
The retained neighbor set is denoted $\mathcal N^r_i$.
\begin{equation}\label{eq:model-transition}
P^r_{ij}=
\begin{cases}
\displaystyle\frac{W^r_{ij}}
 {\sum_{k\in\mathcal N^r_i}W^r_{ik}},
&j\in\mathcal N^r_i,\\[0.4em]
0,&\text{otherwise}.
\end{cases}
\end{equation}
An empty row is zero; a nonempty row sums to one.
$P^r_{ij}$ expresses normalized association strength, not the probability
of the next action: one session may produce several neighbors.

To attenuate rows with low total weight, use
\begin{equation}\label{eq:model-edge-confidence}
d^r_i=\sum_{j\in\mathcal N^r_i}W^r_{ij},\qquad
\rho^r_i=\frac{d^r_i}{d^r_i+\lambda_G},\qquad \lambda_G>0.
\end{equation}
$\lambda_G$ has the same arbitrary units as $d^r_i$.
Coefficient $\rho^r_i$ attenuates low-weight rows;
it does not estimate edge confidence. For a retained edge,
\begin{equation}\label{eq:model-row-regularization}
\rho_i^rP_{ij}^r=\frac{W_{ij}^r}{d_i^r+\lambda_G},
\qquad
0\le \rho_i^rP_{ij}^r\le P_{ij}^r\le1.
\end{equation}
At fixed $W_{ij}^r$, increasing other neighbors' weights attenuates the edge.
The threshold $n_G$ provides additional protection against accidental associations.

Let $\mathcal D_q$ be the latest at most $L_H$ available
\code{detail} events. When truncating history, equal timestamps
are resolved by item identifier; this is a technical selection rule,
not a reconstruction of actual action order.
Normalized history weights are defined as
\begin{equation}\label{eq:model-history-weight}
\alpha_{q\ell}=
\frac{\exp(-(t_q-t_{s(q)\ell})/\tau_H)}
 {\sum_{m\in\mathcal D_q}\exp(-(t_q-t_{s(q)m})/\tau_H)},
\qquad \ell\in\mathcal D_q,\quad\tau_H>0.
\end{equation}
For empty history, sums involving $\alpha_{q\ell}$ are zero.
Cart actions are represented by separate features.

The graph score of candidate $j$ for relation $r$ is
\begin{equation}\label{eq:model-graph-score}
g^r_{qj}=\sum_{\ell\in\mathcal D_q}
\alpha_{q\ell}\rho^r_{i_{s(q)\ell}}
P^r_{i_{s(q)\ell},j}.
\end{equation}
Each relation's score is a separate feature:
$0\le g^r_{qj}\le1$, with $g^r_{qj}=0$ for empty history.
A control variant may use the mixture
$g_{qj}=\sum_r\omega_r g^r_{qj}$, where
$\omega_r\ge0$, $\sum_r\omega_r=1$.
This mixture is not required in the first run:
each relation generates its own candidates and features.

Additional attenuation of individual edges is described in Appendix~\ref{sec:model-edge-optional}.
The first run uses only \eqref{eq:model-graph-score}.

\begin{invariant}[Historical information boundary]
Changing events with time $t\ge r(q)$ does not change
$W^r,P^r,\rho^r$ of this snapshot. Changing events with time
$t>t_q$ does not change $\alpha_{q\ell}$ or $g^r_{qj}$.
\end{invariant}

\subsection{Construction and request cost}

Let $N_{A_G}$ be the number of historical item events,
and $b_G^{\max}$ the maximum size of one time group in the snapshot.
After sorting and grouping, pair enumeration requires on the order of
$N_{A_G}B_Gb_G^{\max}$ operations: for each event, up to $B_G$
subsequent groups of size at most $b_G^{\max}$ are checked.
Observed $b_G^{\max}$ is included in the cost report;
arbitrary splitting of a simultaneous group is prohibited.
The number of retained edges is at most
$|\mathcal R_G|\,|\mathcal I_G|d_G$.
Computing graph candidates for one request accesses
at most $|\mathcal R_G|L_Hd_G$ edges.
Intermediate aggregation memory is measured separately from the final graph.
No dense $|\mathcal I_G|^2$ matrix is created; $B_G,L_H,d_G$
bound the computational cost.

\section{Content features and search context}\label{sec:model-content}

\subsection{Separate representation spaces}

For item $i$, denote the provided description and image vectors
by $v_i^{\mathrm{text}}\in\R^{d_t}$ and
$v_i^{\mathrm{image}}\in\R^{d_v}$.
In catalog notation, $v_i^{\mathrm{text}}=d_i$,
$v_i^{\mathrm{image}}=v_i$.
The coordinates of the two spaces have different meanings even when $d_t=d_v$.
Coveo documentation states that the query vector is compatible
with the item description vector; compatibility with the image vector
is not claimed~\cite{coveo-data}.

For each modality $m\in\{\mathrm{text},\mathrm{image}\}$,
define a mask $M_i^m$. It equals one if the vector exists,
has the expected dimension, finite coordinates, and a norm greater than
a fixed $\eps_v>0$. Define
\begin{equation}\label{eq:model-vector-normalization}
\bar v_i^m=
\begin{cases}
v_i^m/\norm{v_i^m},&M_i^m=1,\\
0,&M_i^m=0.
\end{cases}
\end{equation}
A missing vector is represented by zero and a mask; a zero cosine
between valid vectors remains an observed value.

The history representation is built separately in each space:
\begin{align}\label{eq:model-session-vector}
u_q^m&=\sum_{\ell\in\mathcal D_q}
\alpha_{q\ell}M^m_{i_{s(q)\ell}}\bar v^m_{i_{s(q)\ell}},\\
M_q^m&=\ind\{\norm{u_q^m}>\eps_v\},\qquad
\bar u_q^m=\begin{cases}u_q^m/\norm{u_q^m},&M_q^m=1,\\
0,&M_q^m=0.\end{cases}\nonumber
\end{align}
The norm of the sum is checked because directions may cancel.
The content score and its mask are
\begin{equation}\label{eq:model-content-score}
c^m_{qi}=(\bar u_q^m)^\top\bar v_i^m,\qquad
M^m_{qi}=M_q^m M_i^m.
\end{equation}
For $M^m_{qi}=1$, the score lies in $[-1,1]$.
Retrieval may use a weighted mixture of available modalities only:
\begin{equation}\label{eq:model-content-fusion}
c_{qi}=
\begin{cases}
\displaystyle\frac{\sum_m\gamma_m M^m_{qi}c^m_{qi}}
 {\sum_m\gamma_m M^m_{qi}},&\sum_m\gamma_m M^m_{qi}>0,\\[0.3em]
0,&\text{otherwise},
\end{cases}
\qquad \gamma_m\ge0.
\end{equation}
The ranker receives separate scores and masks.
Variants are compared under identical candidate quotas.

\subsection{Latest available search query}

Let $b(q)$ be the latest search block available by $t_q$,
and $t_{b(q)}$ its time. If there is none, or if
$t_q-t_{b(q)}>\Delta_S$, search context is considered absent.
An unusable latest block under
Section~\ref{sec:data-search-canonicalization} disables search;
there is no return to an earlier block.
Threshold $\Delta_S$ is specified in seconds. A valid query vector
is normalized by rule~\eqref{eq:model-vector-normalization};
the result is denoted $\bar z_q\in\R^{d_t}$,
and the availability mask is $M_q^S$.
\begin{equation}\label{eq:model-search-score}
s_{qi}^{\mathrm{search}}=\bar z_q^\top\bar v_i^{\mathrm{text}},\qquad
M_{qi}^S=M_q^S M_i^{\mathrm{text}}.
\end{equation}
$s_{qi}^{\mathrm{search}}$ measures query--description similarity;
comparison with a graph score requires training or normalization.

Additional search-context features are
\begin{equation}\label{eq:model-search-features}
d_q^S=\begin{cases}
t_q-t_{b(q)},&M_q^S=1,\\
0,&M_q^S=0,
\end{cases}
\qquad M_q^S,\quad
a_q^S=\log\!\left(1+\frac{d_q^S}{t_{\mathrm{unit}}}\right).
\end{equation}
Quantity $d_q^S$ is measured in seconds; $a_q^S$ is dimensionless.
When search is absent, both age and mask are zero; the missing
$t_{b(q)}$ is not evaluated in the zero branch.
Field \code{product_skus_hash} is excluded from features and retrieval:
candidates are selected from $\mathcal I_q$ by similarity to \code{query_vector},
without copying the historical search results.

\begin{invariant}[Feedback after a search query]
Field \code{clicked_skus_hash} is not a feature of the current
search query: it contains subsequent clicks.
Without individual click times, the field is also excluded
from state reconstruction for a later request.
A separate item-view event with a known time
may be included under the normal availability rule.
\end{invariant}

Search-click prediction is allowed as a separate experiment
with its own labels. Absence of a click does not imply a negative preference
and does not change the purchase objective in Section~\ref{sec:business}.

\subsection{Categories and price bucket}

For categories of the candidate and a viewed item with path lengths
$\ell_i^{\mathrm{cat}},\ell_j^{\mathrm{cat}}$, similarity is defined as
\begin{equation}\label{eq:model-category-similarity}
c_{ij}^{\mathrm{cat}}=
\frac{\operatorname{lcp}(i,j)}
 {\max(1,\ell_i^{\mathrm{cat}},\ell_j^{\mathrm{cat}})},
\end{equation}
where $\operatorname{lcp}$ is the number of matching initial categories.
For a missing path, the value is zero and a separate mask is set.
A hash serves as a categorical label.

For ordinal feature \code{price_bucket}, use equality,
the sign and absolute difference of bucket indices, with missing-value masks.

\section{Candidate set construction}\label{sec:model-candidates}

\subsection{Sources and quotas}

All sources are restricted to $\mathcal I_q$ under Section~\ref{sec:requests}:
already purchased items are excluded; viewed items and items observed
in the cart are retained and evaluated in separate slices.
Let $\mathcal C_q^{(a)}$ denote the list from source $a$.
The first run uses six source groups:
\begin{enumerate}
\item the most recent distinct items from available views;
\item neighbors under each of the three graph relations;
\item items most similar in the description space;
\item items most similar in the image space;
\item eligible catalog items similar to the latest search query;
\item historically popular items;
\end{enumerate}
Initial quotas are 50 candidates for each of the three graph relations,
40 for descriptions, 30 for images, 50 for search,
20 recent items, and 10 popular items. Their sum is \(M=300\).
A neural source is added only at the second stage
(Section~\ref{sec:modern-retrieval}); a semantic-code source is introduced
in a separate study (Section~\ref{sec:sid-serving}).
An unavailable source returns an empty list and a reason code.
A content source is enabled only with a valid history
representation; arbitrary sorting of entirely missing scores is not used.

Quotas $m_a\ge0$ are fixed before final evaluation:
\begin{equation}\label{eq:model-candidate-union}
\mathcal C_q^{\mathrm{pre}}=\bigcup_a\mathcal C_q^{(a)},\qquad
|\mathcal C_q^{(a)}|\le m_a,\qquad
\sum_a m_a\le M.
\end{equation}
Duplicates are merged, preserving all source features.
The neural-source quota covers all its subqueries;
a new source redistributes quotas within the existing $M$.
After merging, remaining slots are filled by popularity
from $\mathcal I_q\setminus\mathcal C_q^{\mathrm{pre}}$.
Let $\mathcal C_q^{\mathrm{fill}}$ be the first at most
$M-|\mathcal C_q^{\mathrm{pre}}|$ such items. The final set is
\begin{equation}\label{eq:model-candidate-fill}
\mathcal C_q^{\mathrm{main}}=\mathcal C_q^{\mathrm{pre}}\cup\mathcal C_q^{\mathrm{fill}},
\qquad |\mathcal C_q^{\mathrm{main}}|\le M.
\end{equation}
Filling has a separate feature and is not limited by the initial
popularity quota. It stops at $M$ or catalog exhaustion.
Adding the future correct answer is prohibited.

This is the primary candidate set. Before ranking,
$\mathcal C_q^{\mathrm{fallback}}$ is prepared separately: up to $M$ eligible items, starting
with distinct recent viewed items in descending time order
(ties resolved by identifier), then popularity without duplicates.
Fallback is not merged into the primary set beyond the budget.
The normative set for the branch actually selected is
\begin{equation}\label{eq:model-served-candidates}
\mathcal C_q^{\mathrm{served}}=
\begin{cases}
\mathcal C_q^{\mathrm{main}},&\mathrm{branch}_q=\mathrm{main},\\
\mathcal C_q^{\mathrm{fallback}},&\mathrm{branch}_q=\mathrm{fallback},\\
\varnothing,&\mathrm{branch}_q\in\{\mathrm{empty},\mathrm{error}\}.
\end{cases}
\end{equation}
The branch is fixed when choosing a response, reserving time for mandatory
logging and return. If the selected branch subsequently fails to return in time,
retain its candidates, $L_q=\varnothing$, and a separate failure code:
the loss belongs to delivery, not retrieval. The value \code{error} in the formula
means a failure before selecting a ready branch. In all cases,
$L_q\subseteq\mathcal C_q^{\mathrm{served}}\subseteq\mathcal I_q$
and $|\mathcal C_q^{\mathrm{served}}|\le M$.
Below, $\mathcal C_q$ without a superscript denotes
$\mathcal C_q^{\mathrm{served}}$; for normal ranker training on $B$,
it is $\mathcal C_q^{\mathrm{main}}$. Fallback has its own fixed order
and is not trained using the primary ranker's gradients. An unfinished primary
set is logged as unmeasured, not as an empty diagnostic list.

Historical popularity is defined by the number of distinct
snapshot sessions with a view of the item:
\begin{align}\label{eq:model-popularity}
V_i^{\mathrm{detail}}&=\sum_{s\in A_G}
\ind\{s\text{ contains an available view of }i\},\\
\pi_i^{\mathrm{pop}}&=\frac{V_i^{\mathrm{detail}}+\eta_P}
 {\sum_{j\in\mathcal I_G}(V_j^{\mathrm{detail}}+\eta_P)},
\quad\eta_P>0.
\nonumber\end{align}
At request time, popularity is restricted to $\mathcal I_q$.
An unknown item may change session state and its content
representation but receives no historical edges. Admission of new
items follows the availability rule in Section~\ref{sec:requests}.

\subsection{Neighbor search in vector space}

The first run computes similarity against every eligible vector:
the cost for a modality of dimension $d_m$ is on the order of $|\mathcal I_q|d_m$.
If measured latency is excessive, approximate
nearest-neighbor search (ANN) is tested separately.
Its settings are fixed before final evaluation; a separate request
sample is used to measure
\begin{equation}\label{eq:model-ann-agreement}
\mathrm{Agreement@}m=
\frac1{N_Q^+}\sum_{q:\,|\mathcal C_{q,\mathrm{exact}}|>0}
\frac{|\mathcal C_{q,\mathrm{approx}}\cap
\mathcal C_{q,\mathrm{exact}}|}
 {|\mathcal C_{q,\mathrm{exact}}|}.
\end{equation}
$N_Q^+$ is the number of nonempty exact responses; when $N_Q^+=0$, the metric
is undefined. The empty-response fraction is reported separately.
An identical eligibility filter allows approximation losses to be measured,
but agreement with exact search does not replace purchase evaluation.

\subsection{Candidate provenance features}

For each source, retain its mask, original score,
and reciprocal rank:
\begin{equation}\label{eq:model-source-rank}
o_{qi}^{(a)}=\ind\{i\in\mathcal C_q^{(a)}\},\qquad
u_{qi}^{(a)}=
\begin{cases}
1/(1+\operatorname{rank}_a(q,i)),&o_{qi}^{(a)}=1,\\
0,&o_{qi}^{(a)}=0.
\end{cases}
\end{equation}
Ranks start at one; missing ranks and scores are represented by zero
with a mask. A measurable score is computed for all candidates in the union,
even if the item did not enter the corresponding source's quota.
The transformation $1/(1+\operatorname{rank})$ intentionally gives $1/2$
to the first position; this is a bounded feature, not a probability.

\section{Request--item pair features}\label{sec:model-features}

Vector $x_{qi}\in\R^p$ contains only quantities available at $q$.
The category transformation is fixed with the model;
an item identifier is not used as a real-valued number.
The table defines features for the first run.

{\small
\begin{longtable}{p{0.23\textwidth}p{0.65\textwidth}}
\toprule
Group & Contents \\
\midrule
Graph & $g^r_{qi}$, connection to the latest item group, edge support,
masks for missing neighbors \\
Description and image & $c^m_{qi}$, similarity to the latest item group,
vector masks and empty-history-representation masks \\
Search & $s_{qi}^{\mathrm{search}}$, query age,
availability masks for the query and item description \\
Item history & Number of available views, time since the latest
view, latest add or remove action \\
Session history & Number of events and distinct items, duration
of available history, search presence, number of recorded adds \\
Catalog & Category, category matches, price bucket,
and its difference from buckets of the most recently viewed items \\
Snapshot & Popularity and number of sessions with a recorded purchase of the item \\
Retrieval & Source masks, original scores, and reciprocal ranks \\
\bottomrule
\end{longtable}}

The second stage adds the learned session-vector score,
the maximum intent score, intent-score variance, weight entropy,
search similarity in the learned space, and availability masks.
The semantic-code study adds the log probability
of the code, the number of items with that code, and a mask indicating a found code.
A code probability is not a purchase probability.
When the feature set changes, the ranker is retrained on \(B\);
neural features are not needed to implement the first stage.

Set $\mathcal I_q^{\mathrm{last}}$ contains all distinct items
in the latest \code{detail} time group, without selecting a single last item.
Weight $n_{qj}^{\mathrm{last}}$ equals $1$ in \code{set_v1} and the view
multiplicity of $j$ in that group in \code{multiset_v1}, after removal
of transport retries. For pairwise feature $\phi(i,j)$ with availability
mask $M_{ij}^{\phi}$, its mean and mask are
\begin{align}\label{eq:model-last-group-mean}
Z_{qi}^{\phi}&=\sum_{j\in\mathcal I_q^{\mathrm{last}}}
 n_{qj}^{\mathrm{last}}M_{ij}^{\phi},\qquad
M_{qi}^{\mathrm{last},\phi}=\ind\{Z_{qi}^{\phi}>0\},\nonumber\\
\overline\phi_{qi}&=\begin{cases}
\displaystyle\frac{\sum_{j\in\mathcal I_q^{\mathrm{last}}:\,M_{ij}^{\phi}=1}
 n_{qj}^{\mathrm{last}}\phi(i,j)}{Z_{qi}^{\phi}},&Z_{qi}^{\phi}>0,\\
0,&Z_{qi}^{\phi}=0.
\end{cases}
\end{align}
For descriptions and images, $\phi(i,j)$ is the cosine of valid vectors,
and $M_{ij}^{\phi}=M_i^mM_j^m$. For price buckets, average equality,
$\operatorname{sign}(b_i^{\mathrm{price}}-b_j^{\mathrm{price}})$,
and absolute difference; the mask requires both values to be known.
An unavailable $\phi(i,j)$ is not evaluated.
The graph maximum $\max_{j\in\mathcal I_q^{\mathrm{last}}}\rho_j^rP^r_{ji}$
and maximum category similarity do not depend on multiplicity;
an empty eligible set gives zero and mask $0$.

The item $i$ view count and time since its last view are defined as
\begin{align}\label{eq:model-online-features}
n_{qi}^{\mathrm{detail}}&=
\sum_{e\in\mathcal H_q}\ind\{i_e=i,\ \kappa_e=\mathrm{detail}\},\\
\delta_{qi}^{\mathrm{detail}}&=
t_q-\max\{t_e:e\in\mathcal H_q,\ i_e=i,\ \kappa_e=\mathrm{detail}\}.
\nonumber
\end{align}
With no event, age and mask are zero.
For numerical stability, the count may be transformed as
$\log(1+n)$ and age as $\log(1+\delta/t_{\mathrm{unit}})$.

Cart state and its age are defined in Section~\ref{sec:time}.

A smoothed purchase rate is not used: without confirmed
observation completeness, its denominator is not established. This feature is excluded
from the schema, not filled with zero. Recorded purchase counters are retained.

Statistics use only $A_G$; target encoding from labels
in the ranker-training period is excluded. A new category receives
a separate ``unknown category'' code.

\section{Learning to rank future purchases}\label{sec:model-ranking}

\subsection{Labels, groups, and weights}

The ranker trains on the first requests of sessions $\mathcal S_B^+$
with a recorded target (Section~\ref{sec:eval-observed}).
The model computes score $f(x_{qi})$ with label
$y_{qi}=\ind\{i\in Y_q\}$, $i\in\mathcal C_q$.
Request $q$ forms a comparison group. A zero label means
no observed purchase in the window, not a negative preference.

Positive and remaining candidates form the sets
\begin{equation}\label{eq:model-pair-sets}
\mathcal C_q^+=\{i\in\mathcal C_q:y_{qi}=1\},\qquad
\mathcal C_q^-=\{i\in\mathcal C_q:y_{qi}=0\},\qquad
\mathcal J_q=\mathcal C_q^+\times\mathcal C_q^-.
\end{equation}
One session contributes one group. The weights are
\begin{equation}\label{eq:model-ranking-weights}
w_q^{\mathrm{sess}}=1,\qquad
w_q^{\mathrm{pair}}=\frac{w_q^{\mathrm{sess}}}{\max(1,|\mathcal J_q|)}.
\end{equation}
After multiplication by $\Delta_{qij}$, the sum of pair weights
and gradient norm depend on candidate composition and current ranks;
equal actual contributions from different sessions are not claimed.
A group without a positive item contributes no pairwise gradient,
but a session with a recorded target remains in the evaluation denominator.
The target item is not added to candidates. If there are no training pairs,
report insufficient data, not successful training.

\begin{assumption}[Surrogate comparison under implicit feedback]
\label{ass:model-implicit-negatives}
Each observed purchased candidate is trained to outrank
every candidate without a purchase in the same window. This is an assumption about a useful
surrogate for the historical task, not a statement about exposure
or preference: the unpurchased item may never have been shown.
The historical policy, availability, and popularity affect labels;
without an impression log, this training does not eliminate selection bias.
\end{assumption}
This assumption does not require replacing the task
with positive--unlabeled learning: its identifying assumptions
would also need justification. Policy selection is checked
by an independent historical comparison and then an A/B protocol.

\subsection{Pairwise logistic loss}

For positive $i$ and negative $j$, let
$d_{qij}=f(x_{qi})-f(x_{qj})$.
The basic training objective is
\begin{equation}\label{eq:model-pairwise-loss}
\mathcal L_{\mathrm{pair}}(f)=
\sum_q w_q^{\mathrm{pair}}\sum_{(i,j)\in\mathcal J_q}
\log\!\left(1+\exp(-\beta_R d_{qij})\right),
\qquad \beta_R>0.
\end{equation}
Parameter $\beta_R$ sets the score scale.
Compute the stable $\operatorname{softplus}(-\beta_Rd_{qij})$,
not a direct exponential of a large positive argument.
At $d_{qij}=0$, the derivative with respect to the positive score is
$-w_q^{\mathrm{pair}}\beta_R/2$, so minimization increases the positive score
and decreases the negative score. The logistic transformation of their difference
describes pair order, not the probability of an item purchase.

\subsection{Greater weight for errors near the top of the list}

In LambdaMART, pair weight depends on how much swapping the pair
would change top-position quality~\cite{lambdamart}.
Let $r_{qi}$ be the current position of $i$ among all candidates,
and $K_R$ the ranking-training depth. Ranks start at one;
order is defined by key $(-f(x_{qi}),\operatorname{id}(i))$,
with identifiers compared lexicographically by UTF-8 bytes.
For binary labels, use gain $G_R(0)=0$, $G_R(1)=1$.
Define
\begin{equation}\label{eq:model-rank-discount}
D_R(k)=\frac{\ind\{k\le K_R\}}{\log_2(1+k)},\qquad
Z_q=\sum_{k=1}^{\min(K_R,|\mathcal C_q^+|)}D_R(k).
\end{equation}
For $Z_q>0$, normalized list quality is
$\sum_{i\in\mathcal C_q}y_{qi}D_R(r_{qi})/Z_q$.
Its absolute change after swapping a positive and a negative
item is
\begin{equation}\label{eq:model-lambda-weight}
\Delta_{qij}=
\frac{|D_R(r_{qi})-D_R(r_{qj})|}{Z_q},\qquad
\omega_{qij}=w_q^{\mathrm{pair}}\Delta_{qij}.
\end{equation}
At $Z_q=0$, there are no pairs and gradients are zero.
A pair with both positions below $K_R$ has zero weight. If exactly one
position is above the cutoff, use the full absolute discount difference;
there is no additional pair truncation or normalization by the score
difference.

Weights $\omega_{qij}$ are recomputed from current positions
and held fixed when taking derivatives. This pseudo-gradient rule
does not assume a single smooth function with all these derivatives.
At $\Delta_{qij}=1$, it reduces to pairwise logistic training.

For one pair, denote
$p_{qij}=1/(1+\exp(\beta_Rd_{qij}))$.
Contributions to the first derivative and diagonal second derivative are
\begin{align}\label{eq:model-lambda-derivatives}
g_{qi}^{R,(ij)}&=-\omega_{qij}\beta_Rp_{qij},&
g_{qj}^{R,(ij)}&=+\omega_{qij}\beta_Rp_{qij},\\
h_{qi}^{R,(ij)}&=\omega_{qij}\beta_R^2p_{qij}(1-p_{qij}),&
h_{qj}^{R,(ij)}&=h_{qi}^{R,(ij)}.\nonumber
\end{align}
Derivatives are summed over each candidate's pairs. The trees use
a diagonal approximation, omitting mixed second derivatives.
The group's gradients sum to zero: the absolute score level is unspecified.
For numerical implementation, set $u=\beta_Rd_{qij}$,
$v=\exp(-|u|)$. Then $p=v/(1+v)$ for $u\ge0$ and $p=1/(1+v)$
for $u<0$, while $p(1-p)$ is computed as $v/(1+v)^2$.
This preserves small curvature without subtracting nearby numbers;
no additional positive floor is imposed on $h$.

\subsection{Tree ensemble and regularization}

The ranker is a sum of trees:
\begin{equation}\label{eq:model-boosted-trees}
f_{N_T}(x)=f_0+\nu\sum_{m=1}^{N_T} T_m(x),\qquad 0<\nu\le1.
\end{equation}
A tree leaf outputs one number; $\nu$ limits a new tree's contribution.
Constant $f_0$ does not affect order.

For leaf $\ell$ of the next tree, denote
$G_\ell=\sum_{(q,i)\in\ell}g_{qi}^{R}$ and
$H_\ell=\sum_{(q,i)\in\ell}h_{qi}^{R}$.
With penalty $\lambda_2 w_\ell^2/2$ and $\lambda_2>0$,
the locally optimal leaf value is
\begin{equation}\label{eq:model-newton-leaf}
w_\ell^*=-\frac{G_\ell}{H_\ell+\lambda_2}.
\end{equation}
The formula follows by minimizing the quadratic expression
$G_\ell w+(H_\ell+\lambda_2)w^2/2$.
The reduction in this approximation when splitting a leaf into $L,R$ is
\begin{equation}\label{eq:model-tree-gain}
\mathrm{Gain}=\frac12\left[
\frac{G_L^2}{H_L+\lambda_2}+
\frac{G_R^2}{H_R+\lambda_2}-
\frac{(G_L+G_R)^2}{H_L+H_R+\lambda_2}\right]-\gamma_T.
\end{equation}
$\gamma_T\ge0$ is the penalty for adding a leaf.
A split is allowed when its gain is positive and constraints hold
on depth, leaf count, observation count,
and the sum of second derivatives in each child leaf.
The first run fixes a budget of $200$ iterations; actual
$N_T\le200$ is stored in the run manifest, because the absence of an eligible
split may stop new-tree construction earlier.
Later budget selection uses \(C_{\mathrm{select}}\) under a predefined
plan; after selection, the model is frozen until final evaluation.

\subsection{Executable contract for custom derivatives}
\label{sec:model-ranker-implementation}

Use the Python API \code{lightgbm.train}, version
\code{lightgbm==4.6.0}, with callable
\code{objective=session\_lambdas\_v1}~\cite{lightgbm}.
The function receives raw scores and a \code{Dataset}, and returns
the sums $g_{qi}^R,h_{qi}^R$ from~\eqref{eq:model-lambda-derivatives}.
At every iteration, full ranks, $Z_q$,
and all pairs from~\eqref{eq:model-pair-sets} are recomputed.
This is a custom LambdaMART objective; equivalence to the built-in
\code{objective=lambdarank} is not claimed. The built-in objective's parameters
\code{label\_gain}, \code{lambdarank\_norm},
\code{lambdarank\_truncation\_level}, and \code{sigmoid}
are not set here and do not control the callback:
gain $[0,1]$, $K_R=20$, and $\beta_R=1$ are implemented explicitly within it.

Rows are grouped by request, requests are ordered by
$(s,t_q,\operatorname{id}(q))$, and rows within a group by item
identifier; all orders are part of the data schema. Nonempty group sizes
are passed as \code{Dataset.group}. An empty group creates no rows,
but remains in $|\mathcal Q_s|$ for weighting. Weights
$w_q^{\mathrm{pair}}$ are applied once inside the callback;
\code{Dataset.weight=None}, with no additional row/query weights.
Groups without pairs receive zero derivative arrays.
Sums are computed in FP64, then the callback returns contiguous
FP32 arrays: this is the type accepted by the selected version's C API. Nonfinite
scores, gradients, or Hessians stop training with an error
instead of being replaced by zeros.

The complete first-ranker configuration below is supplemented only
by the fixed feature schema and data; unspecified parameters
use the defaults of version 4.6.0 specifically. The serialized run manifest
contains default values, the derivative-function hash, a version file
with package hashes, and platform/compiler information. A new
library version requires the contract to be checked again.
\begin{verbatim}
objective=session_lambdas_v1 (callable); metric="None"
boosting="gbdt"; data_sample_strategy="bagging"
device_type="cpu"; tree_learner="serial"; num_threads=1
deterministic=true; force_col_wise=true; force_row_wise=false
learning_rate=0.05; num_leaves=31; max_depth=8
lambda_l1=0; lambda_l2=1; min_gain_to_split=0
min_data_in_leaf=100; min_sum_hessian_in_leaf=0.001
bagging_fraction=1; bagging_freq=0; bagging_by_query=false
feature_fraction=1; feature_fraction_bynode=1; extra_trees=false
max_bin=255; min_data_in_bin=3; bin_construct_sample_cnt=200000
feature_pre_filter=false; enable_bundle=false
use_missing=true; zero_as_missing=false; use_quantized_grad=false
boost_from_average=false; linear_tree=false
path_smooth=0; max_delta_step=0
cat_l2=0; cat_smooth=1; max_cat_threshold=32
max_cat_to_onehot=4; min_data_per_group=100
seed=17; data_random_seed=17; feature_fraction_seed=17
bagging_seed=17; extra_seed=17; verbosity=1
train: num_boost_round=200; valid_sets=None; feval=None
       init_model=None; callbacks=None
Dataset: weight=None; init_score=None
predict: raw_score=true; num_iteration=N_T
\end{verbatim}
The initial score is $f_0=0$. Category is the only native
categorical column in the first run: its name is explicitly passed
in \code{categorical\_feature}; the dictionary is built from $A_G$ in
UTF-8 order, with code $0$ for unknown values and $1,2,\ldots$ for known values.
The price bucket is an ordinal number with a separate mask.
The set, order, and types of all remaining numerical features are stored
in the schema; the item identifier is not a feature.
There are no monotonic, interaction, or forced-split constraints.

LightGBM builds histogram-based splits, so enumeration of all
real-valued thresholds is not assumed. Its \code{min\_data\_in\_leaf}
is checked approximately through Hessians: it is a row parameter,
not a query parameter or a strict guarantee of $100$ actual rows.
Actual leaf sizes are published by routing the training rows;
a strict minimum would require a separate
algorithm and a new experiment. With L1, leaf smoothing
\code{path_smooth}, and the \code{max_delta_step} bound disabled,
LightGBM's gain has scale
$2(\mathrm{Gain}+\gamma_T)$ before threshold subtraction, so
\code{min\_gain\_to\_split}$=2\gamma_T$; both are zero
in the first run. \code{cat_l2=0} excludes the additional categorical leaf penalty
absent from~\eqref{eq:model-newton-leaf}.
The chosen positive \code{cat_smooth=1} acts differently:
when grouping categories, it restricts bin eligibility and smooths
the sorting key $G_{\mathrm{bin}}/(H_{\mathrm{bin}}+1)$.
Thus the allowed case $G_{\mathrm{bin}}=H_{\mathrm{bin}}=0$ does not produce
$0/0$ in the sorting key; callback values and the leaf formula are unchanged.
This parameter must not be disabled together with the leaf penalty.
The library's small internal numerical constants
are accommodated by test tolerance rather than bitwise equality.

The acceptance test compares all $g,h$ entries with an independent FP64
calculation: before casting, the callback tolerance is
$10^{-12}+10^{-10}|v_{\mathrm{ref}}|$; after FP32 conversion, it is
$10^{-7}+10^{-6}|v_{\mathrm{ref}}|$. Cases cover several groups,
two positive labels, tied scores, $K_R<|\mathcal C_q|$,
zero-contribution groups, differing candidate counts, and logits $\pm10^4$.
A separate single-step test on a numerical feature with an unambiguous
split checks leaves and doubled gain within $10^{-6}$;
size constraints are explicitly relaxed only in its fixture.
Permuting input rows after canonicalization must preserve
derivatives. Applying weights a second time must be detected
by the test, and built-in LambdaRank is not used as the reference.

\subsection{Negative examples and sparse purchases}

The primary version uses all candidates without negative sampling.
A separate variant with reweighted unpurchased candidates is preregistered
for sensitivity analysis. For each
$i\in\mathcal C_q^+$, instead of uniform weight $1/|\mathcal C_q^-|$,
it uses mixture $\tfrac12u_q(j)+\tfrac14v_q(j)+\tfrac14h_{qi}(j)$:
$u_q$ is uniform over $\mathcal C_q^-$, $v_q$ is proportional to
$1+N_j^{A_G}$ (the number of $A_G$ sessions containing the item), and $h_{qi}$ is uniform
over the first $\min(20,|\mathcal C_q^-|)$ unpurchased items under
fixed category and description similarity to $i$ (category match first,
then cosine, then identifier; unavailable
similarity is zero). The full pair weight in this variant is
$w_q^{\mathrm{sess}}\Delta_{qij}$ multiplied by this mixture
and divided by $|\mathcal C_q^+|$. This is a different surrogate objective,
not an exposure correction. The variant is trained on $B$ and compared
on $C_{\mathrm{select}}$ within the registered budget.
For both variants, publish Hit@20 and ranking quality
for head/tail, defined solely by $A_G$ popularity
(head is the first $\lceil0.1|\mathcal I_G|\rceil$ items
in descending $N_j^{A_G}$ order, with ties resolved by ID;
tail is all remaining items). Compare policy quality,
not numerical values of different loss functions.
If data reduction is unavoidable, pair $(i,j)$ is included
with known positive probability $\xi_{qij}$.
For a fixed current order, an unbiased estimator
of the local gradient sum uses weight
\begin{equation}\label{eq:model-pair-sampling}
\widetilde\omega_{qij}=
\frac{\ind\{(i,j)\text{ is sampled}\}}{\xi_{qij}}
\omega_{qij}.
\end{equation}
$\xi_{qij}$ refers to the pair; ranks and $Z_q$ use the full group.
Recomputing them on the sample does not recover the original LambdaMART.
Small $\xi_{qij}$ increase variance; clipping inverse weights
introduces bias and requires separate designation.

Missing purchases are not reconstructed with invented labels.
Scores from auxiliary next-view and add-to-cart tasks
may be used as features when trained without temporal leakage.

A separate future study may predict a purchase within a session
from available history and candidates. A single scalar prediction for a request
does not itself compare items or different lists: adding
a common value to scores preserves their order. Such a prediction is useful
for diagnostics or auxiliary training but does not estimate
the causal utility of a recommendation list and is not part of the first run.

\subsection{Limits on score interpretation}\label{sec:model-calibration}

Calibration is disabled; scores are not purchase probabilities.
Period $E$ remains reserved; it does not supplement training
or model selection. Probabilistic methods of list construction
require a separate formulation and suitable data.

\section{Final recommendation list construction}\label{sec:model-slate}

\subsection{Primary recommendation method}

The baseline policy and neural extension sort $\mathcal C_q$ by $f(x_{qi})$
and return the first $K_q=\min(K,|\mathcal C_q|)$ distinct items.
Calibration with $a_{\mathrm{cal}}>0$ does not alter this order.
Tied scores use the fixed identifier
order. The first run uses \(K=20\); all comparisons share
the same upper limit $K$ and length rule, not the actual $K_q$.
Short and empty responses are included without artificial padding.

The first run and initial neural variant use \(\lambda_D=0\):
no item-similarity penalty is applied; the calibrator is disabled.
Changes to list diversity and probabilistic constraints are discussed in Appendix~\ref{sec:model-diversity-optional}.

\subsection{Fallback response}

If the ranker is unavailable, a separate fallback
policy is used: the most recent eligible items and historical popularity
in the order of $\mathcal C_q^{\mathrm{fallback}}$ from
\eqref{eq:model-served-candidates}; the first $K$ items are returned.
If no eligible items exist,
an empty list is returned. All eligibility constraints
remain in force on failure; the reason code is included in the request log.
The primary model's calibrator is not applied to fallback scores.
Abstaining from display and changing $K$ define separate policies;
their causal utility is tested under Section~\ref{sec:eval-online},
with a no-widget control if needed.
Recommendation explanations are defined in Section~\ref{sec:explanation}.


\section{Temporal protocol for training and evaluation}\label{sec:eval-temporal}

The protocol is shared by the core system and neural extension:
inputs contain available history, and targets are subsequent behavior.
The neural-training period remains reserved in the first run.

\subsection{Periods and their roles}\label{sec:eval-periods}

Let $A_G,A_N,B,C,E,D$ be disjoint calendar intervals
in the stated order, with $A=A_G\cup A_N$:
\begin{equation}\label{eq:eval-period-order}
\sup A_G<\inf A_N,\quad \sup A_N<\inf B,\quad
\sup B<\inf C,\quad \sup C<\inf E,\quad \sup E<\inf D.
\end{equation}
Gaps and boundary losses are determined by the target horizon and time buffer,
not by the observed target value.
Period $C$ is split in advance into two disjoint subintervals:
\begin{equation}\label{eq:eval-validation-split}
C=C_{\mathrm{smoke}}\cup C_{\mathrm{select}},\qquad
C_{\mathrm{smoke}}\cap C_{\mathrm{select}}=\varnothing,\qquad
\sup C_{\mathrm{smoke}}\le\inf C_{\mathrm{select}}.
\end{equation}
After splitting $C$, the six major periods form seven parts,
whose boundaries are fixed in the run manifest.
Calendar admission~\eqref{eq:eval-admission} is applied separately to each
subinterval: debugging targets do not cross the selection boundary.
Throughout this document, parameter selection ``on $C$'' means only
$C_{\mathrm{select}}$, while implementation validation uses $C_{\mathrm{smoke}}$.

{\small
\begin{longtable}{L{0.10\textwidth}L{0.30\textwidth}L{0.47\textwidth}}
\toprule Period & Purpose & Permitted use \\\midrule\endhead
$A_G$ & Topology and reference data &
Graph, paths, catalog $\mathcal I_G$, popularities, recorded-purchase counters,
and historical feature scales. \\
$A_N$ & Neural retriever training &
Prefixes, recorded future purchases, auxiliary next views;
training the graph encoder, content projections, intents, and search.
Unused by the core system. \\
$B$ & Ranking training &
Candidates and features from the already fixed retriever;
fitting trees to subsequent purchases. \\
$C$ & Debugging and selection independent of debugging &
$C_{\mathrm{smoke}}$: implementation checks and fixes;
$C_{\mathrm{select}}$: one-time selection from a frozen list of variants
(dimensions, epochs, loss weights, quotas, tree count). \\
$E$ & Reserve & Unused in the current formulation. \\
$D$ & Final evaluation &
One independent evaluation of the frozen configuration and declared comparisons. \\
\bottomrule
\end{longtable}}

The start and end of each period, including both $C$ subintervals,
are fixed in the run manifest before training. In the first run, $A_N$ and $E$
are not merged into other periods or used for reference data.

\subsection{Sequence of the first run and final evaluation}
\label{sec:eval-stages}

\begin{enumerate}[leftmargin=*]
\item \textbf{Implementation validation.}
Run the graph, content features, search, and LambdaMART
with a $200$-iteration budget and no parameter tuning. On $C_{\mathrm{smoke}}$, check candidates,
features, ranking, and metrics; measure response time and memory,
and fix errors. $C_{\mathrm{select}},E,D$ remain closed.
These are local measurements; service load acceptance is not required.
The full variant without metadata does not block the first
$C_{\mathrm{smoke}}$; its mandatory deadlines are specified below.
After debugging, freeze the code, features, metrics, configuration list,
random initial states, and tie-breaking rule.
\item \textbf{Final-variant selection.}
Then compare configurations on $C_{\mathrm{select}}$ under the frozen plan
and computational budget. Results do not generate new variants.
Changing the implementation, features, metrics, or list after inspecting
$C_{\mathrm{select}}$ turns it into development data; another
confirmatory selection requires a new, unseen period.
Without such a period, the result remains exploratory.
An example plan is given
in Appendix~\ref{app:ablation}; it is not part of the first run.
\item \textbf{Independent evaluation.}
Before opening $D$, separately freeze the entire final evaluation plan,
check that it matches the selected model, and open $D$ once.
$E$ remains closed.
\end{enumerate}

A neural study is allowed after the baseline's data, formulas,
features, and costs have been checked; proven purchase lift is not required.
Each variant's training respects the period roles in the table above;
$C_{\mathrm{select}}$ selects results and saved states
but supplies no gradients.

The background catalog for training, normalization, and reconstruction is restricted
to $\mathcal I_G$. An unknown item's features may participate in projection
training only when that item is already observed in an $A_N$ prefix.
Its future interactions do not expand the candidate catalog.
The caveat in Assumption~\ref{ass:static-content}
applies to all provided vectors.
After $A_N$, weights and the vector index are fixed. Changing either
requires rebuilding $B$ features and retraining the ranker.

For semantic identifiers, first freeze the encoder
trained on $A_N$, then train the additional module
on that same period. Results on $A_N$ are not an independent evaluation.

When retrieval changes, trees are retrained on $B$.
New events in subsequent periods change only
the current session state, not reference data.

Changing the method based on $D$ results
makes subsequent evaluations on $D$ exploratory.

\subsection{Calendar admission and complete reading}\label{sec:eval-maturity}

Horizon $H$ and buffer $g$ are defined in Section~\ref{sec:requests}.
Every block $J=[a_J,b_J)$ uses the same rule:
\begin{equation}\label{eq:eval-admission}
t_{0,s}\in J,\qquad t_{0,s}+H+g<b_J.
\end{equation}
The rule is identical for sessions with and without a recorded purchase;
the next-view label also cannot cross the boundary.
After admission, IDs are sampled as described in Section~\ref{sec:data}.
All released behavior and search rows of a selected session
are read in full; session parts do not move between periods.
Omitting released rows is a preparation error.
Reading the complete export does not confirm telemetry completeness.

In $A_G$, an additional requirement is
$t_{\mathrm{last},s}^{\mathrm{obs}}+g<b_{A_G}$:
a session with a late tail is excluded in full, regardless of event type.
History includes selected sessions both with and without a purchase.
The catalog and statistics are not expanded with events from subsequent periods.
The time constraint does not prove actual session closure.

Let $N_J^{\mathrm{eligible}}$ be the number of sessions with an eligible first
request in $J$ before calendar exclusions and ID sampling, and
$N_J^{\mathrm{obs}}$ the retained count. Publish both counts,
exclusion reasons, and the fraction of retained sessions:
\begin{equation}\label{eq:eval-retention}
\operatorname{Retention}(J)=
\frac{N_J^{\mathrm{obs}}}{N_J^{\mathrm{eligible}}}.
\end{equation}
With an empty denominator, the ratio is undefined.

\begin{invariant}[Temporal availability with a fixed model]
\label{inv:eval-no-leakage}
For requests in $B,C,D$, features and candidates
depend only on the available prefix, the retrieval package published
before the request, and the stipulated static information.
Changing events after the request while holding the package fixed
does not change these quantities. $C_{\mathrm{smoke}}$ is for debugging,
and $C_{\mathrm{select}}$ for selection; independent final
evaluation of the selected model uses only $D$.
When training on $A_N$ and $B$, future events are allowed
in the target variable but not in the same example's input.
\end{invariant}

Trained weights are not considered available before $A_N$ requests.
Normalizations, negative-example distributions, and contrastive pairs
also have a declared source period.

\subsection{Model versions and final freeze}
\label{sec:eval-snapshots}

The core system uses one fixed package. In the deferred
graph-update mode (Appendix~\ref{sec:freshness}), packages
are published at times $r_1<r_2<\cdots$; a request uses the latest:
\begin{equation}\label{eq:eval-snapshot}
r(q)=\max\{r_j:r_j<t_q\}.
\end{equation}
Graph cutoff, encoder training, and publication times are retained.
Features in $B$ reproduce the intended version sequence;
a new encoder or index cannot be assigned to an old request.

Until an update is ready, the previous package is used and its age
is logged. In a confirmatory experiment, updates are prohibited
or fully specified by its protocol.

Before opening $D$, freeze the data, intervals, request rule,
$H,Q_{\max},K,M$, failure handling, all compared configurations,
primary metric, and interval construction. The selection budget covers
the number of configurations and random initializations, training time,
and peak memory. Equal candidate counts do not imply
equal computational cost.

\subsection{One primary final comparison}
\label{sec:eval-primary-comparison}

When studying an improvement to the core system, before opening $D$,
select one configuration $\pi_1$ on $C_{\mathrm{select}}$ and fix
control $\pi_0$: the graph, item features, search, and LambdaMART ranker.
The primary historical result is the paired difference
$\widehat\Delta_{\mathrm{Hit}}(\pi_1,\pi_0)$ on $D$ and its interval.
Each variant has its own ranker trained on $B$;
eligible data, request schedule, and candidate budget are shared.
Selection on $C_{\mathrm{select}}$ does not continue after inspecting $D$.

Other variants and slices are descriptive. For multiple confirmatory
claims, preregister $m_{\mathrm{hyp}}$ null hypotheses,
valid $p$-value calculations, and an overall level $\alpha$.
For example, Holm's procedure \cite{holm1979} orders
$p_{(1)}\le\cdots\le p_{(m_{\mathrm{hyp}})}$ and sequentially rejects
hypotheses while
\begin{equation}\label{eq:eval-holm}
 p_{(j)}\le\frac{\alpha}{m_{\mathrm{hyp}}-j+1};
\end{equation}
at the first violation, the remaining hypotheses are not rejected.
The procedure permits dependence between valid $p$-values but does not correct
errors in their calculation, including ignoring dependence among a user's sessions.
If only the core system is implemented, evaluate its coverage
and costs on $D$ without a superiority claim.


\section{Evaluation on the historical log}\label{sec:eval-offline}

The historical log can establish which subsequent purchases
would appear in the computed lists. It cannot reveal whether purchases would change
after actually displaying those lists. Thus the primary historical
metric measures coverage of observed purchases, while the causal business metric
is defined separately in Section~\ref{sec:business}.

\subsection{Session population and primary metric}
\label{sec:eval-observed}\label{sec:eval-schedule}

Protocol \code{observed_purchases_v1} evaluates the sole recommendation list
$q=q_s^{(1)}$ with $K=20$ against the recorded target
$Y_q=Y_s^{\mathrm{obs}}$ from~\eqref{eq:request-target}.
Let $\mathcal S_J^{\mathrm{obs}}$ be sessions in period $J$ that pass
calendar admission and sampling under Section~\ref{sec:eval-maturity}.
Their count is $N_J^{\mathrm{obs}}=|\mathcal S_J^{\mathrm{obs}}|$;
boundaries and subsample sizes are fixed before analyzing future purchases,
and shortfalls are not filled with purchasing sessions. For purchase-target
training and evaluation, define
\begin{equation}\label{eq:eval-observed-cohort}
 \mathcal S_J^+=\{s\in\mathcal S_J^{\mathrm{obs}}:
 Y_s^{\mathrm{obs}}\ne\varnothing\},\qquad N_J^+=|\mathcal S_J^+|.
\end{equation}
Selection into $\mathcal S_J^+$ is not passed to candidate retrieval, features, or the decision
to recommend. Historical reference data use all retained sessions
in $A_G$, not only purchasing sessions. Resource use and failures are measured
also on sessions without a recorded target: the outcome is unknown at request time.

For a fixed policy $\pi$:
\begin{align}
 \operatorname{Hit}_s(\pi)
 &=\ind\{L_{q_s^{(1)}}(\pi)\cap Y_s^{\mathrm{obs}}\ne\varnothing\},
 \label{eq:eval-session-hit}\\
 \operatorname{Hit@20}_J(\pi)
 &=\frac{1}{N_J^+}\sum_{s\in\mathcal S_J^+}\operatorname{Hit}_s(\pi),
 \qquad N_J^+>0.
 \label{eq:eval-observed-hit}
\end{align}
This is the fraction of sessions with a recorded future purchase in which
the first list contains at least one target item; the metric lies
in $[0,1]$. Purchase counts and repeated records do not increase a session's contribution.
Changing list length or request count requires a new protocol.

The membership of $\mathcal S_J^+$ is identical across methods. For failures,
late responses, or empty results, the specified fallback applies.
If an eligible list is not delivered on time, $L_q=\varnothing$;
the session remains in the denominator. A fallback hit counts,
and fallback frequency is published separately. Out-of-catalog targets are not removed.
At $N_J^+=0$, the metric is unavailable, not zero.

Telemetry completeness remains unknown even for positive sessions.
Sessions with $Y_s^{\mathrm{obs}}=\varnothing$ are not treated as confirmed
non-purchasing sessions. Evaluation is conditional on a future purchase being recorded:
it does not measure purchase frequency across all sessions, unobserved purchases,
or the causal effect of recommendations. Neither the point estimate nor its interval
eliminates event-recording or historical-policy bias.

Controls for the first evaluation are $A_G$ popularity and recently viewed
items filled with popularity; order and filters are specified
in Section~\ref{sec:model-slate}. Publish paired Hit@20 differences against them.
Metadata sensitivity is checked under~\ref{sec:metadata-sensitivity}.
Selection on $C_{\mathrm{select}}$ maximizes
\eqref{eq:eval-observed-hit}; the budget and tie rule are fixed before
$C_{\mathrm{select}}$, and the selected primary pair before $D$.
Intervals and sufficiency conditions
are specified in Section~\ref{sec:eval-uncertainty}; report contents are in
\ref{app:diagnostic-output}. Tuning on $D$ is prohibited.

\subsection{Exact upper bound for candidate retrieval}
\label{sec:eval-candidate-bound}

Below, for period $J$, every session sum is over $\mathcal S_J^+$,
$N=N_J^+$, $q=q_s^{(1)}$. Define
\begin{align}
 U_s(\pi)&=\ind\{\mathcal C_q^{\mathrm{served}}(\pi)\cap Y_q\ne\varnothing\},
 \label{eq:eval-candidate-hit}\\
 \operatorname{Coverage}_{\mathrm{cand},J}(\pi)
 &=\frac1N\sum_s U_s(\pi).
 \label{eq:eval-candidate-por}
\end{align}
Set $\mathcal C_q^{\mathrm{served}}(\pi)$ is defined
in~\eqref{eq:model-served-candidates}. Branch switching includes
fallback; a failure of an already selected response retains its candidates.
Coverage of a completed $\mathcal C_q^{\mathrm{main}}$ is reported separately
with the fraction of unmeasured sets; it does not replace the operational policy's bound.
The session population is not restricted to successful responses.

\begin{proposition}[Ranking limit under fixed retrieval]
\label{prop:eval-bound}
If $L_q(\pi)\subseteq\mathcal C_q^{\mathrm{served}}(\pi)$, then
\begin{equation}\label{eq:eval-candidate-bound}
 \operatorname{Hit@20}_J(\pi)
 \le\operatorname{Coverage}_{\mathrm{cand},J}(\pi).
\end{equation}
With fixed branches and candidates, $K\ge1$, unrestricted choice
of a subset of size at most $K$, and no delivery failures, the bound
is attainable by a ranker that knows future purchases. A fixed fallback order
or additional list constraints may make it unattainable;
the inequality itself still holds.
\end{proposition}

\begin{proof}
List inclusion implies $\operatorname{Hit}_s\le U_s$.
At $U_s=1$, an ideal ranker includes a target item in the list;
at $U_s=0$, a hit is impossible.
\end{proof}

\subsection{Mandatory decomposition of coverage losses}
\label{sec:eval-loss-diagnosis}

The presence of an eligible target is
\begin{equation}\label{eq:eval-available-hit}
V_s=\ind\{\mathcal I_q\cap Y_q\ne\varnothing\}.
\end{equation}
For sessions in $\mathcal S_J^+$,
$1\ge V_s\ge U_s(\pi)\ge\operatorname{Hit}_s(\pi)$, hence
\begin{equation}\label{eq:eval-loss-decomposition}
1-\operatorname{Hit}_s(\pi)
=(1-V_s)+(V_s-U_s(\pi))+(U_s(\pi)-\operatorname{Hit}_s(\pi)).
\end{equation}
The three disjoint cases are: no eligible targets; retrieval missed
eligible targets; retrieved targets were not delivered in the list.
The mean terms give an exact decomposition of $1-\operatorname{Hit@20}_J$.
Starting on $C_{\mathrm{smoke}}$, publish $N_J^{\mathrm{obs}}$, $N_J^+$,
$\sum_sV_s$, $\sum_sU_s$, $\sum_s\operatorname{Hit}_s$, and the three losses.

Conditional candidate coverage given an eligible target is
\begin{equation}\label{eq:eval-candidate-conditional}
\operatorname{Coverage}_{\mathrm{cand,avail}}
=\frac{\sum_sU_s}{\sum_sV_s},\qquad \sum_sV_s>0.
\end{equation}
A zero denominator means the metric is unavailable.
Report $\sum_sV_s/N$ alongside it to separate catalog coverage
from retrieval quality.

The purchase table distinguishes absence from $\mathcal I_G$
from ineligibility at request time. Its rows are not summed as session
losses: one session may contain purchases of different types.
Source failures, missing features, and deadline overruns
form separate slices, not a fourth loss term.
A match $i\in L_q\cap Y_q$ always means $\operatorname{Hit}_s=1$.

\subsection{Source contribution and overlap}
\label{sec:eval-source-contribution}

The first report uses already computed source results.
Let $\mathcal R_q^{(a)}$ denote the available portion of source $a$'s ordered
eligible-item list before quota $m_a$;
the depth and truncation of this list are recorded. If a source returns
only its post-quota response, pre-quota coverage is marked unmeasured.
Future targets do not participate in search; diagnostic lists
do not replace actual $\mathcal C_q$.
For a nonempty union, define
\begin{align}\label{eq:eval-source-overlap}
O_q^{ab}&=\frac{|\mathcal R_q^{(a)}\cap\mathcal R_q^{(b)}|}
 {|\mathcal R_q^{(a)}\cup\mathcal R_q^{(b)}|},\\
\mathcal R_q^{(a),\mathrm{uniq}}
&=\mathcal R_q^{(a)}\setminus\bigcup_{b\ne a}\mathcal R_q^{(b)}.
\nonumber
\end{align}
For an empty union, overlap is undefined. For each
source, publish the empty-response fraction, list sizes before
and after quota, and future-target coverage before quota, after quota, and in the final
union. Unique items and overlaps are shown separately.
Coverage of any family of measured lists $\mathcal D=(\mathcal D_q)_q$
uses the same session units:
\begin{equation}\label{eq:eval-source-coverage}
\mathsf H(\mathcal D)
=\frac1N\sum_s\ind\{\mathcal D_{q_s^{(1)}}\cap Y_s^{\mathrm{obs}}\ne\varnothing\}.
\end{equation}
Incomplete measurements are not replaced by empty lists; full pre-quota
coverage is then undefined. Overlap is averaged
over defined values with equal session weights;
the missing fraction is reported separately.
The sum of source coverages is not the coverage of their union.
Additional deeper retrieval is optional; its cost
is separated from the working configuration's latency.

To measure source contribution with candidates held fixed, let
$\mathcal O_q(i)$ denote their source set, including fallback filling.
Remove only items obtained exclusively from source $a$:
\begin{align}\label{eq:eval-source-removal}
\mathcal C_q^{-a}
&=\{i\in\mathcal C_q:\mathcal O_q(i)\setminus\{a\}\ne\varnothing\},\\
\Delta_a^{\mathrm{fixed}}
&=\frac1N\sum_s(U_s-U_s^{-a})\ge0.
\nonumber
\end{align}
Here $U_s^{-a}$ is computed by~\eqref{eq:eval-candidate-hit}
using $\mathcal C_q^{-a}$; freed slots are not filled.
A unique item for a request does not necessarily make a unique session contribution:
another source may have found another target item in the same session.
This measures candidate-coverage loss from removing a source from completed
results, not the performance of a reconfigured system.
The reconfigured system is tested in a separate predefined experiment:
redistribute quotas, perform fallback filling, retrain
the ranker on $B$, and compare policies under Appendix~\ref{app:ablation}.
That difference may have either sign. Retraining for every
source is not part of the mandatory first run.

\subsection{Choosing the next change}
\label{sec:eval-next-step}

A research hypothesis is formulated on $C_{\mathrm{smoke}}$ after
checking the implementation and loss table; the list of variants to compare
is fixed before opening $C_{\mathrm{select}}$.
Large $\sum_s(1-V_s)$ calls for analysis of the catalog and target eligibility;
large $\sum_s(V_s-U_s)$, of sources and quotas;
large $\sum_s(U_s-\operatorname{Hit}_s)$, of ranking
and timely list return. Technical causes are checked first.
Low candidate coverage alone does not prove weak retrieval:
target availability and~\eqref{eq:eval-candidate-conditional} must be considered.

The candidate bound uses knowledge of the future. The gap to the delivered list
does not guarantee an attainable ranker improvement and does not estimate
lost money or conversion. To compare ranking alone, hold
$\mathcal C_q$ fixed: a smaller gap caused by worse retrieval
is not an improvement. A new hypothesis, budget, and success criterion
are fixed before testing it; after inspecting $D$, a new final evaluation period is needed.

\subsection{Auxiliary metrics and slices}
\label{sec:eval-diagnostics}

For a request with $|Y_q|>0$, diagnostic recall is
\begin{equation}\label{eq:eval-request-recall}
 R_q@K=\frac{|L_q\cap Y_q|}{|Y_q|}.
\end{equation}
The denominator of the mean over requests with purchases is published separately.
To diagnose ordering, use
\begin{align}
 \operatorname{DCG}_q@K
 &=\sum_{j=1}^{|L_q|}
       \frac{\ind\{L_{q,j}\in Y_q\}}{\log_2(j+1)},
 \label{eq:eval-dcg}\\
 \operatorname{NDCG}_q@K
 &=\frac{\operatorname{DCG}_q@K}
 {\displaystyle\sum_{j=1}^{\min(K,|Y_q|)}1/\log_2(j+1)},
 \quad |Y_q|>0.
 \label{eq:eval-ndcg}
\end{align}
DCG weights are not measured probabilities of viewing positions.
Mean Recall and NDCG use the same $\mathcal S_J^+$.

Report separately coverage of purchased items inside the observed cart
and outside it. Cart state is determined only from the prefix;
it may be incomplete. Group coverages are not added: one session
may contain purchases from both groups.

Mandatory slices are specified before inspecting final results:
available history length, presence of earlier search,
cart presence, historical target-item popularity,
content-vector availability, and whether the snapshot knows the item.
Popularity-group boundaries use $A_G$, not frequencies from period $D$.

Purchases of unknown or ineligible items remain in targets
and denominators, with their fraction reported separately. First observation in the log
does not prove an item is new: its catalog introduction date is unknown.

\section{Uncertainty in historical comparisons}
\label{sec:eval-uncertainty}

First lists are evaluated on $\mathcal S_D^+$: sessions with a recorded
future purchase under Section~\ref{sec:eval-observed}. Below, $N=N_D^+=|\mathcal S_D^+|$,
and sums over $s$ use only this population. Intervals are conditional on
a purchase entering the log; they do not remove bias from incomplete telemetry.

\subsection{Paired session comparison}\label{sec:eval-paired}

Policies $\pi_1$ and $\pi_0$ are evaluated on the same sessions.
For $N>0$, the difference in first-list hit rates is
\begin{equation}\label{eq:eval-paired-diff}
 d_s=\operatorname{Hit}_s(\pi_1)-\operatorname{Hit}_s(\pi_0)\in\{-1,0,1\},\qquad
 \widehat\Delta_{\mathrm{Hit}}=
          \frac{1}{N}\sum_s d_s.
\end{equation}
During resampling, outcomes from both policies are carried together,
preserving their disagreement information for each session.

For a single fixed policy, estimate coverage itself:
\begin{equation}\label{eq:eval-single-policy}
x_s=\operatorname{Hit}_s(\pi),\qquad
\widehat H=\operatorname{Hit@20}_D(\pi)=\frac{1}{N}\sum_s x_s.
\end{equation}
Its interval uses the same calendar method, replacing $d_s$ with $x_s$
in~\eqref{eq:eval-block-stat}--\eqref{eq:eval-block-bootstrap}.
The half-width threshold $h_{\mathrm{single}}>0$ is set separately from
the paired-comparison threshold $h_{\mathrm{pair}}$; both are expressed as proportions.
For $N=0$, estimates are unavailable and the result is \code{insufficient}.

For $N>1$ and approximately independent sessions, the secondary standard error,
conditional on that assumption, is
\begin{equation}\label{eq:eval-paired-se}
 \widehat{\operatorname{se}}(\widehat\Delta_{\mathrm{Hit}})
 =\sqrt{\frac{1}{N(N-1)}
       \sum_s(d_s-\bar d)^2},\qquad
 \bar d=N^{-1}\sum_s d_s.
\end{equation}
The primary Coveo interval uses a temporal bootstrap of consecutive
blocks under Section~\ref{sec:eval-time-block}, subject to
Section~\ref{sec:eval-information}. Independent session resampling
\cite{bootstrap} is only a sensitivity check and does not replace
an unusable primary interval.

For replicate $b$, draw $N$ session indices independently and uniformly
from $\{1,\ldots,N\}$ and compute
\begin{equation}\label{eq:eval-bootstrap}
 \widehat\Delta_{\mathrm{Hit}}^{*(b)}
 =\frac{1}{N}\sum_{j=1}^{N}d_{I_j^{(b)}},
 \qquad b=1,\ldots,B_{\mathrm{boot}}.
\end{equation}
Paired first-list outcomes are carried together. Replicate count
$B_{\mathrm{boot}}=2000$ and the generator's initial state are fixed
in advance. Quantiles are defined in Section~\ref{sec:eval-time-block}.

\subsection{Planning precision for period \texorpdfstring{$D$}{D}}
\label{sec:eval-precision}

Before opening $D$, set the paired-difference interval half-width
$h_{\mathrm{pair}}>0$ in proportions and the planned variance
$v_{\mathrm{plan}}\approx\Var(d_s)$. Justify it using an authorized
pilot on $C_{\mathrm{smoke}}$ after training, without $D$ or future leakage.
The threshold, variance, assumptions, temporal-dependence multiplier, and pilot provenance
enter the immutable pre-$D$ evaluation plan; their absence does not block
\code{preflight} or the first training run. Under independent sessions,
\begin{equation}\label{eq:eval-precision-size}
 N_D^{+,\mathrm{iid}}=
 \left\lceil
 \frac{z_{1-\alpha/2}^{\,2}\,v_{\mathrm{plan}}}
      {h_{\mathrm{pair}}^2}
 \right\rceil,
\end{equation}
where $z_u$ is the level-$u$ standard normal quantile.
This follows from the half-width approximation
$z_{1-\alpha/2}\sqrt{v_{\mathrm{plan}}/N_D^+}$; power for a causal-effect
experiment is calculated separately.

For a single-policy estimate, use the same scheme with substitutions
$d_s\mapsto\operatorname{Hit}_s(\pi)$,
$h_{\mathrm{pair}}\mapsto h_{\mathrm{single}}$, and
$v_{\mathrm{plan}}\approx\Var(\operatorname{Hit}_s(\pi))$.
The primary estimate type is fixed in \code{inference.primary_estimate_type}.
A policy pair and paired-difference threshold are mandatory only for
\code{paired}; \code{single_policy} requires one policy,
$h_{\mathrm{single}}$, and its planned outcome variance.

Variance depends on policy disagreements, not only purchase count.
Zero variance in a small pilot does not justify $v_{\mathrm{plan}}=0$.
With insufficient data, use a safety margin or the rough bound $\Var(d_s)\le1$;
the calculation's provenance and assumptions are included in the protocol.

For temporal dependence, set a multiplier
$\kappa_{\mathrm{dep}}\ge1$ and plan
\begin{equation}\label{eq:eval-precision-dependent}
 N_D^{+,\mathrm{plan}}=
 \left\lceil\kappa_{\mathrm{dep}}N_D^{+,\mathrm{iid}}\right\rceil.
\end{equation}
On the pilot, divide the variance under time-block resampling
by the variance under independent sessions, bounding the result below by one.
With a zero denominator, the ratio is undefined. Block duration
and the safety margin are fixed in advance. The multiplier does not guarantee precision under a regime
change or few blocks; it does not correct the final interval
for independent sessions.
The dependence conditions in Section~\ref{sec:eval-time-block}
remain in force for the primary interval.

Period boundaries account for the forecast number of sessions with a recorded target,
calendar cycles, and exclusion of the last $H+g$ from the start-time window.
Session selection does not depend on target presence. After opening $D$,
boundaries and enrollment rules are not changed to obtain a convenient result.
Publish the actual interval, its half-width, and whether
$h_{\mathrm{pair}}$ or $h_{\mathrm{single}}$ is met.
Additional enrollment requires a new protocol.
For a fixed public log, \texttt{insufficient}
is an acceptable outcome: publish the available estimate and the reason for inadequate
precision without shifting periods or requiring unavailable additional data.
The chosen threshold $h_{\mathrm{pair}}$ or $h_{\mathrm{single}}$,
parameters $v_{\mathrm{plan}},\kappa_{\mathrm{dep}}$,
and their provenance and assumptions become mandatory run-manifest fields
at \code{pre_D}; actual period boundaries are fixed before training.
An unfilled field is not replaced by an invented measurement.

\subsection{Temporal dependence and interpretation limits}
\label{sec:eval-time-block}

Shared promotions, assortment, and demand create temporal dependence among sessions.
Before opening $D$, fix a dense calendar grid with step
$\delta_{\mathrm{grid}}>0$ and an origin anchored to the $D$ boundary.
For $D=[a_D,b_D)$, its $J$ cells cover only the eligible
start-time window $[a_D,b_D-H-g)$, excluding the final calendar buffer $H+g$.
This buffer does not confirm observation completeness.
An integer number of cells is required: boundaries and step are aligned before opening $D$;
unaligned boundaries invalidate the run manifest.
Each session is assigned in full to a cell by $t_{0,s}$. All cells within the window
are retained, including those without sessions with a recorded target;
no artificial tail is added. Let $\mathcal S_j^+$ contain sessions
from $\mathcal S_D^+$ in cell $j$. Record
\begin{equation}\label{eq:eval-block-stat}
 n_j=|\mathcal S_j^+|,\qquad D_j=\sum_{s\in\mathcal S_j^+}d_s.
\end{equation}
The moving-block bootstrap
for stationary dependent observations
\cite{kunsch1989} selects length
$\ell_{\mathrm{boot}}=m\delta_{\mathrm{grid}}$, $m\in\{1,\ldots,J\}$,
using an authorized pilot on $C_{\mathrm{smoke}}$ and domain constraints
without using $D$; block length is fixed before opening $D$.
For each replicate, independently sample
$\lceil J/m\rceil$ starting points uniformly from $\{1,\ldots,J-m+1\}$;
take $m$ consecutive cells from each starting point, concatenate
the resulting sequences, and retain the first $J$ indices
$J_1^*,\ldots,J_J^*$. Thus adjacency is preserved within
each selected block, and the two policies are always carried together.
Use the ratio of sums
\begin{equation}\label{eq:eval-block-bootstrap}
 \widehat\Delta_{\mathrm{Hit}}^*
 =\frac{\sum_{\ell=1}^{J}D_{J_\ell^*}}
              {\sum_{\ell=1}^{J}n_{J_\ell^*}}.
\end{equation}
Both numerator and denominator are recomputed in every replicate.
An unweighted mean of daily metrics with differing session counts has a different meaning.
The primary interval uses the empirical quantiles $\alpha/2$, $1-\alpha/2$
from $B_{\mathrm{boot}}$ replicates of~\eqref{eq:eval-block-bootstrap}.
For every bootstrap interval in this specification, the level-$p$ quantile
is the element at index
$\lceil pB_{\mathrm{boot}}\rceil$ in the sorted sample.
A zero denominator makes a replicate undefined: such cases
are logged, and interval computation receives \texttt{insufficient},
rather than silently discarding them or replacing them with zero.

\begin{assumption}[Validity of the temporal bootstrap]
\label{ass:eval-block-dependence}
Each estimated functional has its own vector of block sums.
For a paired difference, it is $(D_j,n_j)$; for a single coverage estimate,
$(\sum_{s\in\mathcal S_j^+}\operatorname{Hit}_s(\pi),n_j)$.
The corresponding joint process is sufficiently close to stationary
over the eligible start-time window, has finite second moments,
positive mean denominators, and dependence that weakens with lag.
Block length covers the substantial dependence horizon; distant
blocks are not linked by a dominant common shock. Asymptotic justification
requires block length to grow while remaining small relative to observation length.
A joint central limit theorem for component means
and consistency of the chosen bootstrap for the same limit are required;
qualitative stationarity alone does not prove these properties.
The estimated functional must be regular in the limiting means
and have nonzero limiting variation.
\end{assumption}
Coveo has no stable cross-session ID, so calendar blocks
do not prove that repeat-user dependence has been removed.
If another log provides such an ID and user dependence
dominates, user-cluster analysis is chosen before opening the test;
shared calendar shocks require additional treatment of temporal dependence.

Before $D$, the run manifest fixes the grid, block lengths, two additional
lengths for sensitivity, generator, and replicate count.
The necessary technical threshold is
$J_{\mathrm{eff}}=\lfloor J/m\rfloor\ge20$ for the primary calculation;
this is the number of available full blocks, not an estimate of independent observations.
For the fixed partition $\{1,\ldots,m\},\{m+1,\ldots,2m\},\ldots$
into $J_{\mathrm{eff}}$ full blocks, the number of blocks with $\sum n_j>0$
must be at least 20; overlapping bootstrap blocks
are not counted here as new observations. Empty intervals create no information.
For each functional and slice, count only blocks containing its
observations. Failure of the threshold or Assumption
\ref{ass:eval-block-dependence} gives \code{insufficient} for the corresponding
interval. Validity of the paired-difference process does not establish
validity of a single policy's coverage process, or vice versa.

Publish the time series $D_j/n_j$ for $n_j>0$, session counts $n_j$,
boundaries of promotions and assortment changes, and the primary and two predefined
sensitivity intervals. Empty cells are marked as having no
metric. A trend or substantial dependence of the conclusion on block length
limits interpretation; choosing the most convenient interval is prohibited.
Data from $D$ are not used to tune length, remove problematic days,
or replace the primary method. Too few blocks cannot be remedied
simply by increasing bootstrap replicates.
For interval bounds, report Monte Carlo error
using a method chosen in advance; replicate count does not determine
the statistical sufficiency of the original stream.

Both methods evaluate already trained models on this stream.
Uncertainty from retraining and seasonal changes requires several
temporal evaluations with repeated training; overlapping evaluations
are not independent.

\subsection{Information sufficiency and degenerate intervals}
\label{sec:eval-information}

The number of nonempty blocks does not protect against rare policy disagreements.
For the coverage being tested, set $c_s=d_s$ for a pair comparison and
$c_s=\operatorname{Hit}_s(\pi)$ for a single-policy evaluation.
Let $\mathcal B_b$ contain sessions from $\mathcal S_D^+$ belonging to a full
nonoverlapping block of $m$ cells under Section~\ref{sec:eval-time-block}.
For each category $a$ ($-1,0,1$ for a pair; $0,1$ for one policy),
compute
\begin{equation}\label{eq:eval-information-gate}
N_a=\sum_s\ind\{c_s=a\},\qquad
J_a=\sum_{b=1}^{J_{\mathrm{eff}}}
 \ind\{\exists s\in\mathcal B_b:c_s=a\},\qquad
\left|\{a:N_a\ge100,\ J_a\ge20\}\right|\ge2.
\end{equation}
These are initial protective thresholds: a single policy requires at least
100 hits and 100 misses, with each outcome present in at least 20 blocks.
A pair requires two sufficiently represented categories; therefore,
at least 100 disagreements and 20 blocks containing them. Rare categories
are not merged after inspecting $D$ to pass the threshold.

Failure of~\eqref{eq:eval-information-gate}, zero empirical variation,
an undefined replicate, or coincident primary bootstrap interval bounds
gives \code{insufficient}. The point estimate and all counts are published,
but zero width does not confirm precision or permit
$S_{\mathrm{offline}}=\mathrm{passed}$. Policy identity proven for every eligible input
is reported separately as an algebraic result
with proof and artifact hashes, not as measured precision.

For a usable interval $[l,u]$, additionally check
$(u-l)/2\le h_{\mathrm{pair}}$ or $h_{\mathrm{single}}$.
Observation-count, block-count, and half-width thresholds are fixed before $D$.
They exclude clearly uninformative estimates but do not guarantee coverage,
power, or absence of long-range dependence. An interval for rare
outcomes with different assumptions requires a separate advance plan;
it is not substituted for the primary interval after inspecting the result.

\section{Feature storage and updates}\label{sec:operations-features}

\subsection{Training--serving consistency}

Feature \(j\) is defined by a contract shared by training and serving:
\begin{equation}\label{eq:ops-feature-contract}
\mathsf F_j=(\mathrm{entity},\phi_j,\mathrm{sources},
\mathrm{available\_at},\mathrm{version},\mathrm{missing\_rule}).
\end{equation}
The feature version belongs to the pinned package. The package is chosen from
compatible versions published strictly before $t_q$, even if
$a_q>t_q$; within this set, the latest version in the publication order
recorded in the run manifest is selected. A version published between $t_q$ and $a_q$ is prohibited.
If no suitable package exists, \code{version_unavailable} is recorded;
the new package is not used even for fallback.

For a replayed request, verify
\begin{equation}\label{eq:ops-feature-parity}
\max_j\frac{\abs{x^{\mathrm{replay}}_{qj}-x^{\mathrm{serve}}_{qj}}}
{\eps_j+\abs{x^{\mathrm{replay}}_{qj}}}\le\delta_{\mathrm{feat}}.
\end{equation}
Categories, masks, candidates, snapshots, and tie-breaking order
must match exactly; the relative tolerance applies to floating-point
numbers. Parameters, numeric precision, and the index are identical; training transformations,
dropout, and encoder noise are disabled.
Tolerances for FP16 and BF16 are set separately from those for double precision.
The check includes the original ordered lists from each source,
availability masks, quotas, deduplication, and the selected branch;
agreement of the final feature vector alone is insufficient.

\subsection{Numerical stability}\label{sec:ops-numerics}

Equivalent mathematical expressions are implemented using numerically stable
primitives. For finite $x$, use
\begin{align}
\operatorname{softplus}(x)&=\max(x,0)+\log1p(e^{-|x|}),\\
\sigmoid(x)&=\begin{cases}
1/(1+e^{-x}),&x\ge0,\\
e^x/(1+e^x),&x<0.
\end{cases}
\end{align}
The pairwise logistic loss is computed as
$\operatorname{softplus}(-\beta_R d)$; binary cross-entropy is computed
from logits, without evaluating $\log(\sigmoid(x))$ or
$\log(1-\sigmoid(x))$. Sums of exponentials use
\begin{equation}\label{eq:ops-stable-lse}
\operatorname{LSE}(z)=z_*+\log\sum_j e^{z_j-z_*},
\qquad z_*=\max_j z_j,
\end{equation}
or \code{logaddexp} for blockwise accumulation. This rule applies
to attention, softmax, intent weights, contrastive loss, SID,
and catalog normalization. Masks are applied before taking the maximum. An empty admissible domain
is handled by a separate branch: the source is unavailable, and
$-\infty-(-\infty)$ is not evaluated. The value $-\infty$ is allowed only as an explicitly
defined marker for an excluded logit; it is not an item score.
Entropy and KL use stored log probabilities
as specified in~\ref{sec:modern-log-probabilities}: the logarithm of a value rounded to zero
is not computed; derivatives are checked as well.
Exact binomial bounds are computed using beta-distribution quantiles
with analytical boundary cases for zero successes and all successes;
tails use log-CDF, without summing factorials.

The first run uses FP64 for features, scores,
metrics, and statistical intervals. Exception: derivatives of the specified function computed in FP64
are passed to LightGBM through contiguous
FP32 arrays under the contract in~\ref{sec:model-ranker-implementation};
their tolerances are checked separately. Identifiers and counters are not
stored in FP16/BF16; counters use 64-bit integers
with overflow checks. In neural variants, FP16/BF16
are allowed for matrix operations only with at least FP32 accumulation;
norms, softmax/LSE, losses, and derivatives are accumulated in at least FP32,
and final metrics and statistics in FP64. The run manifest records
the numeric types at each stage, libraries, rounding modes, and tolerances.
Initial tolerances for the FP64 check in~\eqref{eq:ops-feature-parity}:
$\eps_j=10^{-8}$ and $\delta_{\mathrm{feat}}=10^{-6}$.
Any other numeric type requires a separate check against FP64 before model selection.

The norm of a finite vector is computed with scaling: when
$v_*=\max_j|v_j|>0$, first compute $w=v/v_*$ and $r=\sqrt{\sum_jw_j^2}$,
and the unit direction is $w/r$. The check $\norm v>\eps$
is evaluated as $v_*>\eps/r$, avoiding the potentially overflowing
product $v_*r$; when $v_*=0$, the zero branch is selected immediately.
Finite individual coordinates do not justify directly summing their squares.

An external modality and a computational error are treated differently.
An absent external vector receives mask 0 and reason
\code{vector_absent}; an incorrect dimension or nonnumeric/nonfinite
coordinates receive mask 0 and \code{vector_invalid}; a norm no greater than
$\eps$ receives mask 0 and \code{vector_zero}. All three cases
use a zero vector of the expected dimension. This transformation is identical
in preparation, training, and serving; other usable
modalities are retained. An error in a required ID, timestamp, or envelope structure
produces \code{invalid_input}, not an absent modality.
NaN/Inf values that unexpectedly arise after this preparation in features,
losses, derivatives, or scores stop training; during serving,
they invalidate the computation and activate the declared fallback.
Replacing an internal error with a zero mask is prohibited.
Tests cover logits $\pm100,\pm10^4$, empty masks,
zero norms, and the largest permitted groups. An FP64 reference,
finite permitted results, and unchanged masks are required;
FP64 absolute/relative tolerances are $10^{-10}/10^{-8}$,
and FP32 tolerances are $10^{-6}/10^{-5}$. Reordering nearly equal
scores is not hidden by a numerical tolerance: the common tie-breaking order applies.

\subsection{Session state updates}

State \(\mathsf U_s(t)\) contains a bounded recent history,
action counters, the last event time, observed cart
state, and search context. For a newly available block \(E_{s,t}\),
\begin{equation}\label{eq:ops-state-update}
\mathsf U_s(t)=\mathcal T\bigl(\mathsf U_s(t^-),E_{s,t};\nu_{\mathrm{features}}\bigr).
\end{equation}
\(\mathcal T\) is deterministic and depends only on the available history.
Redelivery is handled under the contract in
Section~\ref{sec:ops-state-storage}; matching Coveo rows
do not replace a transport event identifier.

For a decaying aggregate of a single action type, a convenient
recurrence is
\begin{equation}\label{eq:ops-decay-update}
x_{si}(t)=e^{-(t-t')/\tau_x}x_{si}(t')+
\sum_{e\in E_{s,t}}\ind\{i(e)=i,\ \kappa(e)=\kappa_x\}.
\end{equation}
Under sequential processing, this is a sum of decaying past events.
An exact sliding window requires an event queue.

Storage retains the last \(L_{\mathrm{hist}}\) general blocks for replay,
a separate queue of the last $L_H$ \code{detail} events for the graph,
the latest admissible search context, and the necessary aggregates
over the entire visible prefix. The general tail does not replace the view queue:
frequent searches must not evict the $L_H$ events that are still needed.
Memory for all queues, aggregates, and delivery markers is accounted for separately.
Removing a block does not reset
aggregates over the entire visible session. Operational state is removed
after inactivity once the required outcome records have been written; state and outcome
retention periods are specified separately.

\subsection{State storage contract}
\label{sec:ops-state-storage}

Process memory is sufficient for replay; serving requires
a key-value store with atomic record updates. No product is prescribed;
retention periods and delivery limits are recorded in the test description.
\begin{enumerate}
\item \textbf{Record key and contents.}
The key $(\nu_{\mathrm{source}},s)$ includes the source identifier
namespace. The record contains $\mathsf U_s$, revision number $v_s$,
version $\nu_{\mathrm{features}}$, the last arrival time,
a completeness flag, and markers for accepted events.
Purchases and the request log are stored independently of the operational record.

\item \textbf{Concurrent access.}
Verifying that the read value $v_s$ is unchanged and writing the new state,
$v_s+1$, and event markers are performed as a single atomic operation.
If the revision differs, the record is read again and the transition is recomputed
within the overall response deadline.
A request uses one confirmed state version and one package
$\mathsf V$, without mixing fields. The snapshot includes only events available to the request
(Section~\ref{sec:time}); if the record contains later
events, the snapshot is reconstructed from the log.
The transformation version matches the package; changing it requires
a declared state transformation or reconstruction.

\item \textbf{Redelivery.}
The block receipt, ID check, and counter/completeness updates are atomic.
Receipts and event markers are retained for at least $\Delta_{\mathrm{retry}}$;
a verified duplicate does not change the logical state, including completeness.
After this period, IDs are checked against the log or the response is
\code{retry_unverifiable}, without treating the delivery as a new event.
The order of transport checks is specified in~\ref{sec:ops-transport-block}.

\item \textbf{Retention and lateness.}
The lifetime after the last arrival,
$T_{\mathrm{state}}$, the lateness limit $\delta_{\mathrm{late}}$,
and outcome retention are specified separately. The tail is bounded by block count and time.
After a record is removed, a late event requires reconstruction
from the log or an incompleteness marker.

\item \textbf{Failure and recovery.}
An unavailable, incompatible, or incomplete record activates the prescribed fallback
with a reason code; unread history is not treated as empty.
If an admissible fallback or a required write is not ready by the deadline,
the result is a technical failure.
Snapshot reconstruction, version checking, and retrying an atomic write
are included in the state budget and are not performed after it expires.
Initial limits: one atomic-update retry,
at most 1000 replayed events, and at most 1 MiB of log data
per request. Exceeding any limit produces
\code{state_replay_limit}; background reconstruction does not delay the response.

\item \textbf{Cached computation results.}
The reuse key depends on the visible prefix,
request time, and transformation package:
\[
(\nu_{\mathcal H_q},t_q,\mathsf V).
\]
Identifier $\nu_{\mathcal H_q}$ identifies the session, visible events,
and search. History and intent weights also depend on time and the model;
a session identifier alone is insufficient for reuse.
\end{enumerate}

\subsection{Requirements for request replay}
\label{sec:ops-replay-trace}

The run manifest specifies one of two modes. \code{ideal_deterministic}
uses $a=t$, frozen data, and normally operating sources; this mode
does not reproduce observed operational failures.
In this mode, evaluations on $C_{\mathrm{smoke}},C_{\mathrm{select}},E,D$
do not require a service implementation or the \code{post_load} stage.
The report states the mode and its assumptions; local time and memory measurements
do not establish operational guarantees.
\code{execution_trace} additionally stores an evaluation trace:
$t_q,a_q$, the package, a reference to the immutable visible prefix/state snapshot,
the selected branch and delivery status, source availability masks and
unavailability reasons, truncation decisions, the seed, and the fault-injection schedule.
The reference includes a hash, but the hash itself does not replace retained inputs.

The trace records the actual number of steps/nodes when search stops
early; if this does not uniquely reproduce the list, that source's list
is stored. Trace completeness and input availability are checked before computing
$\mathcal C_q^{\mathrm{served}}$ and $U_s$. A healthy rerun
does not replace a lost trace: the corresponding candidate set
is marked unmeasured. Trace access and retention follow
Section~\ref{sec:ops-security}; expiration of input retention ends
the ability to replay exactly, but does not invalidate retained aggregate estimates.

\subsection{Arrival time and late events}
\label{sec:ops-event-time-validation}

Input time is an integer number of UTC milliseconds since the Unix epoch
within the signed 64-bit range. Before production operation, the run manifest specifies
the allowed calendar range $[T_{\min}^{\mathrm{in}},T_{\max}^{\mathrm{in}})$;
it must cover the intake period and permitted lateness.
An incorrect type, unit, or range produces \code{invalid_input};
the value is neither rounded nor corrected based on an assumed unit.
A verified duplicate of an accepted block is handled before rechecking
time bounds, and its original timestamps are not changed.
For a new block, before modifying the history, verify
\begin{equation}\label{eq:ops-event-time-window}
-\delta_{\mathrm{late}}\le t(e)-a_{\mathrm{block}}
\le\delta_{\mathrm{future}}.
\end{equation}
$a_{\mathrm{block}}$ is the server time of the atomic decision.
In the initial profile, $\delta_{\mathrm{future}}=0$: there is no waiting buffer.
Violating the upper bound causes quarantine with reason \code{timestamp_future},
and violating the lower bound gives \code{timestamp_late}. The block does not change history,
$t_{\mathrm{last},s}$, or closed time groups. A new trusted
block marks incompleteness as specified in~\ref{sec:ops-transport-block};
fallback remains active until recovery. An error from an untrusted sender
does not change another party's session. Positive clock-ahead tolerance is prohibited in this profile.

An accepted late event changes only responses after its arrival.
State is reconstructed from a checkpoint preceding the tail
and tail events in timestamp order. The checkpoint contains all aggregates and queues;
without it, the complete log is required within budget; otherwise, state is incomplete and fallback applies.
The independent purchase log continues recording outcomes; past responses are unchanged.
Delays, duplicates, and out-of-order delivery are tested separately
from the main historical mode \(a=t\).

\subsection{Partial data availability}\label{sec:ops-partial-data}

A missing feature, a disabled source, and a read failure have
different codes. An item modality that is absent uses its standard
mask if this mode was included in training. A storage failure
is not encoded as an absent modality: a source is disabled only
under a preregistered degradation rule; an unregistered
combination of unavailable sources activates fallback. A partial source list
is allowed only with an explicit completeness flag and the prescribed quota;
an arbitrary portion of a list at deadline expiry is not considered a complete result.

The eligibility catalog is mandatory: an item is served only
after its eligibility is confirmed in the request's version. An older copy
is allowed only within the validity period pinned in the
run manifest. Unknown eligibility excludes an item from both the primary and
fallback branches; inability to verify even the fallback produces
a technical failure. Unread history is not interpreted as a new session.
For each degradation mask, report frequency, response length,
quality, and the number of technical failures.

\subsection{Trust boundaries, input limits, and trace retention}
\label{sec:ops-security}

The threat model includes forged/bulk events, Sybil sessions,
SKU popularity inflation, expensive history reconstruction, and artifact substitution.
Source authentication confirms origin, not the truthfulness
of behavior. Required schema and size checks precede computation;
unusable external vectors are handled as specified in~\ref{sec:ops-numerics}.

\paragraph{Sealed transport block.}\label{sec:ops-transport-block}
The initial production-service profile is \code{sealed_block_v1}.
Each session namespace has one registered
aggregator responsible for combining behavior and search before sending a block.
The ownership registry is fixed before launch; authorization is checked against
the authenticated sender, not the submitted namespace name.
One immutable envelope contains the source namespace, session ID,
event time, a unique \code{block_id}, event count, stable event
IDs, the events themselves, the payload's SHA-256, and \code{closed=true}.
All events have the same $s,t$; the event count equals the lengths of the payload
and ID list, and \code{event_ids[k]=events[k].event_id} for every $k$.
IDs within the envelope are unique. The service computes the hash from the exact bytes of the immutable
payload. The receipt binds the namespace, session, time, block ID,
event count, and this hash; a matching ID alone is insufficient.
A different serialization is not an identical transport retry.
The service accepts only a complete envelope and does not combine fragments
from different requests. An incomplete envelope adds no events;
its rejection for an established trusted session is accounted for
under the completeness rules below.
Receiving the complete envelope and queueing are included in the overall deadline from the first
byte under Section~\ref{sec:ops-deadlines}. A separate 10 ms limit
does not apply to the demonstration service. Aggregator assembly is reflected
in $a(e)-t(e)$ and the separately measured event-to-response latency.
The service records $a_{\mathrm{block}}$ at the atomic decision on a verified
block, including quarantine decisions. For accepted events, $a(e)=a_{\mathrm{block}}$;
rejected events are not included in available history.
Sender time does not replace these timestamps.
When ingestion and a request are combined, the cutoff $a_q$
is recorded after block acceptance; the response deadline is still measured
from the first byte. Recovering retries does not change the original $a(e)$.

Limits: 1 MiB per envelope, 256 UTF-8 bytes per SKU,
1 to 100 raw events per envelope, 100 logical events per block,
1000 logical events per session, and 1000 search-list entries.
Logical counters are determined after deduplicating event IDs
and applying the selected canonicalization. Before accepting a complete block,
the sum of the previous counter and its new logical events is checked; crossing
the limit rejects the entire new block, without partial truncation.
Unused search lists are removed after input validation.

Already closed $(\nu_{\mathrm{source}},s,t)$ groups are neither extended nor
reopened. The same \code{block_id} with different contents, or new IDs
for a closed timestamp, produces \code{closed_block_conflict}, not admissible
lateness. The new unknown portion is quarantined; accepted history
is preserved. Ordinary late delivery is allowed only for a block
not yet accepted and follows~\eqref{eq:ops-event-time-window}.
The historical-data adapter groups the entire known time block
before creating an envelope; the production service does not assume knowledge of future packets.

The rate limit is a token bucket of 10 raw events/s with burst 100;
the cost of the full attempt, including retries, is checked atomically.
Cost is determined by the actual raw event count, not the header;
an empty, rejected, or unparseable attempt costs at least 1.
Connection and sender limits apply before the session is established.
The following distinct transitions then apply:
\begin{enumerate}
\item A verified complete duplicate of an accepted block changes neither counters,
completeness, nor responses. If tokens are insufficient, it may be rejected as
\code{duplicate_rate_rejected}; this is not a history gap. The rule also applies
after $\delta_{\mathrm{late}}$, while the receipt remains verifiable.
\item An envelope containing only known events does not create a new
block. A partial duplicate with new IDs is not handled event by event:
for an already closed timestamp it is a conflict; for a new block, it violates
the mapping between an event ID and its original fields. New unverified events
are not considered accepted; existing records are unchanged.
\item A rejected new block from a trusted aggregator leaves
a gap and incompleteness marker; subsequent requests use fallback
until recovery. An unverifiable retry receives
\code{retry_unverifiable}, without counting the events again;
reconstruction or log verification determines completeness.
\end{enumerate}
Receipt verification is bounded by one indexed read and the
envelope size. If storage is unavailable, the state fallback contract
applies; authentication failure does not change another party's session.
Quarantine of new events is replayed offline with a gap marker:
graph transitions across an unknown segment are prohibited. Input limits
do not use the future session length, change assignment, remove
an admitted session from evaluation/ITT, or disable the independent purchase log.
They do not replace the recommendation-request limit $Q_{\max}$.

Tenant limits on events/s, bytes/s, active sessions, and the daily contribution
of a single SKU must have numeric values before launch. Reaching
a limit triggers throttling or quarantine with a reason code.
For the graph, at most one session contribution to the support
of each pair is retained; for popularity, at most one session contribution per SKU;
additional SKU/source limits are registered as part of
the data transformation. Spikes in edge support and popularity are checked against
previous windows; automatically incorporating suspicious events
into a new package is prohibited until quarantine review. Sybil load,
a single-SKU attack, an oversized block, and mass history reconstruction are mandatory
tests. These measures limit impact but do not guarantee
identification of all fake users. Unknown operational
limits mean the production profile is not ready, not that intake is unlimited.

The minimum required response record contains a pseudonymous session
key, versions, the selected branch, served item identifiers,
the outcome, and technical status. Full candidate lists, scores, and detailed
reasons are stored only in a bounded diagnostic log;
raw queries, vectors, and history are excluded from the ordinary log.
The pseudonymous key is produced by a keyed transformation with namespace
separation; hashing an identifier is not considered anonymization.
For A/B linkage, the assignment key remains stable throughout the
experiment horizon, while diagnostic-log keys rotate independently.

Initial retention policy: state is deleted after 24 hours
of inactivity; $\Delta_{\mathrm{retry}}=24$ hours,
$\delta_{\mathrm{late}}=30$ minutes; the replay log is retained for 48 hours,
and detailed traces for 7 days. Minimal assignment and outcome records
are retained until the prespecified analysis closing date plus 30 days,
then deleted or replaced with aggregates without session keys.
Redelivery markers live independently of state retention and for at least
the retry period. Retention periods and storage owners are included
in the run manifest; these values are a project profile, not
a claim about legally required retention. Access is role-based,
transmission and storage are encrypted, and diagnostic-trace reads are audited.
Deletion covers state, the replay log, diagnostic copies,
and backups under the declared expiry period; recovery
does not restore deleted records without reapplying the deletion registry.
Public reports and user-facing explanations do not expose
session identifiers, raw queries, or action sequences.

\section{Computational budget and serving}\label{sec:operations-serving}

\paragraph{Scope of the demonstration version.}
For a portfolio, one local process with a verified package,
state, repeatable responses, fallback, explanations, and a log is sufficient.
Time boundaries, item eligibility, input limits, and the prohibition
on late publication remain in force. Production load
(Section~\ref{sec:ops-load-profile}), RTO/RPO and recovery of multiple
instances (Section~\ref{sec:ops-recovery}), production signatures, and monitoring
are outside its required acceptance criteria and require a separate task.
The run manifest's production-profile \code{serving}/\code{operations} fields,
\code{post_load}, and $S_{\mathrm{serve}}$ belong to this deferred
operational acceptance, not the demonstration.

\subsection{Preparation, training, and request processing}
\label{sec:ops-compute-stages}

Version preparation, training, and serving are measured separately.
Preparation for the first run creates the graph, content vectors,
and lookup tables; extensions add paths, learned vectors,
approximate nearest-neighbor (ANN) indexes, and the SID tree.
Training transformations and code noise are used only during training.
A request uses prepared paths and item vectors, computing the session
state, candidates, features, and ranking. Item-feature search
in the first run is exact; shortest paths are not recomputed.
Autoregressive SID decoding belongs to the serving budget
of a separately enabled research variant.

\subsection{Version contents}

A consistent package is published:
\begin{equation}\label{eq:ops-package}
\mathsf V=(\nu_D,\nu_{\mathcal I},\nu_G,\nu_{\mathrm{vec}},
\nu_{\mathrm{features}},\nu_f,\nu_{\mathrm{policy}},
\nu_N,\nu_{\mathrm{SID}}).
\end{equation}
The graph, vectors, and ranker are replaced atomically.
Vector-space compatibility is checked by provenance,
not dimension alone.

Version $\nu_N$ combines the indirect-link topology, encoder weights,
the learned item-vector table, and its search index.
Version $\nu_{\mathrm{SID}}$ contains the encoder, codebooks,
decoder, code-to-item mapping, and valid-prefix tree.
A disabled module receives version \code{not_used};
in the first run, this applies to $\nu_N$ and $\nu_{\mathrm{SID}}$.
Replacing a codebook without updating its corresponding tree
and decoder makes the package incompatible.

The package manifest contains the SHA-256 of every artifact, exact library
versions and a lockfile, the build commit, data cutoff and publication time,
builder identifier, and transformation dependencies. In the production service,
the manifest is signed by a trusted publisher (initial profile:
Ed25519); the trust key comes from a separate trusted store,
not the package being verified. Before activation, the signature,
all hashes, schema compatibility, and version authorization are checked. Checking a
hash without authenticating the publisher is insufficient. Artifacts are read
by a declarative loader without executing serialized code;
paths outside the package, unsupported types, and sizes exceeding declared
bounds are rejected. A local research package may use
the \code{local_checksum_only} profile, which explicitly does not verify
the authenticity of a production package. A package that fails verification is not activated.

A request follows this sequence:
\begin{enumerate}
\item validate input structure and pin the published package;
\item obtain the eligible catalog, verify the state's transformation
version, and update the state atomically;
\item for an enabled neural variant, compute the session
representation and search intent weights once;
\item retrieve candidates and their provenance;
\item compute features and scores;
\item apply final-list constraints;
\item record the response, latency, and execution status.
\end{enumerate}
Fallback is included in \(\pi\) at its actually observed usage rate.

\subsection{Memory and operation counts}

Let \(n_I\) be the number of items, \(|\mathcal R_G|\) the number of graph relations,
\(d_G\) the retained neighbor count, \(d_T,d_V\) the dimensions of
content vectors, \(p\) the number of features, \(T_f\) the number of trees,
and \(D_f\) their maximum depth. The main costs are of order
\begin{align}
\mathrm{Mem}_{G}&=O(|\mathcal R_G|n_Id_G),\label{eq:ops-graph-memory}\\
\mathrm{Mem}_{\mathrm{content}}&=O(n_I(d_T+d_V)),\label{eq:ops-content-memory}\\
\mathrm{Work}_{\mathrm{graph},q}&=O(|\mathcal R_G|L_Hd_G),\label{eq:ops-graph-work}\\
\mathrm{Work}_{\mathrm{rank},q}&=O(Mp+MT_fD_f).\label{eq:ops-rank-work}
\end{align}
Exact content search over the catalog requires
\(O(n_I(d_T+d_V))\) scalar operations per request.
If needed, an approximate index is compared with exact search
before opening $D$. A complete graph over \(n_I^2\) pairs is not materialized.

\paragraph{Costs of optional extensions.}
Item vectors of dimension $d$ require $O(n_Id)$ memory
and are recomputed before the index is published. Session representations
are computed once for candidate sources and ranking features.

Indirect-edge preparation is bounded by path length, the number of expansions,
and the number of retained neighbors. Deterministic early termination
is part of the model; its frequency across vertices is reported.

Exact normalization over the catalog requires $O(n_Id)$ operations per request;
an ANN index does not accelerate it. Training also stores intermediate values
for derivatives. Blockwise log-sum-exp accumulation preserves
the exact loss. Negative-item sampling defines a different procedure
(Appendix~\ref{sec:sampled-overview}) and does not eliminate item-representation
recomputation or global regularizers.
Total training time and peak RAM and accelerator memory are measured.
Exhausting the budget is a computational failure;
the number of epochs is not changed without declaration.

\paragraph{ANN acceptance contract.}\label{sec:ops-ann-contract}
Each approximate source has its own quality measured as
\begin{equation}\label{eq:ops-ann-recall}
\operatorname{IndexRecall}@m=\operatorname{Agreement}@m
\end{equation}
from~\eqref{eq:model-ann-agreement} on a fixed probe sample.
This measures agreement of neighbor-search results, not purchases.
Exact and approximate search use identical vectors, filters,
quota $m$, and tie-breaking order. Approximate results are rescored
using exact dot products, deduplicated,
and checked for eligibility; the original fractions of ineligible and duplicate
results are reported, and both are zero in the final list.
Deterministic search expansion after filtering has a specified
limit; exhausting it does not permit removing the filter or increasing $M$.

Additionally, compute set precision with numerator
$|C_q^{\mathrm{approx}}\cap C_q^{\mathrm{exact}}|$ and denominator
$|C_q^{\mathrm{approx}}|$, exact top-1 agreement, and score regret
\[
\operatorname{Regret}_q=
\max(0,r_q(i_q^{\mathrm{exact}})-r_q(i_q^{\mathrm{approx}})).
\]
Here $r_q(i)$ is the dot product of the original query
and item representations for this source, before temperature division and
other scale factors; their norms do not exceed 1. For content
search this is cosine similarity; for neural search, the numerator of its score formula.
The run manifest records \code{regret_score=unscaled_bounded_dot_v1} and references
to both representation formulas. For nonempty responses, regret lies
in $[0,2]$ and is invariant to positive scaling of the ranking score.
When the exact response is nonempty and the approximate response is empty, precision, recall,
and top-1 agreement are zero, while regret is $+\infty$;
such empty responses are therefore prohibited on control requests.
An empty exact response is counted separately and requires an empty
ANN response under the same masks and filter. Stratification is fixed
by available-modality masks and the popularity of the exact top-1:
the popular group contains the first $\lceil0.1|\mathcal I_G|\rceil$ items
ordered by decreasing session count in $A_G$, with ties broken by ID;
the rest form the second group. The boundary is fixed before ANN validation.

Initial engineering thresholds: mean IndexRecall and precision
at least $0.98$, top-1 agreement at least $0.99$, and mean regret
at most $0.01$. In every predeclared nonempty segment,
recall and precision must be at least $0.95$, top-1 agreement at least $0.95$,
and regret at most $0.02$. At least 10\,000 nonempty exact
responses overall and 1000 per segment are required; a smaller sample means
\code{insufficient}, not a passed segment. Across ten repetitions
of identical inputs, lists after exact rescoring must match exactly;
the initial random-generator state, index, thread count,
and tie-breaking order are included in the run manifest.
These are empirical thresholds on a fixed control sample from
$C_{\mathrm{smoke}}$, not population guarantees.
Index parameters are selected before $C_{\mathrm{select}}$.

ANN must also satisfy the common P95/P99, failure, and fallback-response profile
and the total memory limit in Section~\ref{sec:ops-load-profile}, including the index.
On $C_{\mathrm{select}}$, the change in Hit@20 relative to exact search
is reported on the common $\mathcal S_{C_{\mathrm{select}}}^+$.
The allowable decrease $\delta_{\mathrm{ANN}}$ is specified in Hit@20 units
before comparison; without a specified tolerance, ANN is excluded from selection.
This selection filter does not replace the final interval on $D$.
If ANN is disabled, its status is \code{not_used}.

\paragraph{Training-example table.}
The candidate table size, up to \(N_{\mathrm{train}}Q_{\max}M\) rows,
is estimated before launch. If requests or negative candidates are sampled,
inclusion probabilities and the loss are specified; the evaluation catalog
is not reduced to save computation.

\subsection{Latency and fallback branches}
\label{sec:ops-deadlines}

Let \(\ell_{qj}\) be the duration of a sequential stage,
and \(\ell_q^{\mathrm{wait}}\) the waiting time for service. Then
\begin{equation}\label{eq:ops-latency}
\ell_q=\ell_q^{\mathrm{wait}}+\sum_{j=1}^{J}\ell_{qj}.
\end{equation}
For thresholds \(b_j\) and \(b_0\), without assuming independence,
\begin{equation}\label{eq:ops-latency-bound}
\Prob\!\left(\ell_q>b_0+\sum_j b_j\right)
\le \Prob(\ell_q^{\mathrm{wait}}>b_0)+\sum_j\Prob(\ell_{qj}>b_j).
\end{equation}
Stage percentiles do not add up to a percentile of total latency.
Request rate, concurrency, catalog, and the fraction of warm states are fixed.

The fallback list is prepared at the start of the shared candidate procedure.
If the primary branch exhausts its budget before the overall response deadline,
an already prepared eligible branch is used, reserving time
for logging and return; if none exists, an empty list is returned
with a reason code. Reaching the overall deadline does not start new computation.
Evaluation does not remove such a request. Fallback forms a separate branch
$\mathcal C_q^{\mathrm{fallback}}$ and does not extend
$\mathcal C_q^{\mathrm{main}}$. When a ready branch is selected with time reserved
for logging and return, \code{branch} and
$\mathcal C_q^{\mathrm{served}}$ are recorded; mixing the two sets
to improve the candidate bound is prohibited. Always,
$L_q\subseteq\mathcal C_q^{\mathrm{served}}\subseteq\mathcal I_q$
and $|\mathcal C_q^{\mathrm{served}}|\le M$. If no ready branch
is selected, the active set is empty. If the deadline or mandatory write
is not satisfied after selection, $L_q=\varnothing$ and
\code{delivery_status=error}, but the recorded candidate set
is retained: at the session level, delivery loss remains in
$U_s(\pi)-\operatorname{Hit}_s(\pi)$.

For the demonstration, the initial overall timeout
$b_{\mathrm{deadline}}=1000$ ms is fixed before launch and covers
the entire response, including the mandatory write. This is a waiting limit,
not a performance guarantee under load. Individual stages have no hard
deadlines; the reserve for logging and returning fallback is an execution parameter,
not a separate acceptance threshold. The production profile has separate
requirements in Section~\ref{sec:ops-load-profile}.
A monotonic clock enforces the overall deadline with time reserved for returning fallback.
A noninterruptible operator is bounded
in advance or isolated; asynchronous cancellation does not guarantee completion.
Each operation uses the remaining overall time after accounting for the still-required
write and return; no new independent stage deadline is created.
Retries share the same budget. After a branch switch,
a late result from a cancelled operation does not change the response or overwrite
state; a confirmed state write remains valid even
if the response fails. Unfinished operations are cleaned up in an isolated
pool with a fixed limit; new calls are rejected when it
is full. Separate codes are \code{invalid_input},
\code{overload_rejected}, \code{backend_timeout}, \code{cancelled},
\code{state_replay_limit}, and \code{logging_failure}; a single aggregate
\code{error} does not replace these reasons.

SID runs only when sufficient time remains in the overall candidate-retrieval
budget. The step limit and search width are fixed.
When time is exhausted, declared quotas and fallback rules apply,
without a hidden increase in $M$. SID skip frequency and the quality change
due to the time limit are reported separately.
If fallback is not yet ready or the mandatory write is incomplete,
a technical failure code is returned; the request remains
in the evaluation denominators.

The demonstration log links the request, package, candidates, public
response, and outcome. Two-phase commit and RPO=0 are not required; a write failure
or unknown delivery is not declared a successful completion.
Acceptance checks normal primary serving, replay without recomputation,
an actual computation timeout with fallback, and failure with incomplete state.
The HTTP response is checked against the log. Total latency, load, and memory
are published as measurements; P95/P99 and fallback rate have no mandatory
thresholds for the demonstration. Normal CPU-core scheduling is retained.

\subsection{Reproducible load profile}
\label{sec:ops-load-profile}

Deferred production profile: applies only under a separate task
and does not block the demonstration API. Its requirements have not yet been verified by testing.
\begin{center}
\begin{tabular}{p{0.42\linewidth}p{0.51\linewidth}}
\toprule
Parameter & Initial value\\
\midrule
Compute node & 8 virtual CPU cores, 32 GiB RAM,
no GPU during serving\\
Offered load & 100 requests per second\\
Load generator & Response-independent schedule;
at most 32 outstanding generator requests\\
Warm-up & 5 minutes before measurement\\
Measurement & 30 minutes, three separate repetitions\\
Service latency target & P95 $\le50$ ms, P99 $\le100$ ms\\
Total peak RAM of service processes & At most 28 GiB\\
Response deadline & 120 ms from arrival at the service\\
Allowed technical failures & At most $0.1\%$ of all scheduled requests\\
Allowed fallback-response rate & At most $1\%$\\
\bottomrule
\end{tabular}
\end{center}
The report specifies the CPU model, processes and threads, library versions,
numeric type, catalog, index parameters, and memory. Equal core counts
do not guarantee equal performance; GPU execution is tested separately.

\paragraph{Overall memory limit.}\label{sec:ops-memory-limit}
The RAM limit applies to all variants, including exact search,
indexes, local state, queues, and concurrently loaded packages
during updates/rollback. The initial metric is the maximum summed PSS of service
processes: shared pages are accounted for proportionally.
If PSS is unavailable, a conservative sum of peak RSS across
all test processes is allowed; shared pages counted multiple times are not subtracted
using an unverified estimate. The selected measurement tool, sampling interval, and
process set are fixed before testing and are identical across compared
variants. Exceeding the tabulated limit means \code{failed}, regardless of ANN.

Distributions of prefix length, search presence, missing values, and repeat
requests come from $C_{\mathrm{smoke}}$ inputs or a preapproved
production profile. Outcomes from $D$ are not used to select load.
Session states reproduce the declared transitions;
repeating the same warm request does not replace a request stream.
Cold-cache startup and recovery after failure are measured separately.

The generator schedules requests every 10 ms independently of responses;
inputs and initial state are identical for all variants.
A burst profile is declared separately. If the limit of 32 prevents sustaining
100 requests per second, undelivered requests and the actual
rate are recorded; a slowed-down test does not establish the claimed load.
Server concurrency and queue limits are specified separately from the generator limit.
For the initial node, the server allows at most 16 active requests
and 16 waiting requests; input reception and queueing together are limited to 10 ms.
When full, the queue does not grow: the service returns
\code{overload_rejected} without an expensive state read.
Limits are global to the node, not independently repeated in each
process. Rejection and its resulting empty response are included
in operational denominators; an assigned session remains in ITT.
A recommendation request may be retried only with the same idempotency key
and within the original deadline; a server response does not trigger
an unlimited automatic retry loop. Event delivery follows
Section~\ref{sec:ops-transport-block} and does not restart an old response deadline.

Service latency is measured from the first request byte at its ingress
to readiness of the complete response, including waiting, feature retrieval,
and the mandatory write. External network and client processing belong
to a separate end-to-end latency. Scheduled send time,
actual send time, service arrival, and completion receive
separate timestamps; generator delay is not hidden
inside an incomplete measurement.

Requests are selected by scheduled send time and observed
until completion or deadline, including those unfinished at the end of the window.
For acceptance percentiles, failure, nondelivery, or deadline expiry
receives latency $+\infty$; actual times to error are retained.
A fallback response uses its measured latency and branch marker.
Percentiles over successful primary responses alone are supplementary diagnostics.
The empirical quantile at level $p\in(0,1]$ for $n>0$ values
is defined as $x_{(\lceil pn\rceil)}$ after sorting, including $+\infty$.
For an empty window, rates and quantiles are undefined: \code{insufficient},
not zero. The status log is reconciled with an independent ingress counter and
the generator schedule; missing records count as failures and are not removed
from the denominator. This rule also applies to operational windows $W$.

Each of the three repetitions must meet latency,
failure, and fallback-response targets; pooling repetitions must not hide
a failed run. Throughput,
maximum queue length, actual concurrency, and CPU/RAM usage are reported separately.
Load acceptance precedes the A/B test; during the experiment,
latencies and failures continue to be monitored as guardrail metrics.

\subsection{Recovery and rollback}\label{sec:ops-recovery}

Before publication, the last verified package version
and its state schema are retained. The initial rollback RTO requirement is 60 s
from the operator/guardrail-controller decision to readiness of the verified
version on all instances; RPO for confirmed assignment
and outcome writes is zero. These are targets tested through fault injection,
not reported measurements. Each request is pinned to one version;
rollback does not mix artifacts or change previously stored responses.
Incompatible state is reconstructed or migrated under a
verified rule; eligible fallback applies until it is ready.
Returning to an old version with an expired catalog is prohibited.
Loading two versions is included in peak memory; if this is impossible,
the switch is staged with the same RTO and availability checks.
Testing covers process loss, unavailable state,
mandatory-write failure, package corruption, and rollback to a different
feature version. Measured RTO/RPO and failed-request codes
are reported separately from steady-state load results.

\section{Monitoring quality and data changes}\label{sec:operations-monitoring}

\subsection{Immediately observable metrics}

For a request window \(W\), define
\begin{align}
\widehat e_W&=\frac1{|W|}\sum_{q\in W}\ind\{\mathrm{error}_q=1\},\\
\widehat f_W&=\frac1{|W|}\sum_{q\in W}\ind\{\mathrm{fallback}_q=1\},\\
\widehat u_W&=\frac1{|W|}\sum_{q\in W}\ind\{\text{the prefix contains an unknown item}\},\\
\widehat m_{W,j}&=\frac1{|W|}\sum_{q\in W}\ind\{x_{qj}\text{ is missing}\}.
\end{align}
Response length, the fraction of empty responses,
latency percentiles, the candidate-count distribution,
snapshot age \(t_q-r(q)\), and the incompatible-version rate are measured separately.
For the neural component, monitoring covers rates of disabled modalities,
vector norms, the distribution of intent weights, and ANN disagreement
with exact search. The semantic module additionally tracks
the number of used codes, frequencies of code collisions between items,
truncation of large groups, and the quality gap between soft training
and final discrete serving.
Lower latency achieved through more empty responses
is not considered an improvement without accounting for the other constraints.

\subsection{Distribution changes}

Bin boundaries for a numeric feature are determined on training
data; missingness forms a separate category. Let \(p_b\) and \(q_b\)
be smoothed proportions in the reference and current windows, \(b=1,\ldots,B_h\):
\begin{equation}\label{eq:ops-hist}
p_b=\frac{n_b+\eps_h}{n+B_h\eps_h},\qquad
q_b=\frac{m_b+\eps_h}{m+B_h\eps_h},\qquad h_b=\frac{p_b+q_b}{2}.
\end{equation}
The Jensen--Shannon divergence is defined as
\begin{equation}\label{eq:ops-js}
\operatorname{JS}(p,q)=\frac12\sum_b p_b\log\frac{p_b}{h_b}
+\frac12\sum_b q_b\log\frac{q_b}{h_b}.
\end{equation}
It is zero for identical distributions and does not exceed
\(\log2\). It indicates an input change; it does not prove
that recommendation quality has deteriorated. Thresholds are set using sequential
control windows, without adjustment after inspecting the problematic period.

\paragraph{Registry and decision rule.}\label{sec:ops-drift-rule}
The registry includes input features and missingness, unknown items,
source masks, source quotas and list sizes, the fraction of unique candidates,
scores and final-list length, the selected branch, and failure reasons.
Categories and numeric boundaries are frozen from early data;
smoothing in~\eqref{eq:ops-hist} uses $\eps_h=1/2$.
The reference window contains at least 10\,000 requests; current windows are
nonoverlapping UTC hours with at least 1000 requests and 200 sessions.
Lower volume returns \code{insufficient}; the window is not declared
stable. A separate counter of insufficient telemetry
raises a notification after three consecutive such windows.

Thresholds are fixed jointly for the entire registry. Before operation,
on at least 100 control hourly windows preceding
$C_{\mathrm{select}}$ and not overlapping the reference window, compute
$A_w=\max_j\operatorname{JS}_{w,j}/0.05$. Window $w$ uses
only a previously available reference. The threshold
$h_{\mathrm{drift}}=\max(1,Q_{0.99}(A_w))$ is fixed together
with the reference, feature set, and nearest-rank empirical-quantile rule.
Control windows from $C_{\mathrm{smoke}}$ are allowed; the registry
and calibration rule are specified before opening them. Insufficient control windows
mean unverified monitoring, not permission to tune thresholds
on $D$. On calibration windows, record separately the exceedance rate
$A_w>h_{\mathrm{drift}}$ and the alert frequency of the complete rule below;
these are in-sample quantities, not false-alarm estimates.
Such an estimate requires separate windows and a justified assumption
that no real shift occurred. Time of day, weekday, and window
dependence are considered in diagnostics; future alert frequency is not guaranteed.

An alert is raised when $A_w>h_{\mathrm{drift}}$ in two consecutive
sufficient windows; it is cleared after three windows with
$A_w<0.8h_{\mathrm{drift}}$. Repeat notifications for the same cause
are sent at most once per hour, but new failures are recorded immediately.
The maximum over the full registry accounts for multiplicity in empirical
calibration; individual JS values are not interpreted as $p$-values
and are not subjected to FDR correction as significance tests.
Changing the registry, boundaries, or reference requires recalibration on early windows.
The notification recipient and investigation owner are recorded in the run manifest.

A distribution-change alert triggers checks of completeness, versions, sources,
and output distributions, without automatic retraining.
Separate 5-minute operational windows with at least 1000 requests
are compared with the initial technical-failure limit of $0.1\%$
and fallback limit of $1\%$. Exceedance in two consecutive windows
activates a previously verified fallback/rollback and pauses
enrollment of new experimental units; already assigned units remain in ITT.
Initial guardrail requirements, regardless of window volume:
the same protection is triggered by 20 consecutive
technical failures of valid admitted requests or three consecutive failed
internal health checks. Checks run
every 5 s with separate diagnostic IDs and a 120 ms deadline; they do not
modify user state or enter business-metric denominators.
Request sequences use arrival order;
an unfinished request receives its outcome when its deadline expires.
Success resets the corresponding counter. Errors from untrusted or
invalid inputs do not trigger a global rollback of a healthy version.
An unauthentic package is blocked before activation; a single attempt to
load it does not stop the active service. A detected eligibility violation
in actual serving blocks it immediately.
During A/B testing, the rules of Section~\ref{sec:eval-online} apply:
a conclusion from an ordinary fixed-horizon analysis after a guardrail
stop is not considered confirmatory.

Coveo dates are shifted by an integer number of weeks. Weekday
order is preserved. No link to actual holidays is established
without the original calendar dates.

\subsection{Metrics with delayed outcomes}

Historical evaluation updates Hit@20 for recorded targets
under the common calendar admission rule. Unknown coverage remains unknown.
Production-service conversion is computed only with verified
complete outcome logging for mature cohorts defined by assignment time.
Counting early positive outcomes while negative outcomes remain incomplete is prohibited.

In operation, current conversion can be observed as
\(\widehat p_W=\sum_{s\in W}Y_s/|W|\).
Its change relative to the previous week mixes the model effect,
traffic composition, and external circumstances. The causal effect
is estimated against a concurrently running control group.
Segments for confirmatory analysis of actual exposures are defined before variant
assignment. For example, search use after assignment cannot
support a claim about a causal effect
``among those who used search.''

\subsection{Model updates}

Data changes trigger diagnostics: log completeness,
version compatibility, target availability, and changes in candidate sources.
A new model is trained only on history that passes calendar admission
and the complete-read check, undergoes the same
checks, and is published as a new version. Crossing
the \(\operatorname{JS}\) threshold alone does not authorize automatic model replacement.

During a confirmatory A/B test, behavioral lookup tables
and the model are frozen by default. Training a shared model on responses
from both variants creates cross-group interference; this requires
a separate experimental design. Operational fixes
that do not alter the prescribed policy are documented.

\section{Explaining results and limits of interpretation}\label{sec:explanation}

For each recommended item, candidate-source provenance,
source scores, and the snapshot version are retained in a diagnostic trace
under the access and retention rules of Section~\ref{sec:ops-security}.
Related events are replaced by minimal reason codes;
their raw sequence is not included in the public explanation.
A graph explanation identifies the relation, the prefix item with the largest
contribution, and support $n^r_{ij}$. A content explanation identifies the modality
used and measured similarity. A search explanation refers
to the available query and compatible vector space.

An explanation describes the computation, not the cause of a purchase. It is limited to
Coveo fields: an anonymized vector does not reveal an item's name, purpose,
properties, or compatibility with a particular phone model.

The first variant uses templates based on the computation trace.
A language model does not extend the available facts; ranking based on one
requires a catalog with text and separate validation of quality, cost,
and feature availability.

\section{Result completeness and acceptance criteria}\label{sec:acceptance}

Historical evaluation is accepted under the \code{observed_purchases_v1} protocol
in Section~\ref{sec:eval-observed}. Success does not establish purchase growth
and does not, by itself, authorize an A/B test.

\subsection{Result states}

The state is represented by the tuple
\begin{equation}\label{eq:acceptance-state}
S=(S_{\mathrm{data}},S_{\mathrm{math}},S_{\mathrm{offline}},
S_{\mathrm{serve}},S_{\mathrm{business}}).
\end{equation}
Each component takes a value in
\(\{\mathrm{not\_run},\mathrm{passed},\mathrm{failed},
\mathrm{insufficient}\}\). For business impact,
\(\mathrm{not\_observed}\) is also used if no randomized
test has been conducted. Failure and insufficient data are distinct states.

\begin{longtable}{L{0.16\textwidth}L{0.73\textwidth}}
\toprule Component & Required conditions for \code{passed} \\\midrule\endhead
$S_{\mathrm{data}}$ & The run manifest is complete for the current stage; hashes and schema
match; the canonicalization rule is fixed; session overlap across periods
is empty; selected sessions have been fully read from the data release, calendar
admission and recorded targets are computed correctly. Unknown telemetry
completeness is preserved, including for sessions with purchases; exclusions are listed. \\
$S_{\mathrm{math}}$ & All enabled-component tests required
at the current stage under~\ref{app:tests} pass within the specified tolerances, including derivatives,
mask completeness, time boundaries, and the hit-loss identity. \\
$S_{\mathrm{offline}}$ & The frozen protocol has been executed on $D$, with
$N=N_D^+>0$ and an independent complete report on $\operatorname{Hit@20}_D$,
including the sensitivity comparison
without metadata. Replay mode and trace completeness satisfy
Section~\ref{sec:ops-replay-trace}. Every reported estimate
meets the assumptions, information and nondegeneracy thresholds
in~\ref{sec:eval-information}, the minimum block count, and required half-width:
$h_{\mathrm{pair}}$ for a pair, $h_{\mathrm{single}}$ for one policy,
both expressed as proportions.
The sign of the difference is reported separately: \code{passed} means the
estimate is usable, not that the model is superior. \\
$S_{\mathrm{serve}}$ & Each of the three load runs satisfies
all thresholds in~\ref{sec:ops-load-profile}, including the total limit and peak
memory accounting under~\ref{sec:ops-memory-limit}; feature-parity,
failure, recovery, integrity, and input-limit checks pass.
Partial source failure is allowed only within the registered
fallback policy and its limits. \\
$S_{\mathrm{business}}$ & The complete authorized A/B analysis is performed without
assignment or logging violations, the lower bound exceeds the prespecified
business threshold, and guardrail constraints are satisfied. \\
\bottomrule
\end{longtable}
A detected contract violation or measured breach of a
technical threshold yields \code{failed}; a correct procedure
with insufficient sessions containing a recorded target, blocks, or precision yields \code{insufficient}.
For the business component, with a usable interval $[\ell_B,u_B]$ and all
other conditions satisfied, the status is \code{passed} when
$\ell_B>\delta_{\mathrm{accept}}$, \code{failed} when
$u_B\le\delta_{\mathrm{accept}}$, and \code{insufficient} when
$\ell_B\le\delta_{\mathrm{accept}}<u_B$;
the threshold is defined in~\ref{sec:eval-acceptance}. If both
a violation and insufficient data are known, \code{failed} takes precedence;
all reasons are retained.
\code{not_run} means the stage has not yet been executed.

Run stages are defined in~\ref{sec:eval-stages}. While $D$ remains closed,
$S_{\mathrm{offline}}=\mathrm{not\_run}$, regardless of debugging success.
For local \code{ideal_deterministic} operation, state
$S_{\mathrm{offline}}$ does not depend on $S_{\mathrm{serve}}$;
Until service acceptance checks are run, $S_{\mathrm{serve}}=\mathrm{not\_run}$.
The report includes the loss decomposition in~\ref{sec:eval-loss-diagnosis}
and source contributions in~\ref{sec:eval-source-contribution}; technical failures
form a separate breakdown, not an additional term in the hit-loss decomposition.

A pilot A/B test additionally requires verified
complete experiment outcome logging, denoted $O_{\mathrm{AB}}=1$:
\begin{equation}\label{eq:acceptance-pilot}
\mathsf{Ready}_{\mathrm{pilot}}=
\ind\{O_{\mathrm{AB}}=1\}\,
\ind\{S_{\mathrm{data}}=S_{\mathrm{math}}=
S_{\mathrm{offline}}=S_{\mathrm{serve}}=\mathrm{passed}\}.
\end{equation}
For this formula, $S_{\mathrm{serve}}$ must be obtained specifically
under the operational profile: tenant and item limits are filled in,
the package signature is verified, retention and access rules are active,
and monitoring is calibrated on enough windows. Success under the local
\code{local_checksum_only} profile or \code{insufficient} monitoring
does not authorize a pilot. The acceptance profile is stated explicitly in the report.
Demonstrated business superiority additionally requires a positive
decision on \(\Delta B\) under Section~\ref{sec:eval-acceptance}
and all guardrail constraints. Power is a property of the prespecified
plan, not of the result computed after the experiment.
A complete specification does not establish that the system has passed testing.

\subsection{Single-request result}

\begin{equation}\label{eq:acceptance-request-result}
\mathsf R_q=(q,s,t_q,\mathsf V,L_q,\mathcal C_q,
\mathrm{branch},\mathrm{delivery\_status},
\mathrm{scores},\mathrm{sources},\mathrm{status},\mathrm{latency}).
\end{equation}
This is a logical diagnostic structure, not a requirement to retain
all candidates and scores for every production-service request indefinitely.
The minimum required record and bounded diagnostic log
are separated under~\ref{sec:ops-security}. Recorded targets are joined
only during evaluation; production-experiment outcomes are joined after maturation.
Candidates are replayed using the pinned mode
and trace from~\ref{sec:ops-replay-trace}. Without a complete trace, actual
degradation and $\mathcal C_q^{\mathrm{served}}$ are not considered measured;
a healthy recomputation does not automatically reconstruct them.
The historical report retains the session set, recorded targets,
the unknown status of telemetry completeness, selection parameters,
intervals, failures, and cost. The run identifier combines data,
configuration, code, and dependency versions. Changing the initial random-generator
state does not create an independent visitor sample.

\subsection{Complete computational formulation}

For configuration $\xi$, historical artifacts
$\mathsf A_G$, retrieval parameters $\theta_\xi$, and ranker $f_\xi$ are constructed sequentially.
Denote their fitting procedures by
$\mathcal T_G,\mathcal T_N,\mathcal T_R$:
\begin{align}
\mathsf A_G&=\mathcal T_G(D_{A_G};\xi),\\
\theta_\xi&=
\begin{cases}
\varnothing,&\text{neural extension disabled},\\
\mathcal T_N(D_{A_N};\mathsf A_G,\xi),&\text{extension enabled},
\end{cases}\\
\mathcal C_q^\xi&=\mathcal G_\xi(\mathcal H_q,\mathsf A_G,\theta_\xi,\mathcal I_q),\\
x_{qi}^\xi&=\Phi_\xi(\mathcal H_q,i,\mathsf A_G,\theta_\xi),\\
f_\xi&=\mathcal T_R(\{(x_{qi}^\xi,y_{qi}):s\in\mathcal S_B^+,
q=q_s^{(1)},\ i\in\mathcal C_q^\xi\}),\\
\widehat\xi&\in\argmax_{\xi\in\Xi_{\mathrm{feasible}}}
\operatorname{Hit@20}_{C_{\mathrm{select}}}(\pi_\xi).
\end{align}
The history for $\mathcal T_G$, first requests from sessions with a recorded
target for $\mathcal T_R$, and their targets are defined in Section~\ref{sec:eval-observed}.
In the first run, $\Xi_{\mathrm{feasible}}=\{\xi_0\}$:
one prespecified configuration is executed, without tuning
on $C_{\mathrm{select}}$ results. Period $A_N$ remains reserved;
its data do not augment the baseline system's graph or lookup tables.
Configurations for subsequent comparisons and the selection rule are fixed
before opening $C_{\mathrm{select}}$.
$\Xi_{\mathrm{feasible}}$ includes only variants that pass the declared
constraints. When comparing the original model with one new variant,
an exact Hit@20 tie retains the original model. Time and memory
are reported separately; load runs are unnecessary for tie-breaking.
A Hit@20 tie tolerance is not tuned to the results:
integer hit totals are compared on common sessions.
If no variants are eligible, selection receives status \code{insufficient}
or \code{failed} according to the established reason, and no winner is assigned.
The final list is sorted by $f_{\widehat\xi}$ unless
a separate diversification procedure was selected in advance. Period $D$ is excluded from every
fitting or selection procedure.
For the SID variant, $\theta_\xi$ includes an additional encoder,
decoder, and code-to-item mapping trained before $B$.

The business objective~\eqref{eq:business-optimization} is tested
with actual exposures under Appendix~\ref{sec:eval-online}.

\clearpage\part{Neural extension}\label{part:neural}
\section{Rationale and sequence for the neural extension}\label{sec:research}

This part follows the baseline run on p.~\pageref{sec:first-run}.
It adds one learned retrieval component; multiple intents and item
vectors are components of it, not separate mandatory stages.

\subsection{Relationship to work from 2024--2026}

The table lists the adopted principles. The formulas below define
an original adaptation for Coveo. Source status was checked
as of September 19, 2026.

\begin{longtable}{L{0.26\textwidth}L{0.64\textwidth}}
\toprule Work & Principle and role in the project \\\midrule\endhead
SPARE, ACM TORS, 2025 \cite{spare2025} &
Shortest paths and compact graph propagation.
Relation types, their costs, and the purchase target are adapted. \\
MiaSRec, SIGIR 2024 \cite{miasrec2024} &
Multiple session intents. Search-dependent intent weights
are added in this model. \\
MACL, Pattern Recognition, 2026 \cite{macl2025} &
Contrastive modality alignment; Coveo provides precomputed
vectors. \\
DIGER, SIGIR 2026 \cite{diger2026}; GRID, 2025 \cite{grid2025} &
Learned semantic codes and comparison of coding methods.
Used only in a separate study. \\
\bottomrule
\end{longtable}

The results of these studies concern their own data, products, and objectives;
they do not specify an expected purchase uplift on Coveo.

\subsection{Configuration and training chronology}

The neural policy $\pi_{\mathrm{main}}$ extends
$\pi_{\mathrm{base}}$ with learned graph and content vectors,
session intents, search weights, and a new candidate source.
The research policy $\pi_{\mathrm{SID}}$ additionally uses
semantic codes. Both preserve the overall limit $M$.

Training order and calendar admission of targets are specified in~\ref{sec:modern-time}.
The graph is published at $r_G$ before $A_N$, and learned retrieval at $r_N$ before $B$.
Fitting on $A_N$ is not independent evaluation on that period.
Before ranker training on $B$, the encoder and item vectors
are frozen; tree gradients do not flow into them.

\subsection{Implementation order and controlled extensions}
\label{sec:research-stages}

The first training run uses the configuration in~\ref{app:neural-config}.
One and two intents are compared according to the plan in~\ref{app:ablation}.
If the budget for the exact full-catalog loss is exceeded,
the sampled mode in~\ref{sec:sampled-overview} is allowed; its quality is checked
separately. It receives no additional candidate quota.

To control for history processing, Appendix~\ref{sec:sequence-overview}
retains the item graph and input information while replacing the session encoder.
A variant with additional cart actions or a sequence of
search queries is tested separately: it also changes the input set.
The value of the graph itself is tested by disabling the graph branch.

Diagnostics in~\ref{sec:eval-next-step} motivate the next experiment.
All comparisons specify the control, changed component,
number of runs, final hit rate, uncertainty of the difference, and cost.
A smaller gap to the candidate bound is not, by itself,
a criterion for improvement.

\subsection{Hypotheses following diagnostics}\label{sec:research-hypotheses}

The following variants are outside the required recipe; their
benefit is not assumed.
\begin{description}
\item[Search history.]
When losses occur in sessions with multiple queries, examine the number of
available queries, changes in their vectors, and consistency with views.
All features use available history; a change in interest
is not treated as an error in the search signal. Original text length
cannot be reconstructed from the supplied vector.
\item[Hard negatives.]
Investigate changes to the sampling distribution under the conditions in
Appendix~\ref{sec:sampled-training}; compare final policies,
not numerical values of different approximate losses.
\item[Category smoothing.]
When losses occur on rare items, test a mixture of the graph row
and a normalized category distribution. For weighted $W_{ij}$,
a Bayesian interpretation requires a separate observation model.
A full prior must not be expanded into a dense catalog matrix:
a bounded support or a separate candidate source is needed.
\end{description}

A neural ranker is considered after diagnosing losses
with retrieval fixed (Section~\ref{sec:model-ranking}).
Predicting purchase probability and calibrating it require a different observation
set (Section~\ref{sec:model-calibration}).
Catalog expansion requires a separate new-item admission rule
from Appendix~\ref{sec:freshness}; metadata alone do not establish an availability date.

\section{Learned model of graph, content, and session intents}
\label{sec:modern-model}

Data and training rules are specified in~\ref{sec:modern-time}, quotas
and settings in~\ref{app:neural-config}.
The sampled denominator in~\ref{sec:sampled-overview}
does not change the representation architecture.

\subsection{What is learned and when}
\label{sec:modern-time}

Catalog $\mathcal I_G$, graphs $P^r$, paths, and statistics are built from $A_G$
and fixed. Parameters $\theta_N$ are trained on $A_N$ using first requests
under calendar admission~\eqref{eq:eval-admission}. Ranker training,
configuration selection, and final evaluation follow the common order:
\begin{equation}\label{eq:modern-periods}
A_G\prec A_N\prec B\prec C\prec E\prec D.
\end{equation}
Symbol $\prec$ denotes strict temporal precedence with a calendar
gap of $H+g$. Period $E$ is reserved and unused.
No session belongs to two periods; telemetry completeness remains unknown.

Neighbors are determined only from $A_G$: a graph built from all of $A_N$ would reveal
future links to its early prefixes. The epoch count, dimension, and loss
coefficients for the first training run are fixed. Subsequent comparisons on $C_{\mathrm{select}}$
follow a prespecified plan; every variant is trained only on $A_N$
without further training on $C$ or $E$.

\subsection{Shortest paths in the fixed graph}
\label{sec:modern-paths}

To search for indirect links, combine known relations:
\begin{equation}\label{eq:modern-path-base}
\bar P_{ij}=\sum_{r\in\mathcal R_G}\zeta_rP^r_{ij},\qquad
\zeta_r\ge0,\quad\sum_r\zeta_r=1.
\end{equation}
Coefficients $\zeta_r$ are part of the configuration; row sums may be
less than one because of empty relations. For $\bar P_{ij}>0$,
define the dimensionless edge length
\begin{equation}\label{eq:modern-path-edge}
\ell_{ij}^{G}=c_G-\log\bar P_{ij},\qquad c_G>0.
\end{equation}
A strong link reduces length, while $c_G$ penalizes each hop.
This geometry does not define purchase-sequence probabilities.

For a directed path $p=(i=i_0,i_1,\ldots,i_h=j)$, its length
is the sum of edge lengths. Bound the number of hops by $H_P$:
\begin{equation}\label{eq:modern-distance}
D_{ij}^{G}=
\min_{\substack{p:i\leadsto j\\1\le |p|\le H_P}}
\sum_{(u,v)\in p}\ell_{uv}^{G}.
\end{equation}
If no path exists, the minimum is $+\infty$; $i=j$ is excluded.
All lengths are positive, so a shortest path has no cycle.
Distance is directed: generally $D_{ij}^{G}\ne D_{ji}^{G}$.
For each vertex, retain at most $d_P$ items
with the smallest $D_{ij}^{G}\le R_P$, $R_P>0$.
The resulting set $\mathcal N_i^{P}$ and distances are frozen.
Equal distances are resolved by identifier order.

Bounded search tracks the state ``vertex, number of hops.''
Ordinary search with a single label per vertex may miss
a valid path under the $H_P$ constraint.
The approximate variant expands at most $B_P$ states
per source vertex, in increasing length order with fixed tie-breaking.
When the budget is exhausted, the valid neighbors found are retained
with a truncation flag. Their distances $\widetilde D_{ij}^{G}$ are lengths
of discovered paths, not necessarily exact minima.
Below, $D_{ij}^{G}$ denotes the configuration-specified variant:
exact or approximate. It is shared by the model version and computed in advance.

Illustrative example: $\bar P_{ab}=0.8$, $\bar P_{bc}=0.7$,
$\bar P_{ac}=0.1$, $c_G=0.1$. Edge $a\to c$ has length approximately
$2.403$, while path $a\to b\to c$ has length $0.780$; when $H_P\ge2$,
the link to $c$ may be retained through $b$.

\subsection{Learned links and separate item representations}
\label{sec:modern-item}

For $i\in\mathcal I_G$, the learned identifier vector $b_i^N\in\R^d$
is initialized as specified in Section~\ref{sec:modern-training-recipe}. Introduce normalization
\begin{equation}\label{eq:modern-normalize}
\mathsf n_\eps(v)=\frac{v}{\max(\eps_N,\norm{v})},\qquad\eps_N>0.
\end{equation}
It is defined at zero as well; the result's norm does not exceed one.
The usable-direction mask is $\ind\{\norm{v}>\eps_N\}$.

The historical distribution over retained paths and its learned
adjustment are defined as follows:
\begin{align}\label{eq:modern-edge-attention}
\nu_{ij}^{P}
&=\frac{\exp(-D_{ij}^{G}/\tau_P)}
{\sum_{v\in\mathcal N_i^P}\exp(-D_{iv}^{G}/\tau_P)},\\
A_{ij}^{N}
&=\frac{\nu_{ij}^{P}\exp(\eta_N\langle
\mathsf n_\eps(b_i^N),U_N\mathsf n_\eps(b_j^N)\rangle)}
{\sum_{v\in\mathcal N_i^P}\nu_{iv}^{P}\exp(\eta_N\langle
\mathsf n_\eps(b_i^N),U_N\mathsf n_\eps(b_v^N)\rangle)}.
\nonumber
\end{align}
Here $\tau_P>0$ and $\eta_N\ge0$ are configuration parameters,
and $U_N\in\R^{d\times d}$ is learned on $A_N$.
The formulas apply only to nonempty $\mathcal N_i^P$;
otherwise, all $A_{ij}^{N}$ are zero.
Weights are learned over fixed links.

\paragraph{Log probabilities.}\label{sec:modern-log-probabilities}
On a nonempty admissible support, first compute
$\widetilde z_j=z_j-\max_v z_v$, then
$\operatorname{logsoftmax}(z)_j=\widetilde z_j-\log\sum_v e^{\widetilde z_v}$.
All differences $z_j-\operatorname{LSE}(z)$ below are computed this way,
not through the logarithm of rounded softmax values.
For~\eqref{eq:modern-edge-attention}, retain
\begin{align}\label{eq:modern-edge-log-probabilities}
\ell_{ij}^{P}&=-D_{ij}^{G}/\tau_P
-\operatorname{LSE}_{v\in\mathcal N_i^P}(-D_{iv}^{G}/\tau_P),\\
z_{ij}^{A}&=\ell_{ij}^{P}+\eta_N\langle
\mathsf n_\eps(b_i^N),U_N\mathsf n_\eps(b_j^N)\rangle,\nonumber\\
\ell_{ij}^{A}&=z_{ij}^{A}
-\operatorname{LSE}_{v\in\mathcal N_i^P}(z_{iv}^{A}),\qquad
A_{ij}^N=\exp(\ell_{ij}^{A}).\nonumber
\end{align}
Rounding an exponential to zero is allowed; its logarithm is not reconstructed
from that zero. Entropy and KL below use the stored log probabilities
and are differentiated through the same computation. Masked positions are excluded
before operations; empty support produces a separate zero branch with zero
gradient. Multiplying zero by infinity and replacing NaN with zero are prohibited.

The four branch representations are
\begin{align}\label{eq:modern-item-branches}
e_i^{\mathrm{id}}&=\mathsf n_\eps(b_i^N),&
e_i^{G}&=\mathsf n_\eps\!\left(
\sum_{j\in\mathcal N_i^P}A_{ij}^{N}B_Nb_j^N\right),\\
e_i^{T}&=\mathsf n_\eps(T_N\bar v_i^{\mathrm{text}}),&
e_i^{V}&=\mathsf n_\eps(V_N\bar v_i^{\mathrm{image}}).
\nonumber
\end{align}
Matrix dimensions are $B_N\in\R^{d\times d}$,
$T_N\in\R^{d\times d_t}$, and $V_N\in\R^{d\times d_v}$.
The identifier mask is $\ind\{i\in\mathcal I_G\}$;
the graph mask is $\ind\{\mathcal N_i^P\ne\varnothing\}$;
text and image use $M_i^{\mathrm{text}}$
and $M_i^{\mathrm{image}}$. Denote these four masks by $m_i^b$,
$b\in\{\mathrm{id},G,T,V\}$.
An unavailable branch produces a zero vector and a zero mask.

Let $\chi_i$ contain modality masks and
$\log(1+V_i^{\mathrm{detail}})$ from $A_G$.
When at least one branch is available, define
\begin{equation}\label{eq:modern-item-fusion}
\gamma_{ib}^{N}=
\frac{m_i^b\exp(c_b+v_b^\top\chi_i)}
{\sum_a m_i^a\exp(c_a+v_a^\top\chi_i)},\qquad
e_i=\mathsf n_\eps\!\left(\sum_b\gamma_{ib}^{N}e_i^b\right).
\end{equation}
Coefficients $c_b,v_b$ are learned. If all masks are zero,
$e_i=0$; define $m_i^N=\ind\{\norm{e_i}>\eps_N\}$.
Zero directions are excluded only from neural retrieval.

Coveo vectors are fixed; the projections and listed parameters are learned.
An unknown identifier in the prefix may use the supplied
vectors under the static-metadata assumption. Its identifier
and graph branches are absent, and historical frequency is zero;
the candidate catalog is not expanded.

The learned search representation uses the same text projection:
\begin{equation}\label{eq:modern-query-vector}
z_q^N=\mathsf n_\eps(T_N\bar z_q),\qquad
m_q^{SN}=M_q^S\ind\{\norm{T_N\bar z_q}>\eps_N\}.
\end{equation}
Query and image are compared only in the learned shared space,
not by taking the dot product of their original vectors.

\subsection{Multiple intents and selection by search context}
\label{sec:modern-intents}

From $\mathcal D_q$, retain events with at least one available input
branch, that is, $\sum_bm_i^b>0$, obtaining $\mathcal D_q^N$.
This condition depends on the data, not the learned norm of $e_i$.
Let $n_q^N$ be the number of distinct
items in this set, $J_I\ge1$ the upper bound on the number of
intents, and $J_q=\min(J_I,n_q^N)$.
Each $k=1,\ldots,J_I$ has a learned vector
$c_k^N\in\R^d$. For event $\ell$, define
\begin{align}\label{eq:modern-history-transform}
r_{q\ell}^{N}&=\tanh\!\left(
F_Ne_{i_{s(q)\ell}}
+a_N^t\log\!\left(1+\frac{t_q-t_{s(q)\ell}}{t_{\mathrm{unit}}}\right)
+a_N^f\log(1+n_{q,i_{s(q)\ell}}^{\mathrm{detail}})
\right),\\
z_{qk\ell}^{H,N}
&=-\frac{t_q-t_{s(q)\ell}}{\tau_H}+(c_k^N)^\top r_{q\ell}^{N},
\qquad k\le J_q,\nonumber\\
\beta_{qk\ell}^{N}
&=\exp\!\left(z_{qk\ell}^{H,N}
-\operatorname{LSE}_{v\in\mathcal D_q^N}(z_{qkv}^{H,N})\right),
\nonumber\\
h_{qk}&=\mathsf n_\eps\!\left(
\sum_{\ell\in\mathcal D_q^N}\beta_{qk\ell}^{N}e_{i_{s(q)\ell}}
\right).
\nonumber
\end{align}
Matrix $F_N$ has dimension $d\times d$, and $a_N^t,a_N^f$ are
vectors of dimension $d$. Function $\tanh$ is applied coordinatewise.
This normalizes the temporal weights in~\eqref{eq:model-history-weight}
with an intent adjustment, computed without intermediate rounding of
$\alpha_{q\ell}$: mask $\mathcal D_q^N$ is applied before stable LSE.
Frequencies and ages come from the prefix; equal timestamps have
equal ages, and permutations within a group do not change the sums.

Dynamic intent selection through entmax from the full MiaSRec
is not reproduced. Upper bound $J_I$ limits cost;
for a single item, $J_q=1$. Directions are latent vectors,
not established categories of customer goals.

Let $\mathcal K_q=\{1,\ldots,J_q\}$ be the direction indices.
If this set is empty but $M_q^S=1$, define
$\mathcal K_q=\{1\}$ and $h_{q1}=z_q^N$.
If the input search vector is also absent, the set remains empty.
A zero direction retains the training request with zero scores:
shrinking norms must not exclude examples from the loss.
For nonempty $\mathcal K_q$, intent weights are
\begin{align}\label{eq:modern-intent-weights}
\delta_{qk}^{N}&=a_k^I+(v_k^I)^\top h_{qk}
+\beta_Sm_q^{SN}e^{-d_q^S/\tau_{SN}}
\langle z_q^N,h_{qk}\rangle,\\
\omega_{qk}&=\frac{\exp(\delta_{qk}^{N})}
{\sum_{v\in\mathcal K_q}\exp(\delta_{qv}^{N})},\qquad
u_q^I=\sum_{k\in\mathcal K_q}\omega_{qk}h_{qk}.
\nonumber
\end{align}
Parameters $a_k^I,v_k^I$ are learned;
for entropy, retain $\ell_{qk}^{I}=\delta_{qk}^N-
\operatorname{LSE}_{v\in\mathcal K_q}(\delta_{qv}^N)$ and
$\omega_{qk}=\exp(\ell_{qk}^{I})$.
$\beta_S=\log(1+e^{b_S})\ge0$ is parameterized by learned $b_S$,
and $\tau_{SN}>0$ is measured in seconds.
When $m_q^{SN}=0$, the entire search term is zero;
$d_q^S$ is fully defined in~\eqref{eq:model-search-features}.
Weight $\omega_{qk}$ specifies a direction's relative importance,
not a purchase probability. With only one direction, its weight is always
one, so add a direct search contribution with a bounded coefficient:
\begin{align}\label{eq:modern-search-gate}
\ell_q^g&=\begin{cases}
b_g+(v_g)^\top u_q^I+(w_g)^\top z_q^N+a_g^t a_q^S,
&m_q^{SN}=1,\\
0,&m_q^{SN}=0,
\end{cases}\\
g_q^{SN}&=m_q^{SN}\bar g_N\sigmoid(\ell_q^g),
\qquad 0\le\bar g_N\le1,\nonumber\\
u_q^N&=(1-g_q^{SN})u_q^I+g_q^{SN}z_q^N.
\nonumber
\end{align}
Here $\sigmoid(x)=(1+e^{-x})^{-1}$;
$b_g,a_g^t\in\R$ and $v_g,w_g\in\R^d$ belong to $\theta_N$.
In the absence of a valid search direction, $m_q^{SN}=0$,
giving $g_q^{SN}=0$, including when $M_q^S=1$ but the projection is zero.
Degenerate search does not reduce the scale of a usable history vector.
Limit $\bar g_N$ is specified by the configuration.
For empty $\mathcal K_q$, define $u_q^I=0$.
Training-example admission uses the original mask $M_q^S$, not
learned $m_q^{SN}$; reducing the projection norm does not exclude the example
from the loss. The direction threshold creates a piecewise differentiable computation:
the mask is fixed during gradient computation, and the derivative
is not claimed to be continuous at the threshold.
With search alone, $h_{q1}=z_q^N$ and the formula preserves
$u_q^N=z_q^N$. Without search, $u_q^N=u_q^I$.
The convex combination preserves $\norm{u_q^N}\le1$ and can change
the direction even when $\omega_{q1}=1$. When $\bar g_N=0$, search affects
intent weights and the baseline source; for search-only history,
the fallback rule $h_{q1}=z_q^N$ still applies.

Illustrative example: for orthogonal directions with equal initial
weights, a search vector aligned with the first, with $\beta_S=\log 3$
and zero age, yields weights $3/4$ and $1/4$.

\subsection{Candidate retrieval and ranking features}
\label{sec:modern-retrieval}

For $\tau_N>0$, the neural score is
\begin{align}\label{eq:modern-score}
s_{qi}^{N}
&=\frac{1-g_q^{SN}}{\tau_N}
\sum_{k\in\mathcal K_q}\omega_{qk}\langle h_{qk},e_i\rangle
+\frac{g_q^{SN}}{\tau_N}\langle z_q^N,e_i\rangle\\
&=\frac{\langle u_q^N,e_i\rangle}{\tau_N}.
\nonumber
\end{align}
Temperature $\tau_N$ changes the loss scale while preserving item order.
The norm constraints imply $|s_{qi}^{N}|\le1/\tau_N$.

The primary new source returns the first $m_N^{\mathrm{mix}}$ items
by $s_{qi}^N$ among $i\in\mathcal I_q$ with $m_i^N=1$.
The item-vector matrix is fixed after training.
For practically zero $u_q^N$, the source returns an empty list
with a mask, avoiding arbitrary selection among equal zero scores.

An additional source selects up to $m_{Nk}$ items
by $\langle h_{qk},e_i\rangle$ for each active intent.
This preserves items associated with a direction of small $\omega_{qk}$.
Lists are treated as separate sources but use one matrix
$e_i$: the mixture $u_q^N$ and all $h_{qk}$ need only different queries
to the same shared index. A separate index per intent is unnecessary.
The total neural quota $m_N$ and baseline quotas satisfy
\begin{equation}\label{eq:modern-quota}
m_N^{\mathrm{mix}}+\sum_{k=1}^{J_I}m_{Nk}\le m_N,\qquad
m_N+\sum_{a\in\mathcal A_{\mathrm{base}}}m_a\le M.
\end{equation}
Inactive directions and directions with norm no greater than $\eps_N$
have empty lists during candidate retrieval.
The overall budget is not increased; deduplication, popularity
backfill, and eligibility follow Section~\ref{sec:model-candidates}.

For every candidate in the union, the ranker additionally receives
$s_{qi}^{N}$, $\max_{k\in\mathcal K_q}\langle h_{qk},e_i\rangle$,
similarity $\langle z_q^N,e_i\rangle$ with mask $m_q^{SN}m_i^N$,
coefficient $g_q^{SN}$,
the number of directions with norm above $\eps_N$, and the following two statistics:
\begin{align}\label{eq:modern-ranking-features}
H_q^I&=-\sum_{k\in\mathcal K_q}\exp(\ell_{qk}^{I})\ell_{qk}^{I},\\
V_{qi}^{I}&=\sum_{k\in\mathcal K_q}\omega_{qk}
\left(\frac{\langle h_{qk},e_i\rangle}{\tau_N}
-\frac{\langle u_q^I,e_i\rangle}{\tau_N}\right)^2.
\nonumber
\end{align}
Entropy describes how distributed the weights are; variance describes differences
in item scores. Its center $u_q^I$ excludes the direct search contribution,
which is not an intent in the mixture.
Without usable vectors, values are zero with masks.
Provenance and rank features for the new source are retained separately.
The final $f(x_{qi})$ is retrained on $B$ for every architecture
being compared. The neural score is not called a conversion probability.

\subsection{Primary training objective: recorded future purchases}
\label{sec:modern-objective}

Let $\mathcal Q_N$ contain the first requests from selected sessions
$\mathcal S_{A_N}^{\mathrm{obs}}$. The primary loss uses only sessions
$\mathcal S_{A_N}^{+}$ with a recorded future purchase;
the others are available to the auxiliary tasks below. For a request, define
$Y_q=Y_{s(q)}^{\mathrm{obs}}$, $Y_q^G=Y_q\cap\mathcal I_q$, and
$a_q^{\mathrm{buy}}=\ind\{|Y_q^G|>0\}$.
When $a_q^{\mathrm{buy}}=1$, the target distribution is
$t_{qi}^{\mathrm{buy}}=\ind\{i\in Y_q^G\}/|Y_q^G|$.
Distinct items in recorded purchases receive equal shares regardless
of unknown quantity and price.
The fixed training-input mask is
$m_q^N=\ind\{|\mathcal D_q^N|>0\ \text{or}\ M_q^S=1\}
\ind\{|\mathcal I_q|>0\}$.
Set $\mathcal D_q^N$ and the mask depend on input data,
not on learned norms.
With empty $\mathcal K_q$ and no search, $u_q^N=0$.
For any nonempty $\mathcal I_q$, the normalized model is defined as
\begin{equation}\label{eq:modern-softmax}
p_N(i\given q)=\frac{\exp(s_{qi}^{N})}
{\sum_{j\in\mathcal I_q}\exp(s_{qj}^{N})}.
\end{equation}
The exact denominator includes the entire nonempty request catalog,
including zero item representations. When $u_q^N=0$, the distribution
is uniform and the example is retained. Requests without available input
are excluded before training; their fraction is recorded.

Each session contributes one request, so $w_q^N=1$;
additionally, set $a_q^{\mathrm{buy}}=0$ when $m_q^N=0$.
The primary loss is
\begin{equation}\label{eq:modern-buy-loss}
\mathcal L_{\mathrm{buy}}^N=
-\frac{\displaystyle\sum_{q\in\mathcal Q_N}w_q^Na_q^{\mathrm{buy}}
\sum_{i\in\mathcal I_q}t_{qi}^{\mathrm{buy}}\log p_N(i\given q)}
{\displaystyle\sum_{q\in\mathcal Q_N}w_q^Na_q^{\mathrm{buy}}}.
\end{equation}
A request term with $a_q^{\mathrm{buy}}=0$ is zero without
evaluating undefined $t_{qi}^{\mathrm{buy}}$.
If the denominator is zero, purchase training is impossible;
the variant receives an insufficient-data status.
The loss learns to select an item from a recorded purchase, not to estimate
purchase probability or causal effect. Sessions without a recorded
target are available to auxiliary tasks but do not enter Hit@20 or become
confirmed negative outcomes.
Purchases outside $\mathcal I_G$ do not expand the catalog; their fraction is recorded,
and they remain possible misses during evaluation.

For one request, the score derivative has a simple form:
\begin{equation}\label{eq:modern-softmax-gradient}
\frac{\partial\left[-\sum_i t_{qi}\log p_N(i\given q)\right]}
{\partial s_{qj}^{N}}=p_N(j\given q)-t_{qj}.
\end{equation}
Parameter derivatives are computed by the chain rule.
For stability, the maximum score is subtracted; the denominator may
be computed in blocks. Item sampling and approximate-loss bias
are specified in Appendix~\ref{sec:sampled-overview}.

\subsection{Auxiliary next-view prediction}
\label{sec:modern-next}

Let $t_q^P$ be the time of the first recorded purchase after $t_q$
in the same session ($+\infty$ if no record exists).
The target is the next recorded view, with no claim
that actions between records are complete. The admissible time set is
\begin{equation}\label{eq:modern-next-time}
\mathcal T_q^+=\{t_e:\ e\in s(q),\ \kappa_e=\mathrm{detail},\
t_q<t_e\le\tau_s,\ t_e<t_q^P\}.
\end{equation}
If $\mathcal T_q^+\ne\varnothing$, set $t_q^+=\min\mathcal T_q^+$;
items in the entire group at that time form $D_q^+$.
The last view and a view exactly at the horizon are allowed;
a view simultaneous with a purchase is not. For empty $\mathcal T_q^+$,
set $D_q^+=\varnothing$ without computing $t_q^+$.
After intersection with $\mathcal I_q$, introduce mask
$a_q^{\mathrm{next}}$ for a nonempty target with $m_q^N=1$, and the uniform distribution
$t_{qi}^{\mathrm{next}}$ over that target.
Loss $\mathcal L_{\mathrm{next}}^N$ is defined by
formula~\eqref{eq:modern-buy-loss}, replacing
$(a^{\mathrm{buy}},t^{\mathrm{buy}})$ with
$(a^{\mathrm{next}},t^{\mathrm{next}})$.
With an empty denominator, this auxiliary loss is zero
with a flag indicating no training examples.
Weight $\lambda_{\mathrm{next}}\ge0$ is initially fixed,
then selected on $C_{\mathrm{select}}$ by Hit@20, including the zero-weight variant.
Next-view accuracy does not replace this criterion.

\subsection{Contrastive alignment of two representations}
\label{sec:modern-contrastive}

Two representations of one available prefix are constructed
by independently disabling available branches
$\mathrm{id},G,T,V$ with specified probabilities.
Each item retains at least one originally available branch;
if all are disabled, the first of them is restored
in a fixed order. Disabling changes the masks in
\eqref{eq:modern-item-fusion}, after which the session computation is repeated.
Events, their times, and the search vector are preserved.
Original texts and images are not required.

Denote the results of $\mathsf n_\eps$ normalization by
$\tilde u_q^{(1)},\tilde u_q^{(2)}$.
Identical masks do not create a pair; a zero learned vector
is retained in the loss with zero scores.
Each batch contains at most one request per session.
Signature $\mathsf{sig}(q)$ contains only encoder inputs:
items, relative times, frequencies, and supplied vectors.
In request batch $\mathcal B_N$, define the comparison set
\begin{equation}\label{eq:modern-contrastive-negative}
\mathcal V_q=\{v\in\mathcal B_N:
s(v)\ne s(q),\quad \mathsf{sig}(v)\ne\mathsf{sig}(q),\quad
Y_v^G\cap Y_q^G=\varnothing\}.
\end{equation}
This set includes only requests with differing random-transformation
masks; duplicate signatures among comparisons
are retained once in a fixed order.
Labels from $A_N$ are used only in the loss, excluding matching purchases
from contrasted examples. Different purchases do not establish
different interests.

For $\mathcal V_q\ne\varnothing$ and $\tau_C>0$, define
\begin{equation}\label{eq:modern-contrastive-loss}
\ell_q^{1\to2}=-\log
\frac{\exp(\langle\tilde u_q^{(1)},\tilde u_q^{(2)}\rangle/\tau_C)}
{\exp(\langle\tilde u_q^{(1)},\tilde u_q^{(2)}\rangle/\tau_C)
+\sum_{v\in\mathcal V_q}
\exp(\langle\tilde u_q^{(1)},\tilde u_v^{(2)}\rangle/\tau_C)}.
\end{equation}
Loss $\mathcal L_C^N$ is the weighted mean of
$(\ell_q^{1\to2}+\ell_q^{2\to1})/2$ over eligible requests
with weights $w_q^N$ normalized by their sum.
Without eligible pairs, the loss is zero, with a note in the report.
With random batches, the expected loss is minimized;
batch composition and size are specified by the configuration.

\paragraph{Reproducible selection of contrastive comparisons.}
Rule \code{contrastive_sampler=label_filtered_v1} is part of the training
run manifest; it is not a learned parameter. Batch order,
random masks, and loss normalization are specified in
Section~\ref{sec:modern-training-recipe}. After forming a batch,
only the filters in~\eqref{eq:modern-contrastive-negative} are applied;
masks and counts of excluded comparisons by reason are retained.

In a preregistered comparison,
\code{contrastive_sampler=label_free_v1} removes only the condition
$Y_v^G\cap Y_q^G=\varnothing$, preserving all other rules and the budget.
This measures the effect of using targets to select comparisons.
A request without a recorded purchase can serve as a comparison even without this
filter; absence of a record does not confirm different interests. These rules
apply only to training period $A_N$; candidates and inputs
in $B,C,D$ do not use future purchases when constructing representations.

The representation is assumed to tolerate partial modality loss.
This may fail when a distinguishing feature is lost, so a variant
without the contrastive component is required. Sessions in the denominator are not
confirmed negative exposures; preserving meaning when vectors are disabled
and increasing conversion are not guaranteed.

\subsection{Regularization and final training objective}
\label{sec:modern-training}

To limit deviations of graph weights from their historical reference,
introduce
\begin{equation}\label{eq:modern-graph-regularizer}
\mathcal L_G^N=
\frac1{|\mathcal I_G^P|}
\sum_{i\in\mathcal I_G^P}\sum_{j\in\mathcal N_i^P}
\exp(\ell_{ij}^{A})(\ell_{ij}^{A}-\ell_{ij}^{P}),\qquad
\mathcal I_G^P=\{i:\mathcal N_i^P\ne\varnothing\}.
\end{equation}
For empty $\mathcal I_G^P$, set $\mathcal L_G^N=0$.
On shared positive support, the penalty is nonnegative and vanishes
when rows match. It limits changes to historical links
without establishing causality.

The final objective is
\begin{equation}\label{eq:modern-total-loss}
\min_{\theta_N}\quad
\mathcal L_{\mathrm{buy}}^N
+\lambda_{\mathrm{next}}\mathcal L_{\mathrm{next}}^N
+\lambda_C\mathcal L_C^N
+\lambda_G^N\mathcal L_G^N
+\lambda_\theta\norm{\theta_N}^2,
\qquad\lambda_{\mathrm{next}},\lambda_C,\lambda_G^N,\lambda_\theta\ge0.
\end{equation}
The primary loss has weight one. Other weights do not represent
the monetary value of events: they are initially specified in Appendix~\ref{app:parameters},
then selected on $C_{\mathrm{select}}$ under the comparison plan. Zero weight disables a task.
The squared norm of $\theta_N$ is the sum of squared scalar parameters;
dimensions are fixed before training.

\subsection{Reproducible training recipe for extensions}
\label{sec:modern-training-recipe}

Profile \code{neural_training_v2} applies to enabled neural,
sequential, and SID modules. Reference environment: NumPy 2.1.3, PyTorch 2.6.0,
CPU, FP32 for parameters, losses, and optimizer state; one computation
thread and one inter-op thread, with deterministic algorithms.
AMP, graph compilation, gradient accumulation, learning-rate scheduling, and early
stopping are disabled. Exact builds, hardware, and dependency hashes
are pinned in the version file under Section~\ref{app:run-manifest};
bitwise agreement across different environments is not assumed.

\paragraph{Initialization and randomness.}
Use \code{Generator(PCG64(seed))}. For module $m=0,1,2$
(main neural model, sequential model, SID), separate
streams have seed $17+100m+j$: $j=0$ for initialization,
$j=1$ for example order, $j=2$ for branch masks,
$j=3$ for SID noise, $j=4$ for negative items,
$j=5$ for search dropout, and $j=6$ for random control SID codes.
A disabled consumer makes no calls; states are not reset
between epochs. Call order is specified in the corresponding sections.
The parameter schema contains names, shapes, roles, and shared references:
tensors are traversed by UTF-8 name, and coordinates in row-major order;
the shared text/search matrix is initialized once.
Random values are generated in FP64 by a \code{uniform} call for the whole
tensor, then cast once to FP32.

A linear-transformation matrix of shape $d_{\mathrm{out}}\times
d_{\mathrm{in}}$ receives independent
$\mathrm{Uniform}[-a,a]$ values, $a=\sqrt{6/(d_{\mathrm{in}}+d_{\mathrm{out}})}$.
Identifier vectors $b_i^N$, intent vectors $c_k^N$, and codebook vectors
$c_{\ell b}^{\mathrm{SID}}$ receive
$\mathrm{Uniform}[-\sqrt{3/d_{\mathrm{emb}}},\sqrt{3/d_{\mathrm{emb}}}]$,
where $d_{\mathrm{emb}}$ is their dimension;
rows are ordered by item ID or numeric index.
Sequential-model positions $p_b$ receive
$\mathrm{Uniform}[-0.02,0.02]$. All other vectors and scalars are zero,
except LN scale factors $\gamma$, which are one,
and $b_g$ from~\eqref{eq:app-search-init}.
Matrix exception: $U_N=0$, so initially $A_{ij}^N=\nu_{ij}^P$.
The same rules apply to the sequential model's own parameters;
blocks absent from its architecture are not created.
These rules also specify $b_S=0$, decoder biases, and time/frequency
coefficients; intents do not use identical zero initialization.

\paragraph{Batches and loss scaling.}
Plans for all epochs are constructed before training. For each epoch,
sessions sorted by UTF-8 are shuffled using
\code{permutation}; their first requests are split into batches of size 64.
Every eligible request appears exactly once per epoch;
the short final batch is retained. The batch plan and batch count
$T_{\mathrm{ep}}$ are stored before training.
For contrastive masks, iterate over requests in batch order, views
$1,2$, unique input-prefix items by UTF-8 ID, and branches
$\mathrm{id},G,T,V$. Each branch receives one
\code{random()} call, including absent branches; an available branch is disabled
when the value is below the specified probability. Restoration
after complete dropout follows Section~\ref{sec:modern-contrastive}.

Let $\ell_q^r$ be the request term in the full loss for task
$r\in\{\mathrm{buy},\mathrm{next}\}$, already including $w_q^N$,
the mask, and the global denominator in~\eqref{eq:modern-buy-loss}.
On batch $\mathcal B$, minimize
\begin{equation}\label{eq:modern-batch-loss}
\mathcal L_{\mathcal B}
=T_{\mathrm{ep}}\sum_{q\in\mathcal B}
(\ell_q^{\mathrm{buy}}+\lambda_{\mathrm{next}}\ell_q^{\mathrm{next}})
+\lambda_C\mathcal L_{C,\mathcal B}^N
+\lambda_G^N\mathcal L_G^N+\lambda_\theta\norm{\theta_N}^2.
\end{equation}
Global denominators are not recomputed per batch.
At fixed parameters, the mean of the first terms over all batches
equals the full purchase and auxiliary loss, including short
batches. Contrastive loss is normalized within each batch according to its
definition; the graph penalty is computed over the entire fixed graph
at every step. This exact cost is included in the training budget.

\paragraph{Parameter updates.}
Use \code{torch.optim.Adam} with one group containing all unfrozen
parameters, zero initial moments, and settings
\begin{verbatim}
lr=0.001; betas=(0.9,0.999); eps=1e-8; weight_decay=0
amsgrad=false; foreach=false; fused=false; maximize=false
capturable=false; differentiable=false
\end{verbatim}
Each step computes the full gradient of $\mathcal L_{\mathcal B}$, including
the explicit L2 penalty; gradients are dense, with zero for absent contributions.
Its global Euclidean norm is computed in FP64. If the norm exceeds 5,
the entire gradient is multiplied by $5/\norm{g}$ before updating Adam moments;
for zero norm, no division is performed. L2 is not added again
through the optimizer. Nonfinite losses, gradients, or parameters abort
training. Exactly one Adam step is performed after each batch;
its moments and step counter continue across epochs.
The main and sequential models are trained for three epochs.
A checkpoint stores parameters, moments, step number, batch plan,
and all generator states.

\subsection{Checking intent distinctness}
\label{sec:modern-intent-diagnostics}

Several directions may describe one interest. Their coincidence
is not ruled out by the loss. Let
$\mathcal K_q^+=\{k\in\mathcal K_q:\norm{h_{qk}}>\eps_N\}$,
$n_q^I=|\mathcal K_q^+|$, and
$\widehat h_{qk}=h_{qk}/\norm{h_{qk}}$ for these directions.
For $n_q^I\ge2$ and prespecified $0<\delta_I<1$, define
\begin{align}\label{eq:modern-intent-duplicate}
\operatorname{Sim}_q^I
&=\frac{2}{n_q^I(n_q^I-1)}
\sum_{\substack{k<l\\k,l\in\mathcal K_q^+}}
\langle\widehat h_{qk},\widehat h_{ql}\rangle,\\
\operatorname{Dup}_q^I
&=\frac{2}{n_q^I(n_q^I-1)}
\sum_{\substack{k<l\\k,l\in\mathcal K_q^+}}
\ind\{\langle\widehat h_{qk},\widehat h_{ql}\rangle\ge1-\delta_I\}.
\nonumber
\end{align}
The first metric is mean similarity, the second the fraction of nearly coincident
pairs. They are undefined when $n_q^I<2$; report the fraction of such requests,
norms, and the frequency of zero directions so that zeroing cannot hide degeneration.

Let $\mathcal C_{qk}^I$ be the candidates of a single intent.
For pairs with a nonempty union, measure
\begin{equation}\label{eq:modern-intent-overlap}
O_{qkl}^I=
\frac{|\mathcal C_{qk}^I\cap\mathcal C_{ql}^I|}
{|\mathcal C_{qk}^I\cup\mathcal C_{ql}^I|}.
\end{equation}
For a request, average $O_{qkl}^I$ over eligible pairs $k<l$;
if no such pairs exist, the result is undefined.
Metrics are averaged over first requests on which they are defined.
Missing-value rates use all first requests; missing values are not replaced by zeros.
Separately compare $\operatorname{Coverage}_{\mathrm{cand}}$
and final Hit@20 for $J_I=1$ and $J_I=2$
with the same total neural-source quota.
Each variant's ranker is retrained.
A high $H_q^I$ is possible even for identical directions.

Breakdowns use history length and search presence. The penalty
$\sum_{k\ne l}\langle c_k^N,c_l^N\rangle^2$ is not part of the initial objective:
it is also minimized when all $c_k^N=0$, without ensuring distinct intents.
Added complexity is justified by measured direction overlap
without additional coverage, not by the number of viewed items alone.

\subsection{Search-signal robustness}
\label{sec:modern-search-robustness}

Limit $\bar g_N$ bounds the direct contribution of $z_q^N$.
It does not bound the entire search influence: parameter $\beta_S$
in intent weights remains learned and has no upper bound.
Distributions of $g_q^{SN}$, $\max_k\omega_{qk}$,
$H_q^I$, and purchase hit rate by query age are retained.
The maximum is taken over $\mathcal K_q$; for an empty set,
it is undefined and accompanied by a missingness mask.
To compare the query with history, use the vector without search
mixing
\begin{equation}\label{eq:modern-history-probe}
v_q^{\mathrm{hist}}=
\mathsf n_\eps\!\left(\sum_{\ell\in\mathcal D_q^N}
\alpha_{q\ell}e_{i_{s(q)\ell}}\right).
\end{equation}
The cosine between $z_q^N$ and $v_q^{\mathrm{hist}}$ is computed only
when both norms exceed $\eps_N$. Disagreement may indicate a change of interest. This diagnostic uses historical data;
search after A/B assignment does not define a subgroup
for a confirmatory causal comparison.

If diagnostics show sensitivity to search, a separate
variant with random dropout of this signal during training is allowed.
For every prefix and step, independently draw
$D_q^S\sim\operatorname{Bernoulli}(1-p_S)$, $0\le p_S<1$, and use
\begin{equation}\label{eq:modern-search-dropout}
\widetilde M_q^S=D_q^SM_q^S
\end{equation}
in all search branches of the new encoder: projection and its mask,
intent weights, direct contribution, and the search-only fallback
representation. A disabled vector is replaced with zero.
The same $D_q^S$ applies to both contrastive views
of this prefix; baseline sources and target labels are unchanged.
For $p_S>0$, stream $j=5$ is called once through \code{random()}
for every request in batch order, including requests without search;
$D_q^S=\ind\{U\ge p_S\}$. No calls occur when $p_S=0$ or during serving.

Admission $m_q^N$ and the training-weight denominator are computed before random
dropout. If no input remains afterward, $u_q^N=0$ and the example
is retained with a uniform probability distribution over items.
The expectation of the selected loss over this randomness is minimized.
The initial configuration has $p_S=0$; a nonzero value is selected
on $C_{\mathrm{select}}$ only in a separate comparison. During serving, the signal is not disabled. This robustness check
does not prescribe suppressing search when interest changes and does not guarantee
diversity or increased hit rate.

\subsection{Cost and failure behavior}
\label{sec:modern-cost}

Let $N_G=|\mathcal I_G|$, and
$d_0\le|\mathcal R_G|d_G$ be the maximum out-degree
of the combined baseline graph.
Stored paths require $O(N_Gd_P)$ memory.
Bounded search expands at most $N_GB_P$ states
and examines at most $N_GB_Pd_0$ outgoing edges;
queue cost depends on the implementation.
Before truncation, exact search can be substantially more expensive in the worst case.
The fraction of vertices with truncation is reported alongside candidate quality.

For one full recomputation of item representations,
projections require on the order of
$N_G[d(d_t+d_v)+d^2]$ operations,
and graph aggregation after pretransforming
vectors requires on the order of $N_Gd_Pd$.
Once the model is fixed, the result is stored in an $N_G\times d$ matrix
and an index; the graph neural network is not recomputed during serving.
Session computation requires on the order of
$L_Hd^2+J_IL_Hd+d d_t$ operations
if the current search query must be projected again.
An exact catalog scan requires $O(N_Gd)$;
approximate search is validated under the same exact-search agreement
protocol as the baseline content sources.

With a full denominator, neural-training cost for
$N_Q^N$ requests per epoch includes $O(N_Q^NN_Gd)$;
blockwise computation reduces memory, not the number of products.
Additional costs are graph-representation and projection recomputation at every
parameter update and contrastive pairs. Cached vectors
cannot replace representations from the current encoder.
Time, memory, and epoch limits are fixed before training;
exceeding them is a computational failure.

New-source failures are handled only under the common contracts in
\ref{sec:ops-partial-data} and~\ref{sec:ops-deadlines}. Unregistered
degradation activates $\mathcal C_q^{\mathrm{fallback}}$ with its fixed
order, without the primary model's ranker.

\begin{invariant}[Fixed knowledge and available prefix]
With $\theta_N$ and $A_G$ artifacts fixed, changing events
after $t_q$ does not change $h_{qk}$, $\omega_{qk}$, $g_q^{SN}$,
$u_q^I$, $u_q^N$, or $s_{qi}^N$.
Permuting rows within one time group does not change the result.
Disabling all new sources restores the baseline policy
when its original quotas and ranker are restored.
\end{invariant}

\clearpage\part{Appendices: computational variants and additional studies}\label{part:appendices}
\appendix
\section{Approximate training with item sampling}
\label{sec:sampled-training}
\label{sec:sampled-overview}

This is an optional computational mode of the neural extension;
the baseline system does not use it. The reference is the full
denominator in~\eqref{eq:modern-softmax}. All positive items
are included exactly; the rest use weighted random
sampling \cite{candidatesampling}. Bias requires validation
\cite{sampledsoftmax2025}; MIDX is not reproduced.
Agreement in gradients and quality with the exact mode is not assumed.
NCE, which solves a different task of distinguishing observations from noise, is not used.

\subsection{Positive items and sampling distribution}
\label{sec:sampled-proposal}

Consider one eligible request $q$ from $A_N$ and one task
$a\in\{\mathrm{buy},\mathrm{next}\}$ with a nonempty target.
Denote its item set by $T_q^a$:
$T_q^{\mathrm{buy}}=Y_q\cap\mathcal I_q$,
$T_q^{\mathrm{next}}=D_q^+\cap\mathcal I_q$.
Eligibility masks and request weights are preserved
from Sections~\ref{sec:modern-objective}--\ref{sec:modern-next}.
On $T_q^a$, the original uniform distribution $t_{qi}^a$ is defined.
Complement $\mathcal N_q^a=\mathcal I_q\setminus T_q^a$ contains
negative items for the training task; absence from the target does not establish
that an item was viewed and rejected.

Let $N_G=|\mathcal I_G|>0$, and let $\pi_j^{\mathrm{pop}}$ be
the popularity distribution from~\eqref{eq:model-popularity},
built only from $A_G$. For fixed $\delta\in(0,1]$,
define
\begin{equation}\label{eq:sampled-base-proposal}
Q_G(j)=(1-\delta)\pi_j^{\mathrm{pop}}+\frac{\delta}{N_G},
\qquad j\in\mathcal I_G.
\end{equation}
The uniform component gives every catalog item
positive probability. If $\mathcal N_q^a\ne\varnothing$,
the conditional distribution for this request is
\begin{equation}\label{eq:sampled-proposal}
Q_q^a(j)=
\frac{Q_G(j)}
{\sum_{v\in\mathcal N_q^a}Q_G(v)},
\qquad j\in\mathcal N_q^a.
\end{equation}
Allowed mass $b_q^a=\sum_{j\in\mathcal N_q^a}Q_G(j)$
is computed at the start of each enabled task with a nonempty complement
by direct compensated Kahan summation in FP64, in UTF-8 ID order;
subtracting excluded-item mass from one is not used.
In the mathematical model, sampling from $Q_G$ and rejecting forbidden
items gives $Q_q^a$ and requires on average $m/b_q^a$ attempts for $m$
items; small mass reduces the benefit and is recorded.
Distribution $Q_G$ is independent of learned weights and fixed from $A_G$.
Labels exclude targets only during training on $A_N$.

In ideal sampling, independently draw $m\ge1$ items
$J_1,\ldots,J_m$ with replacement from $Q_q^a$.
Denote the number of times item $j$ is drawn by
$C_{qj}^a=\sum_{b=1}^{m}\ind\{J_b=j\}$.
Duplicates retain their multiplicity; otherwise, the sum's expectation changes.
The sample is refreshed each time a request is presented;
stream $j=4$ from profile~\ref{sec:modern-training-recipe} is used.
In batch order, process the purchase task and then the next-view
task for each request. Cumulative sums of $Q_G$ in UTF-8 ID order
are built once in FP64, divided by the final sum, and stored as $F$
with final element 1.

Before sampled-mode training, verify that the computed
$Q_G(j)$ are finite and positive and that $0=F_0<F_1<\cdots<F_{N_G}=1$.
With items indexed in this order, the grid
$U\in\{0,1,\ldots,2^{53}-1\}/2^{53}$ from an FP64 \code{random()}
call corresponds to masses
\[
\rho_j=2^{-53}\left(\left\lceil2^{53}F_j\right\rceil
-\left\lceil2^{53}F_{j-1}\right\rceil\right).
\]
Every item must satisfy $\rho_j>0$ and
$|\rho_j-Q_G(j)|\le10^{-6}Q_G(j)$; this is a project tolerance,
and the comparison uses $Q_G(j)$ computed in FP64.
Violation of these conditions, or nonfinite or nonpositive $b_q^a$
for a nonempty complement, causes failure
\code{negative_sampling_precision_unsupported} before consuming the sampling
stream or updating parameters. No item is removed from support,
and probabilities are not silently set to zero. This check does not establish exact sampling
from real-valued $Q_G$: the expectation identities and limiting statements
below concern ideal sampling; the implementation is its numerical
approximation with the stated tolerance for base probabilities.

Each attempt makes one \code{random()} call and
\code{searchsorted(F, U, side="right")};
forbidden items are rejected until $m$ admissible results are obtained.
Probabilities are computed in FP64. Disabled tasks, empty targets,
empty complements, and exact mode do not consume this stream.
The sample is not reduced when the training budget is exhausted: the run fails.

\subsection{Denominator estimation and derivatives}
\label{sec:sampled-loss}

For brevity, set $s_i=s_{qi}^{N}$.
The exact denominator and its random estimate are
\begin{align}\label{eq:sampled-normalizer}
Z_q&=\sum_{i\in T_q^a}e^{s_i}
       +\sum_{j\in\mathcal N_q^a}e^{s_j},\\
\widehat Z_q^a&=\sum_{i\in T_q^a}e^{s_i}
       +\frac1m\sum_{b=1}^{m}
        \frac{e^{s_{J_b}}}{Q_q^a(J_b)}.
\nonumber
\end{align}
Each positive target appears once, even when there are more than $m$.
If $\mathcal N_q^a=\varnothing$, no sampling occurs:
$\widehat Z_q^a=\sum_{i\in T_q^a}e^{s_i}=Z_q$.
For empty $T_q^a$, the original zero task mask applies;
no positive target is created artificially.

For a fixed model,
\begin{equation}\label{eq:sampled-unbiased-normalizer}
\E\!\left[\frac{e^{s_{J_b}}}{Q_q^a(J_b)}\right]
=\sum_{j\in\mathcal N_q^a}e^{s_j},
\qquad
\E[\widehat Z_q^a]=Z_q.
\end{equation}
The approximate loss for one request is
\begin{equation}\label{eq:sampled-query-loss}
\widehat\ell_q^a
=\log\widehat Z_q^a-\sum_{i\in T_q^a}t_{qi}^a s_i.
\end{equation}
All positive targets participate in the average.

Set $\kappa_i=1$ for $i\in T_q^a$ and
$\kappa_j=C_{qj}^a/(mQ_q^a(j))$ for $j\in\mathcal N_q^a$.
For a fixed sample,
\begin{equation}\label{eq:sampled-gradient}
\widehat p_i=\frac{\kappa_i e^{s_i}}{\widehat Z_q^a},
\qquad
\frac{\partial\widehat\ell_q^a}{\partial s_i}
=\widehat p_i-t_{qi}^a.
\end{equation}
Outside $T_q^a$, set $t_{qi}^a=0$. An unsampled negative item has $\kappa_i=0$;
$\sum_i\widehat p_i=1$, but this is the distribution of the training computation,
not the full catalog. Parameter derivatives follow the chain rule
with $Q_q^a$ and sample indices fixed.

To compute $\log\widehat Z_q^a$ stably, use scores
$s_i+\log\kappa_i$ for $\kappa_i>0$, subtracting the maximum.
The expanded correction is
\begin{equation}\label{eq:sampled-adjusted-score}
\widetilde s_i=
\begin{cases}
s_i, & i\in T_q^a,\\
s_i+\log C_{qi}^a-\log m-\log Q_q^a(i),
& i\in\mathcal N_q^a,\ C_{qi}^a>0.
\end{cases}
\end{equation}
For this scheme, $s_i-\log Q(i)$ is insufficient: it ignores multiplicity,
sample size, and exact inclusion of positive targets. The different corrections
for mandatory targets and sampled items matter~\cite{logq2025}.
Corrections apply only to the training denominator;
candidate retrieval uses the original $s_{qi}^N$.

Unbiasedness in~\eqref{eq:sampled-unbiased-normalizer}
does not imply an unbiased loss. Concavity of the logarithm gives
\begin{align}\label{eq:sampled-bias}
\E[\log\widehat Z_q^a]&\le\log Z_q,\\
\E[\nabla_\theta\widehat\ell_q^a]
&\ne\nabla_\theta\ell_q^a
\quad\text{in general}.\nonumber
\end{align}
Here $\theta=\theta_N$, and
$\ell_q^a=\log Z_q-\sum_i t_{qi}^as_i$.
Exact and approximate losses cannot be compared directly
when selecting an architecture.
For fixed parameters, a finite catalog, and positive
$Q_q^a$, the law of large numbers gives
$\widehat Z_q^a\to Z_q$ as $m\to\infty$.
The loss and its gradient at differentiable points converge as well.
Convergence at a fixed model does not guarantee identical
results from the two training processes.

\subsection{Training integration and cost}
\label{sec:sampled-protocol}

Loss $\widehat{\mathcal L}_a^N$ is obtained by replacing
the single-request loss in~\eqref{eq:modern-buy-loss}
with $\widehat\ell_q^a$, preserving $w_q^N$ and $a_q^a$.
In~\eqref{eq:modern-total-loss}, only
$\mathcal L_{\mathrm{buy}}^N$ and, when used,
$\mathcal L_{\mathrm{next}}^N$ are replaced.
Samples for these two tasks are constructed using their own targets.
Contrastive loss, graph regularization, and the parameter penalty are unchanged.
The approximation does not apply to the separate full-catalog SID loss.

The denominator requires at most $|T_q^a|+m$ dot products per request.
Direct computation of the allowed mass additionally requires
$O(|\mathcal N_q^a|)$ scalar operations per enabled task.
History, positive and sampled items, and their graph neighbors
are also recomputed using current parameters; stale
vectors are not substituted. The graph regularizer over all rows
and other shared components retain their full cost.

For the first approximate run, use $\delta=0.1$ and $m=512$;
$m\in\{1024,2048\}$ are tested if the selection budget on $C_{\mathrm{select}}$ allows it.
Any $m$ denotes sampling with replacement; full enumeration is exact mode.
Comparing them on a small shared catalog and identical sessions checks
computation and the cost of approximation, without guaranteeing transfer to the full catalog.

Modes follow protocol~\ref{sec:modern-time} with identical catalog,
quotas, and evaluation labels; Hit@20 and $\operatorname{Coverage}_{\mathrm{cand}}$
are computed without negative-item sampling.
The report includes epoch and total training time, update count,
peak memory, duplicate rate in the sample, and purchase hit rate.
Time savings are assessed together with the change in hit rate.

Hard negative items are allowed only as a separate hypothesis.
Applying~\eqref{eq:sampled-normalizer} requires the known actual
sampling probability of every item and positive support over all of
$\mathcal N_q^a$, for example through a uniform component. During gradient
computation, the distribution is fixed; dependence on parameters
and a different differentiation method require a separate derivation.
Arbitrarily adding items retrieved by an index while retaining the old $Q_q^a$
invalidates the justification of the denominator estimate.

\subsection{Computational checks}
\label{sec:sampled-checks}

Normative examples, derivatives, and boundary cases are collected
in X02 (Section~\ref{app:tests-extra}).

\section{Sequential model for a controlled comparison}
\label{sec:sequence-control}
\label{sec:sequence-overview}

After the neural stage, weighted history aggregation is compared
with a single causal attention block inspired by SASRec~\cite{sasrec2018}.
This compact adaptation replaces the session encoder while retaining the data,
purchase target, catalog, and ranker type. The item graph remains
in both models; its contribution is tested by separately disabling the graph branch.

\subsection{Identical inputs and groups of simultaneous events}
\label{sec:sequence-input}

The model uses the history $\mathcal D_q^N$ of available \code{detail} events
from Section~\ref{sec:modern-intents}, with the same $L_H$, handling
of unknown identifiers, and assumption about static vectors.
Sequences of cart, search, and page-view actions
require a separate comparison: they change both the architecture and the available inputs.
The ranker retains the cart features.

The item encoder is defined by
\eqref{eq:modern-edge-attention}--\eqref{eq:modern-item-fusion},
but is trained from scratch on $A_N$ jointly with the sequential encoder.
Denote its output by $e_i^{\mathrm{seq}}$.
The graph, paths, and catalog from $A_G$ are retained; the trained parameters
of the main model are not transferred. The search projection $z_q^{\mathrm{seq}}$ is obtained from
\eqref{eq:modern-query-vector} with its own trainable matrix
$T_{\mathrm{seq}}$, shared by the text branch and search.

Partition $\mathcal D_q^N$ into nonempty timestamp groups
$\mathcal G_{q1},\ldots,\mathcal G_{qT_q}$ with times
$t_{q1}<\cdots<t_{qT_q}\le t_q$, where $T_q\le L_H$.
Index $b$ denotes a group, not a row of the source file.
Define its input vector:
\begin{align}\label{eq:sequence-token}
n_{qb,i}^{\mathrm{detail}}
&=\sum_{\ell:\,e_{s(q)\ell}\in\mathcal H_q,\ t_{s(q)\ell}\le t_{qb}}
\ind\{i_{s(q)\ell}=i,\ \kappa_{s(q)\ell}=\mathrm{detail}\},\\
\bar e_{qb}&=\frac1{|\mathcal G_{qb}|}
\sum_{\ell\in\mathcal G_{qb}}e_{i_{s(q)\ell}}^{\mathrm{seq}},\nonumber\\
f_{qb}&=\frac1{|\mathcal G_{qb}|}
\sum_{\ell\in\mathcal G_{qb}}
\log(1+n_{qb,i_{s(q)\ell}}^{\mathrm{detail}}),\nonumber\\
x_{qb}&=F_{\mathrm{seq}}\bar e_{qb}+p_b
+a_{\mathrm{seq}}^t
\log\!\left(1+\frac{t_q-t_{qb}}{t_{\mathrm{unit}}}\right)
+a_{\mathrm{seq}}^f f_{qb}
+a_{\mathrm{seq}}^m\log|\mathcal G_{qb}|.\nonumber
\end{align}
Here $F_{\mathrm{seq}}\in\R^{d\times d}$;
$p_b,a_{\mathrm{seq}}^t,a_{\mathrm{seq}}^f,a_{\mathrm{seq}}^m\in\R^d$
are trained on $A_N$, and $p_1,\ldots,p_{L_H}$ is the group position table.
Frequency includes events through the end of the group; age is measured at $t_q$.
The contents of subsequent groups are unavailable. Row order within a group
does not affect the input; zero learned vectors do not remove events.

\subsection{A single causal attention block}
\label{sec:sequence-attention}

For $v\in\R^d$, define coordinate normalization as
\begin{equation}\label{eq:sequence-layernorm}
\operatorname{LN}_{\gamma,\beta}(v)
=\gamma\odot
\frac{v-\bar v\boldsymbol 1}
{\sqrt{d^{-1}\sum_{a=1}^d(v_a-\bar v)^2+\eps_{\mathrm{LN}}}}
+\beta,\qquad
\bar v=d^{-1}\sum_{a=1}^dv_a.
\end{equation}
The symbol $\odot$ denotes elementwise multiplication;
$\gamma,\beta\in\R^d$ are trainable; the initial $\eps_{\mathrm{LN}}=10^{-5}$
is included in the full configuration of the sequential variant.
The three applications $\operatorname{LN}_1,\operatorname{LN}_2,
\operatorname{LN}_3$ below have distinct trainable parameters.

For each existing group, compute the query, key, and value:
\begin{equation}\label{eq:sequence-qkv}
\tilde x_{qb}=\operatorname{LN}_1(x_{qb}),\qquad
Q_{qb}=W_Q\tilde x_{qb},\quad
K_{qb}=W_K\tilde x_{qb},\quad
V_{qb}=W_V\tilde x_{qb},
\end{equation}
where $W_Q,W_K,W_V\in\R^{d\times d}$.
A single attention head uses a causal mask: a group can access
itself and preceding groups:
\begin{align}\label{eq:sequence-attention}
a_{qbc}
&=\frac{\exp(Q_{qb}^{\top}K_{qc}/\sqrt d)}
{\sum_{v=1}^{b}\exp(Q_{qb}^{\top}K_{qv}/\sqrt d)},
&&1\le c\le b\le T_q,\\
a_{qbc}&=0,&&c>b.\nonumber
\end{align}
Dummy groups introduced by batch padding are excluded from the computation.
Attention to the group itself ensures a nonempty denominator.
Residual connections and an elementwise nonlinearity are then applied:
\begin{align}\label{eq:sequence-ffn}
y_{qb}&=x_{qb}+W_O\sum_{c=1}^{b}a_{qbc}V_{qc},\\
h_{qb}^{\mathrm{seq}}
&=y_{qb}
+W_2\operatorname{ReLU}\!\left(
W_1\operatorname{LN}_2(y_{qb})+b_1\right)+b_2.\nonumber
\end{align}
Here $W_O\in\R^{d\times d}$,
$W_1\in\R^{d_f\times d}$, $W_2\in\R^{d\times d_f}$,
$b_1\in\R^{d_f}$, $b_2\in\R^d$, and
$\operatorname{ReLU}(v)_a=\max(0,v_a)$.
The starting value $d_f=2d$ is a computational choice.

The history representation is read at the end of the observed prefix:
\begin{equation}\label{eq:sequence-readout}
u_q^{H,\mathrm{seq}}
=\mathsf n_\eps\!\left(
\operatorname{LN}_3(h_{qT_q}^{\mathrm{seq}})\right),
\qquad T_q>0.
\end{equation}
If $T_q=0$ and $M_q^S=1$, set
$u_q^{H,\mathrm{seq}}=z_q^{\mathrm{seq}}$;
if both history and search are absent, set it to zero.
Search is added using the same bounded mixing rule
\eqref{eq:modern-search-gate}, replacing $u_q^I$ with
$u_q^{H,\mathrm{seq}}$ and $z_q^N$ with $z_q^{\mathrm{seq}}$.
The mixing coefficient is trained from scratch with the same $\bar g_N$
and initialization. The result $u_q^{\mathrm{seq}}$, when only
search is present, equals $z_q^{\mathrm{seq}}$. Zero norms do not change
the fixed mask of the training example.

\subsection{Training, candidates, and comparison conditions}
\label{sec:sequence-training}

For the same $i\in\mathcal I_q$, define
\begin{equation}\label{eq:sequence-score}
s_{qi}^{\mathrm{seq}}
=\frac{\langle u_q^{\mathrm{seq}},e_i^{\mathrm{seq}}\rangle}{\tau_N}.
\end{equation}
Purchase targets, session weights, and auxiliary next views
are the same as in Sections~\ref{sec:modern-objective}--\ref{sec:modern-next}.
Training minimizes the sum~\eqref{eq:modern-total-loss},
with the encoder and its parameters replaced, following
the recipe in~\ref{sec:modern-training-recipe}.
Graph regularization acts on this model's own item encoder.
In the first comparison, both architectures use $\lambda_C=0$,
the same normalization method, number of negative items, and training
budget. If the contrastive component is subsequently enabled, the example
transformations and pair selection rules are also identical. The architectures
may have different parameter counts.

After training, the item vectors are fixed and indexed.
One query $u_q^{\mathrm{seq}}$ to the index returns up to $m_N=90$
items according to~\eqref{eq:sequence-score}
whose $e_i^{\mathrm{seq}}$ norm exceeds $\eps_N$.
This is the entire neural quota; the baseline sources retain 210 items
under the overall limit $M=300$.
If the query norm is at most $\eps_N$, the neural source is empty;
deduplication and filling of the remainder follow the main policy.
Candidate eligibility and the prohibition on inserting targets follow
\ref{sec:model-candidates}.

The ranker receives the score~\eqref{eq:sequence-score}, masked search
similarity, the mixing coefficient, and candidate provenance and rank.
Multi-intent features are not attributed to this model.
Ranker retraining and selection follow~\ref{app:ablation};
the architectures are compared using the same criterion and a predefined budget.

\subsection{Cost and testable properties}
\label{sec:sequence-cost}

Grouping requires on the order of $L_Hd$ operations on precomputed vectors.
For one prefix, the block requires
$O(T_qd^2+T_q^2d+T_qdd_f)$ operations
and $O(T_q^2+T_qd+T_qd_f)$ intermediate memory.
With the starting $d_f=2d$, the latter expression is $O(T_q^2+T_qd)$.
Search projection, item vector reads,
index lookup, and final ranking are accounted for separately.
The item graph and its paths are not recomputed for each request.
Training cost includes the same recomputation of item representations
and the same selected method for computing the denominator.
Standard key caching across requests is not assumed here:
event age depends on $t_q$, and truncating the history changes positions.

Tests verify that the response is unchanged when rows within one timestamp
group are permuted or events after $t_q$ are modified,
that attention ignores batch padding, and that the overall budget $M$ is preserved.
Permuting items across different times may change the result.
Purchase coverage, training time, memory, and measured latency are published.

\section{Aging of relations and an updated graph}\label{sec:freshness}

An updated graph is a separate policy that accounts for changes in relations
between eligible items. The main comparison retains the $A_G$ snapshot
and the temporal protocol in~\ref{sec:eval-temporal}.

\subsection{Two distinct time scales}\label{sec:freshness-weights}

Parameter $\tau_G$ in~\eqref{eq:model-session-edge} describes the distance
between actions within one session. Define the aging of recorded sessions
using a window $T_{\mathrm{win}}>0$ and a scale $\tau_{\mathrm{age}}>0$ in seconds.

Let $c_k$ be the data cutoff for the $k$th update,
and let $r_k>c_k$ be the actual publication time of the completed package.
Denote by $\mathcal S^{\mathrm{obs}}(c_k)$ the selected historical
sessions satisfying $t_{0,s}+H+g<c_k$ and
$t_{\mathrm{last},s}^{\mathrm{obs}}+g<c_k$.
All published records of a selected session have been read and are available by $c_k$;
a session with a later tail is excluded in its entirety. These conditions do not establish
telemetry completeness. The same rule applies to sessions with and without a recorded
purchase: a known target does not permit early inclusion.

Retain the pairs from~\eqref{eq:model-event-pairs} for which both events
fall within the calendar window and both items belong to
$\mathcal I_G$. For $i,j\in\mathcal I_G$, define
\begin{equation}\label{eq:freshness-pairs}
\mathcal P^{r,\mathrm{fr}}_{sij}(c_k)
=\{(\ell,m)\in\mathcal P^r_{sij}(c_k):
c_k-T_{\mathrm{win}}\le t_{s\ell}<t_{sm}<c_k\}.
\end{equation}
For an eligible pair, introduce
\begin{equation}\label{eq:freshness-pair-weight}
\xi_{s\ell m}(c_k)=
\exp\!\left(-\frac{t_{sm}-t_{s\ell}}{\tau_G}
            -\frac{c_k-t_{sm}}{\tau_{\mathrm{age}}}\right).
\end{equation}
The weights and support are
\begin{align}
w^{r,\mathrm{fr}}_{sij}(c_k)
&=
\begin{cases}
\displaystyle\max_{(\ell,m)\in\mathcal P^{r,\mathrm{fr}}_{sij}(c_k)}
\xi_{s\ell m}(c_k),
&\mathcal P^{r,\mathrm{fr}}_{sij}(c_k)\ne\varnothing,\\
0,&\text{otherwise},
\end{cases}\label{eq:freshness-session}\\
W^{r,\mathrm{fr}}_{ij}(c_k)
&=\sum_{s\in\mathcal S^{\mathrm{obs}}(c_k)}
w^{r,\mathrm{fr}}_{sij}(c_k),\nonumber\\
n^{r,\mathrm{fr}}_{ij}(c_k)
&=\sum_{s\in\mathcal S^{\mathrm{obs}}(c_k)}
\ind\{w^{r,\mathrm{fr}}_{sij}(c_k)>0\}.
\label{eq:freshness-weight}
\end{align}
The bound $0\le w^{r,\mathrm{fr}}_{sij}\le1$ is preserved and, consequently,
$0\le W^{r,\mathrm{fr}}_{ij}\le n^{r,\mathrm{fr}}_{ij}$.
Support counts distinct sessions in the current window, not all sessions
over the lifetime of the graph. Sparsification, threshold $n_G$,
normalization, and row regularization apply to these new quantities.
Equality~\eqref{eq:model-row-regularization} still holds.
A larger $\tau_{\mathrm{age}}$ weakens aging but does not remove
the hard window boundary.

\subsection{Publishing a consistent version}\label{sec:freshness-publication}

The catalog $\mathcal I_G$ is retained. New items require a separate
eligibility rule based on first observation time, parameter initialization,
and compatibility checks. A static metadata file does not establish
past item availability.

The graph, bounded paths, item vectors,
indexes, features, and a compatible ranker are published together.
A new graph can change vectors even with unchanged neural weights.
The previous package remains in use until the new one passes validation; a request selects
a version by publication time $r_k<t_q$, not by data cutoff $c_k$.

In the first mode, the trained model parameters are fixed.
The ranker is trained and validated in advance on earlier periods
by replaying the adopted graph update rule.
This validation is performed separately from that of the static model.

In the second mode, the required components are retrained on new
historical blocks while preserving their roles. The package is published after
all stages. Parameters, training targets, and validation results are available
before $r_k$; calendar eligibility of targets retains the $H+g$ margin.

\subsection{Comparison protocol and checks}\label{sec:freshness-evaluation}

The updating policy is evaluated on a separate subsequent interval
$D_{\mathrm{roll}}$. Before it is opened, the update schedule,
$T_{\mathrm{win}},\tau_{\mathrm{age}}$, training mode, publication delay
rule, eligible data for each release, and fallback package are fixed.
These quantities are selected on earlier intervals.
Period $D$ of the main static comparison is not used
for their tuning or retraining.
Training on early sessions in $D_{\mathrm{roll}}$ is allowed only after
calendar eligibility is satisfied; the results are unavailable to preceding requests.
The evaluation concerns a predefined learning policy.

The static and updating policies are replayed on the same
requests from $D_{\mathrm{roll}}$ with identical $K,M,H,Q_{\max}$.
To separate the effect of decay from the benefit of new history, weights
are also compared on the same data window.
Published results include Hit@20, $\operatorname{Coverage}_{\mathrm{cand}}$, data age $t_q-c_k$,
publication delay $r_k-c_k$, recomputation cost, and the number of missed
updates. In an A/B test, training uses a predefined shared external stream
or separate data from each group. Mixing these data requires an analysis
of interactions between groups.

X01 checks time, weights, and the immutability of the published graph
(Section~\ref{app:tests-extra}).

\section{Candidate generator based on semantic codes}
\label{sec:sid-overview}

A semantic identifier (SID) encodes an item as a sequence
of indices into small codebooks. An additional source after the neural
stage predicts these indices one at a time.
The principles borrowed from DIGER and GRID are listed in~\ref{sec:research}.

The experiment uses Coveo vectors, positional soft quantization,
and a recurrent decoder: each next index depends on the previous ones.
It differs from DIGER in using no language encoder, residual
quantization, or adaptive noise adjustment. Code assignments are trained
with a separate full-catalog recommendation loss; the graph
encoder, ranker, and business criterion are retained.

\subsection{Training data and fixed representations}
\label{sec:sid-data}

The catalog $\mathcal I_G$ and frequencies are fixed using $A_G$.
The neural module trained on $A_N$ is frozen;
its $e_i\in\R^d$ and $u_q^N\in\R^d$ become SID inputs.
SID is also trained on $A_N$, which does not serve as a generalization check.
Dimensions, epoch counts, and coefficients are selected on $C_{\mathrm{select}}$,
without fitting trainable parameters to this period.

Let $x_i^{\mathrm{SID}}$ combine normalized description
and image vectors with availability masks; a missing modality yields
a zero block. Categories and raw images are not required.
The fixed input is
\begin{equation}\label{eq:sid-fixed-input}
 z_i^{\mathrm{SID}}=[e_i;x_i^{\mathrm{SID}}].
\end{equation}
Concatenation $[\,;\,]$ uses static data
subject to the limitation in Section~\ref{sec:data}.

For a request $q$ in the training block, set
$Y_q^G=Y_q\cap\mathcal I_q\cap\mathcal I_G$.
The set $\mathcal Q_N^{\mathrm{SID}}$ contains the first requests of sessions
$\mathcal S_{A_N}^{+}$ with $Y_q^G\ne\varnothing$. Set
$w_q^{\mathrm{SID}}=1$ and
$W_N^{\mathrm{SID}}=\sum_{q\in\mathcal Q_N^{\mathrm{SID}}}w_q^{\mathrm{SID}}$.
Each session has total weight one. If the catalog is empty
or $W_N^{\mathrm{SID}}=0$, training is skipped with status
\code{insufficient}. Requests without recorded purchases do not generate
training targets and do not enter Hit@20. Purchases outside $\mathcal I_G$
do not expand the catalog and remain possible misses during evaluation.

\subsection{Soft assignments and the hard code}
\label{sec:sid-codes}

Let $L_{\mathrm{SID}}$ be the code length,
$V_{\mathrm{SID}}\ge2$ the size of each codebook,
and $d_{\mathrm{SID}}$ the dimension of its vectors.
For position $\ell$, train a matrix $A_\ell^{\mathrm{SID}}$,
a bias $b_\ell^{\mathrm{SID}}$, and vectors
$c_{\ell b}^{\mathrm{SID}}\in\R^{d_{\mathrm{SID}}}$,
$b=1,\ldots,V_{\mathrm{SID}}$. The assignment logits are
\begin{equation}\label{eq:sid-logits}
 a_{i\ell b}^{\mathrm{SID}}
 =-\norm{A_\ell^{\mathrm{SID}}z_i^{\mathrm{SID}}
       +b_\ell^{\mathrm{SID}}-c_{\ell b}^{\mathrm{SID}}}^{2}.
\end{equation}
Tied maxima are resolved by the smaller codebook entry index.
At training step $t$, draw independent
$U_{i\ell b}\sim\mathrm{Uniform}(0,1)$ and
$G_{i\ell b}=-\log(-\log U_{i\ell b})$.
For numerical stability, $U$ is clipped to
$[\eps_U,1-\eps_U]$, where $0<\eps_U<1/2$.
Soft assignments are given by
\begin{equation}\label{eq:sid-soft}
 \rho_{i\ell b}^{\mathrm{SID}}(t)=
 \frac{\exp\{(a_{i\ell b}^{\mathrm{SID}}
                 +\sigma_t^{\mathrm{SID}}G_{i\ell b})/
                    \tau_t^{\mathrm{SID}}\}}
 {\sum_{v=1}^{V_{\mathrm{SID}}}
  \exp\{(a_{i\ell v}^{\mathrm{SID}}
                 +\sigma_t^{\mathrm{SID}}G_{i\ell v})/
                    \tau_t^{\mathrm{SID}}\}}.
\end{equation}
Here $\tau_t^{\mathrm{SID}}>0$ controls the softness of the distribution,
and $\sigma_t^{\mathrm{SID}}\ge0$ is the noise scale.
All soft losses within a step use the same noise for an item.
Noise is reproduced using stream $j=3$ of the profile in
\ref{sec:modern-training-recipe}: each joint step traverses the entire
catalog in UTF-8 ID order, then $\ell,b$, using \code{random()} and
$\eps_U=10^{-7}$. Noise logarithms are computed in FP64 and the result is cast
to FP32. Schedule~\eqref{eq:app-sid-schedule} uses
$t=1,\ldots,T_{\mathrm{SID}}$, where $T_{\mathrm{SID}}$ is the number of steps
in three joint epochs, computed from the saved batch plan.
During pretraining, $\sigma=0,\tau=1$; no noise is generated.

The positional vectors and their concatenation are
\begin{equation}\label{eq:sid-soft-vector}
 \widetilde c_{i\ell}^{\mathrm{SID}}
 =\sum_b\rho_{i\ell b}^{\mathrm{SID}}c_{\ell b}^{\mathrm{SID}},
 \qquad
 \widetilde v_i^{\mathrm{SID}}
 =[\widetilde c_{i1}^{\mathrm{SID}};\ldots;
   \widetilde c_{iL_{\mathrm{SID}}}^{\mathrm{SID}}].
\end{equation}
Concatenation preserves the distinction between positions. The serving index
uses only deterministic codes
\begin{equation}\label{eq:sid-hard}
 k_{i\ell}^{\mathrm{SID}}=\argmax_b a_{i\ell b}^{\mathrm{SID}},
 \qquad
 \mathbf k_i^{\mathrm{SID}}
 =(k_{i1}^{\mathrm{SID}},\ldots,k_{iL_{\mathrm{SID}}}^{\mathrm{SID}}).
\end{equation}
Noise is disabled during serving and historical evaluation.

\subsection{Losses and gradient flow to assignments}
\label{sec:sid-loss}

The trainable matrix $R^{\mathrm{SID}}$ maps the session vector
to the space of concatenated codes. For $\tau_R^{\mathrm{SID}}>0$,
define the full-catalog distribution
\begin{align}\label{eq:sid-full-prob}
 a_{qi}^{R,\mathrm{SID}}
 &=\frac{(R^{\mathrm{SID}}u_q^N)^\top
                  \widetilde v_i^{\mathrm{SID}}}
                {\tau_R^{\mathrm{SID}}},\\
 p_{qi}^{R,\mathrm{SID}}
 &=\frac{\exp(a_{qi}^{R,\mathrm{SID}})}
 {\sum_{j\in\mathcal I_q\cap\mathcal I_G}
             \exp(a_{qj}^{R,\mathrm{SID}})}.
\end{align}
The future-purchase loss is
\begin{equation}\label{eq:sid-rec-loss}
 \mathcal L_R^{\mathrm{SID}}
 =-\frac{1}{W_N^{\mathrm{SID}}}
 \sum_{q\in\mathcal Q_N^{\mathrm{SID}}}
 \frac{w_q^{\mathrm{SID}}}{|Y_q^G|}
       \sum_{i\in Y_q^G}\log p_{qi}^{R,\mathrm{SID}}.
\end{equation}
The remaining items form the normalization catalog;
the absence of a purchase record does not imply a negative response to an impression.

For one request and target distribution
$\pi_{qi}^{\mathrm{SID}}=\ind\{i\in Y_q^G\}/|Y_q^G|$,
the derivative with respect to the logit is
$p_{qi}^{R,\mathrm{SID}}-\pi_{qi}^{\mathrm{SID}}$.
Since
\begin{equation}\label{eq:sid-soft-gradient}
 \frac{\partial\rho_{i\ell b}^{\mathrm{SID}}}
      {\partial a_{i\ell v}^{\mathrm{SID}}}
 =\frac{\rho_{i\ell b}^{\mathrm{SID}}}
            {\tau_t^{\mathrm{SID}}}
       \bigl(\ind\{b=v\}-\rho_{i\ell v}^{\mathrm{SID}}\bigr),
\end{equation}
there is a derivative chain
$\mathcal L_R^{\mathrm{SID}}\to\widetilde v_i^{\mathrm{SID}}
\to\rho_{i\ell b}^{\mathrm{SID}}\to a_{i\ell b}^{\mathrm{SID}}
\to A_\ell^{\mathrm{SID}},c_{\ell b}^{\mathrm{SID}}$.
This derivative trains soft assignments but does not differentiate
the discrete $\argmax$. An improvement in the soft loss does not guarantee
an improvement in hard retrieval.

Linear decoders $D_m^{\mathrm{SID}}$, $m\in\{\mathrm{text},\mathrm{image}\}$,
are used to preserve modality information.
Let $N_m=\sum_{i\in\mathcal I_G}M_i^m$. Define
\begin{equation}\label{eq:sid-reconstruction}
 \mathcal L_X^{\mathrm{SID}}
 =\sum_{m:N_m>0}\frac{1}{N_m}
 \sum_{i\in\mathcal I_G}M_i^m
 \norm{D_m^{\mathrm{SID}}\widetilde v_i^{\mathrm{SID}}
                        -\bar v_i^m}^{2}.
\end{equation}
If a modality is entirely absent, its term is zero.
The supplied vectors are reconstructed because the original text
and images are unavailable.

Average code utilization and the balancing loss are defined as
\begin{equation}\label{eq:sid-balance}
 \bar\rho_{\ell b}^{\mathrm{SID}}
 =\frac{1}{|\mathcal I_G|}\sum_{i\in\mathcal I_G}
          \rho_{i\ell b}^{\mathrm{SID}},
 \qquad
 \mathcal L_B^{\mathrm{SID}}
 =\frac{1}{L_{\mathrm{SID}}}\sum_{\ell,b}
  \bar\rho_{\ell b}^{\mathrm{SID}}
            \log(V_{\mathrm{SID}}\bar\rho_{\ell b}^{\mathrm{SID}}).
\end{equation}
For soft assignments, gradient computation retains
$r_{\ell b}=\operatorname{LSE}_{i\in\mathcal I_G}
(\log\rho_{i\ell b}^{\mathrm{SID}})-\log|\mathcal I_G|$,
where $\log\rho$ is obtained directly via \code{logsoftmax}.
The balancing term is computed as $\exp(r_{\ell b})
(\log V_{\mathrm{SID}}+r_{\ell b})$ following~\ref{sec:modern-log-probabilities}.
In control variants with fixed codes, balance is computed from
nonzero integer frequencies; zero cells are skipped,
and the balancing-loss gradient is zero.
The regularizer does not guarantee diversity of hard codes: uniform
soft assignments yield zero loss, yet after tie-breaking
the entire catalog may receive one code. Its utilization and collisions
are checked under~\ref{sec:sid-validation}; excessive regularization
may destroy useful item groups.

\subsection{Autoregressive decoder and training order}
\label{sec:sid-decoder}

The decoder state $h_{q\ell}^{\mathrm{SID}}\in\R^{d_h^{\mathrm{SID}}}$
is initialized from the session vector. For a partially generated code
$k_1,\ldots,k_{\ell-1}$, define
\begin{align}\label{eq:sid-decoder}
 h_{q1}^{\mathrm{SID}}&=\tanh(U^{\mathrm{SID}}u_q^N+b_0^{\mathrm{SID}}),\\
 P_\ell^{\mathrm{SID}}(b\given q,k_{<\ell})
 &=\frac{\exp\{(O_\ell^{\mathrm{SID}}h_{q\ell}^{\mathrm{SID}}
                         +o_\ell^{\mathrm{SID}})_b\}}
 {\sum_v\exp\{(O_\ell^{\mathrm{SID}}h_{q\ell}^{\mathrm{SID}}
                         +o_\ell^{\mathrm{SID}})_v\}},\\
 h_{q,\ell+1}^{\mathrm{SID}}
 &=\tanh(F^{\mathrm{SID}}h_{q\ell}^{\mathrm{SID}}
       +T^{\mathrm{SID}}c_{\ell k_\ell}^{\mathrm{SID}}
       +p_\ell^{\mathrm{SID}}).
\end{align}
The function $\tanh$ is applied to each coordinate.
All matrices and biases in these formulas are trainable.
The probability of a complete sequence is the product
of the conditional probabilities of its positions.

At the start of an epoch, hard codes~\eqref{eq:sid-hard}
are computed for all items and fixed as the targets for that epoch.
When training the next position, the decoder receives the correct
previous code. The loss is
\begin{equation}\label{eq:sid-gen-loss}
 \mathcal L_G^{\mathrm{SID}}
 =-\frac{1}{W_N^{\mathrm{SID}}L_{\mathrm{SID}}}
 \sum_{q\in\mathcal Q_N^{\mathrm{SID}}}
 \frac{w_q^{\mathrm{SID}}}{|Y_q^G|}
 \sum_{i\in Y_q^G}\sum_{\ell=1}^{L_{\mathrm{SID}}}
 \log P_\ell^{\mathrm{SID}}
       (k_{i\ell}^{\mathrm{SID}}\given q,k_{i,<\ell}^{\mathrm{SID}}).
\end{equation}
Loss~\eqref{eq:sid-gen-loss} does not differentiate target indices;
the connection to $A_\ell^{\mathrm{SID}}$ and $b_\ell^{\mathrm{SID}}$ is provided by
\eqref{eq:sid-rec-loss}. The gradient through the decoder input vectors
changes tables $c_{\ell b}^{\mathrm{SID}}$ and possible future assignments,
but is not a derivative with respect to a code index.

The joint objective is
\begin{align}\label{eq:sid-objective}
 \mathcal L^{\mathrm{SID}}
 &=\mathcal L_G^{\mathrm{SID}}
  +\lambda_R^{\mathrm{SID}}\mathcal L_R^{\mathrm{SID}}
  +\lambda_X^{\mathrm{SID}}\mathcal L_X^{\mathrm{SID}}
  \nonumber\\
 &\quad+\lambda_B^{\mathrm{SID}}\mathcal L_B^{\mathrm{SID}}
  +\lambda_2^{\mathrm{SID}}\norm{\theta_{\mathrm{SID}}}^{2},
\end{align}
where $\lambda_R^{\mathrm{SID}}>0$, the other coefficients are nonnegative,
and $\theta_{\mathrm{SID}}$ contains only this module's parameters.
The profile in~\ref{sec:modern-training-recipe} applies to all three phases.
In pretraining, only $A_\ell^{\mathrm{SID}}$,
$b_\ell^{\mathrm{SID}}$, $c_{\ell b}^{\mathrm{SID}}$, and $D_m^{\mathrm{SID}}$ are trained;
all other SID parameters are frozen. One epoch is run on
$\mathcal L_X^{\mathrm{SID}}+\lambda_B^{\mathrm{SID}}\mathcal L_B^{\mathrm{SID}}$:
the catalog, initially sorted by UTF-8 ID, is shuffled by calling
\code{permutation} on the ordering stream; batches contain up to 256 items.
The sum of item reconstruction terms with global $N_m$
is multiplied by the number of batches in the epoch; balance is computed over the full catalog.
All parameters $\theta_{\mathrm{SID}}$ are then trained for three joint
epochs on~\eqref{eq:sid-objective},
in batches of up to 64 requests from distinct sessions. Purchase and generative
terms retain global weights and the factor $T_{\mathrm{ep}}$,
as in~\eqref{eq:modern-batch-loss}; reconstruction, balance, and L2
are computed in full at every step.
At the end, the encoder and code tables are fixed, all
hard codes are recomputed, and only the decoder is then trained further
on \eqref{eq:sid-gen-loss} using the same $A_N$ and the final targets.
In this final phase, only $U^{\mathrm{SID}},b_0^{\mathrm{SID}}$,
$O_\ell^{\mathrm{SID}},o_\ell^{\mathrm{SID}},F^{\mathrm{SID}}$,
$T^{\mathrm{SID}},p_\ell^{\mathrm{SID}}$ from~\eqref{eq:sid-decoder} are trained.
It lasts two epochs with the same batching rule
and global normalization; the other losses are disabled.
At the start of each phase, the Adam moments and step counter are reset,
while generator states continue. The lists of trainable and frozen
tensors for each phase are saved in the run manifest; the original neural network
remains frozen throughout all SID phases.
After decoder training, the entire module is frozen before processing $B$.

\subsection{Valid sequences, collisions, and output size}
\label{sec:sid-serving}

Different items may receive the same code.
For a sequence $\mathbf k$, define its item bucket
\begin{equation}\label{eq:sid-bucket}
 \mathcal B^{\mathrm{SID}}(\mathbf k)
 =\{i\in\mathcal I_G:\mathbf k_i^{\mathrm{SID}}=\mathbf k\}.
\end{equation}
In the mapping table, items are sorted by popularity from $A_G$,
then by identifier. The served subset
$\mathcal B_R^{\mathrm{SID}}(\mathbf k)$ contains the first $R_{\mathrm{SID}}$
items; the excluded fraction is reported separately. An item is uniquely
identified by the pair ``code, position in the table,'' since codes may repeat.

A prefix tree contains catalog codes: an edge specifies the next
index, and a path to a leaf represents an existing code. The decoder uses only
edges of this tree.
At depth $\ell$, at most $B_{\mathrm{SID}}$ paths are retained
with the largest sums of original log probabilities
from \eqref{eq:sid-decoder}. Probabilities are not renormalized
after an edge is prohibited, because this would change the path scores being compared.
Ties are resolved lexicographically. Bounded beam width
does not guarantee the globally best $B_{\mathrm{SID}}$ sequences.

Let $\mathcal K_q^{\mathrm{SID}}$ be the retrieved leaves.
From the set
\begin{equation}\label{eq:sid-candidates}
 \mathcal U_q^{\mathrm{SID}}
 =\mathcal I_q\cap\bigcup_{\mathbf k\in\mathcal K_q^{\mathrm{SID}}}
                     \mathcal B_R^{\mathrm{SID}}(\mathbf k)
\end{equation}
select the first $m_{\mathrm{SID}}$ distinct items by score
$\sum_\ell\log P_\ell^{\mathrm{SID}}
(k_{i\ell}^{\mathrm{SID}}\given q,k_{i,<\ell}^{\mathrm{SID}})$,
then by the same historical popularity and identifier.
Consequently, $|\mathcal C_q^{(\mathrm{SID})}|
\le\min(m_{\mathrm{SID}},B_{\mathrm{SID}}R_{\mathrm{SID}})$.
Quota $m_{\mathrm{SID}}$ is included in the existing limit $M$;
the usual sources and the common filling rule are retained.

Additional features may include source ranks,
the code log probability, and the mapping table size. Code probability
does not estimate purchase probability or the causal effect of an impression. Ranker
retraining follows~\ref{app:ablation}, candidate constraints follow
\ref{sec:model-candidates}, and failure handling follows~\ref{sec:modern-cost}.

The generator is enabled for serving only after measured validation
of latency, memory, and cost on the target node. If requirements are violated,
the source is excluded; research results are retained separately.

\subsection{Diagnostics, comparisons, and computational limits}
\label{sec:sid-validation}

After noise is disabled, measurements include the proportions
$f_{\ell b}^{\mathrm{SID}}=|\mathcal I_G|^{-1}
\sum_i\ind\{k_{i\ell}^{\mathrm{SID}}=b\}$,
the effective number of codes
$\exp\{-\sum_b f_{\ell b}^{\mathrm{SID}}\log f_{\ell b}^{\mathrm{SID}}\}$,
the fraction of distinct complete codes, and the fraction of items beyond
$R_{\mathrm{SID}}$. Zero entropy terms are defined as zero.
The soft loss is compared with the quality of the hard source.
Utilization, collision, and runtime thresholds are fixed before the run;
the regularizer does not guarantee compliance.

Diagnostics cover the entire $\mathcal I_G$ before applying the
$R_{\mathrm{SID}}$ limit. Let $N_G=|\mathcal I_G|>0$,
$b_{\mathbf k}=|\mathcal B^{\mathrm{SID}}(\mathbf k)|$,
with the sum below taken over all nonempty mapping tables:
\begin{align}
\operatorname{Col}_{\mathrm{item}}^{\mathrm{SID}}
&=\frac1{N_G}\sum_{\mathbf k}
b_{\mathbf k}\ind\{b_{\mathbf k}>1\},\label{eq:sid-catalog-collisions}\\
\operatorname{Drop}_{R}^{\mathrm{SID}}
&=\frac1{N_G}\sum_{\mathbf k}
\max(0,b_{\mathbf k}-R_{\mathrm{SID}}).\nonumber
\end{align}
These measures distinguish nonunique codes from item exclusion;
zero $\operatorname{Drop}_{R}^{\mathrm{SID}}$ does not rule out collisions.
If $N_G>V_{\mathrm{SID}}^{L_{\mathrm{SID}}}$, at least
$N_G-V_{\mathrm{SID}}^{L_{\mathrm{SID}}}$ assignments repeat;
capacity is increased for a control without repetitions.
The reverse inequality permits unique codes but does not guarantee them.

Three variants are required with the same $M,L_{\mathrm{SID}},V_{\mathrm{SID}}$
and tuning budget: random fixed codes, codes obtained after
reconstruction without recommendation fine-tuning, and jointly trained
codes. In the first two variants, the assignment table is fixed:
code vectors may be trained as decoder input representations,
but changing them does not recompute the assigned indices.
The vector formulas for these variants replace soft assignments
with $\rho_{i\ell b}^{\mathrm{SID}}=\ind\{b=k_{i\ell}^{\mathrm{SID}}\}$.
The random variant requires
$V_{\mathrm{SID}}^{L_{\mathrm{SID}}}\ge|\mathcal I_G|$.
Stream $m=2,j=6$ calls
\code{choice(C, size=N_G, replace=False, shuffle=True)} once, where
$C=V_{\mathrm{SID}}^{L_{\mathrm{SID}}}\le2^{63}-1$ and $N_G=|\mathcal I_G|$.
This uniformly samples distinct numbers from $0,\ldots,C-1$ without replacement.
The numbers are assigned to items
in identifier order and represented in base
$V_{\mathrm{SID}}$, padded with leading zeros to length
$L_{\mathrm{SID}}$; one is added to each digit.
This tests the value of semantic organization. When collision
rates differ, the accessible fraction of the catalog is also compared.

X04--X05 check derivatives, training phases, output validity, and collisions
(Section~\ref{app:tests-extra}).

Code table memory is
$O(L_{\mathrm{SID}}V_{\mathrm{SID}}d_{\mathrm{SID}})$,
and index memory is $O(|\mathcal I_G|L_{\mathrm{SID}})$.
For a minibatch of $b_N$ requests, full-catalog normalization requires
$O(b_N|\mathcal I_G|L_{\mathrm{SID}}d_{\mathrm{SID}})$ operations
after the soft item vectors have been computed. The catalog is processed
in blocks with exact accumulation of the normalization sum: this saves memory
but does not reduce the number of dot products.
Decoding cost is bounded by
$O(L_{\mathrm{SID}}B_{\mathrm{SID}}
[(d_h^{\mathrm{SID}})^2+V_{\mathrm{SID}}d_h^{\mathrm{SID}}
+d_h^{\mathrm{SID}}d_{\mathrm{SID}}])$,
and the number of item mappings read is
$B_{\mathrm{SID}}R_{\mathrm{SID}}$.
Codebook size and beam width are varied under the same latency limit.
Admission to the served configuration is checked under~\ref{sec:sid-serving};
historical coverage and causal effect are evaluated separately.

\section{Additional retrieval and slate construction methods}\label{sec:model-extensions}

These variants are not part of the first run. Each changes a separate part
of the system and is compared with a fixed baseline under the same candidate limit.

\subsection{Downweighting edges supported by few sessions}\label{sec:model-edge-optional}

A separate research comparison may further reduce
the contribution of edges supported by few sessions:
\begin{align}\label{eq:model-edge-ablation}
\chi^r_{ij}&=\frac{n^r_{ij}}{n^r_{ij}+\lambda_{\mathrm{edge}}},
\qquad \lambda_{\mathrm{edge}}\ge0,\nonumber\\
g^{r,\mathrm{edge}}_{qj}
&=\sum_{\ell\in\mathcal D_q}\alpha_{q\ell}
\rho^r_{i_{s(q)\ell}}\chi^r_{i_{s(q)\ell},j}
P^r_{i_{s(q)\ell},j}.
\end{align}
Coefficient $\chi^r_{ij}$ is computed only for retained edges,
where $n^r_{ij}\ge1$; setting $\lambda_{\mathrm{edge}}=0$ restores
the original score. This variant retains the same $P^r$, topology,
paths, and candidate quota, but changes graph scores and requires
retraining the ranker on $B$. It may suppress useful rare
relations, so it is not enabled automatically in the main configuration.

\subsection{Additional slate diversity}\label{sec:model-diversity-optional}

An additional policy may reduce excessive similarity between items.
The following probabilistic variant is outside the current formulation:
it would require separately defined and validated estimates
$\widehat p_{qi}$ of purchase probability. Scores $f(x_{qi})$ are not substitutes for them.
Let $S_{ij}\in[0,1]$ be the symmetric similarity of valid
content vectors: negative cosine similarities are replaced by zero,
and available modalities are combined in a normalized mixture.
If no modality is shared, $S_{ij}=0$ with an unknown-value mask.
In this case, zero does not establish that the items differ.
For a set $L\subseteq\mathcal C_q$, define the surrogate
\begin{equation}\label{eq:model-slate-surrogate}
U_q(L)=\sum_{i\in L}\widehat p_{qi}
-\frac{\lambda_D}{\max(1,K-1)}
\sum_{\{i,j\}\subseteq L}S_{ij},\qquad\lambda_D\ge0.
\end{equation}
The first term favors high observed purchase frequencies,
while the second penalizes repeated inclusion of similar items.
$U_q$ is not measured in monetary units and is not conversion.
The main version does not compute this objective: the list is sorted
by the ranker score, with $\lambda_D=0$. The probabilistic extension
requires a separate training and selection plan and a new final period;
tuning on the results of $D$ is prohibited.

For $\lambda_D>0$, items are added greedily:
\begin{equation}\label{eq:model-greedy-slate}
i_k=\argmax_{i\in\mathcal C_q\setminus L_{k-1}}
\left[\widehat p_{qi}-
\frac{\lambda_D}{\max(1,K-1)}\sum_{j\in L_{k-1}}S_{ij}\right],
\quad L_k=L_{k-1}\cup\{i_k\},\quad L_0=\varnothing.
\end{equation}
Addition continues through $k=K_q$, including when the increment
is negative: the protocol fixes the output length.
This algorithm is a heuristic; the global maximum
of~\eqref{eq:model-slate-surrogate} is not guaranteed.
Sequential construction costs on the order of
$K|\mathcal C_q|$ similarity computations when accumulated sums are updated.

\subsection{Why the sum of item probabilities is not conversion}

Even if the true probabilities of a future purchase
of every item under the same condition $\mathcal H_q$ were known, purchases
of different items are dependent events. Only the general bounds
below hold for the probability of at least one purchase of an item in $L$:
\begin{equation}\label{eq:model-union-bounds}
\max_{i\in L}p_i\le
\Prob\!\left(\bigcup_{i\in L}\{y_{qi}=1\}\given\mathcal H_q\right)
\le\min\!\left(1,\sum_{i\in L}p_i\right).
\end{equation}
The formula $1-\prod_{i\in L}(1-p_i)$ requires conditional independence,
which is generally not justified for substitute and complementary items.
Estimates $\widehat p_{qi}$ are additionally conditional on the score and the candidate
procedure, so substituting them into these bounds is only
illustrative and does not produce guaranteed probability bounds.
None of these formulas estimates the causal increase in purchases
from changing the recommendation policy.

\subsection{Multiple graph steps}

The one-step graph admits a testable extension using a random
walk with restart at the session history.
Let $T$ be a row-normalized mixture of relation matrices;
zero rows are replaced by the historical distribution $\pi^{\mathrm{pop}}$.
Let $h_q$ be the initial weight distribution over history
items, obtained by summing $\alpha_{q\ell}$.
On the shared snapshot catalog, define
\begin{equation}\label{eq:model-ppr}
z_q^{(n+1)}=\eta_Rh_q+(1-\eta_R)T^\top z_q^{(n)},
\qquad z_q^{(0)}=h_q,\quad 0<\eta_R\le1.
\end{equation}
At each step, part of the mass returns to the known history,
while the remainder follows edges. If the known history is empty,
this source is disabled. Items unknown to the graph receive no
fictitious initial vertices; the remaining weights are renormalized.

Because $T$ is row-stochastic, the transformation contracts
distance in the sum-of-absolute-values norm by a factor of $1-\eta_R$.
Consequently, a unique limiting distribution exists,
and for $n\ge1$ it satisfies
\begin{equation}\label{eq:model-ppr-error}
\norm{z_q^{(n)}-z_q^*}_1\le2(1-\eta_R)^n.
\end{equation}
This bounds computational error, not recommendation quality.
Transitions are performed on the sparse matrix;
the contribution of replaced zero rows is computed by one multiplication
of vector $\pi^{\mathrm{pop}}$ by their total mass.
A full step with an explicit vector costs $O(N_G+E_{\mathrm{sparse}})$,
where $N_G$ is the catalog size and $E_{\mathrm{sparse}}$ the number of retained
edges before filling zero rows. A linear scan is needed even with no edges,
so repeated scans of the full catalog may exceed the latency
budget. The source is accepted only after measuring its additional
purchase coverage and cost relative to one step.

\subsection{Behavioral item vectors}

Item2Vec can be studied as an additional source
using co-occurring interactions~\cite{item2vec}.
Let $\mathcal P_A$ be item pairs from temporal and positional
windows in period $A_N$, restricted to catalog $\mathcal I_G$. Each item has
two learned vectors $u_i^{B},v_i^{B}\in\R^{d_B}$: a center vector and a context vector.
One possible training objective is
\begin{align}\label{eq:model-item2vec}
\mathcal L_B=&-\sum_{(i,j)\in\mathcal P_A}
\left[\log\sigmoid((u_i^{B})^\top v_j^{B})
+m_B\E_{k\sim Q_B}\log\sigmoid(-(u_i^{B})^\top v_k^{B})\right]
\nonumber\\
&+\frac{\lambda_B}{2}\sum_i(\norm{u_i^{B}}^2+\norm{v_i^{B}}^2),
\qquad Q_B(k)=\frac{n_k^{3/4}}{\sum_jn_j^{3/4}}.
\end{align}
Here $n_k$ is the number of interactions with the item in $A_G$,
$m_B$ is the number of random contexts per positive pair,
and $\lambda_B>0$ is regularization. A random context helps
distinguish characteristic neighborhoods from background frequency;
it does not assert that the user rejected item $k$.
The expectation is replaced by an average over independent
samples from the fixed $Q_B$.

Let $\mathcal D_q^B\subseteq\mathcal D_q$ contain views of items
with learned behavioral vectors. When this set is nonempty,
the query representation is
\begin{equation}\label{eq:model-item2vec-query}
u_q^B=\frac{\sum_{\ell\in\mathcal D_q^B}
\alpha_{q\ell}u_{i_{s(q)\ell}}^B}
{\sum_{\ell\in\mathcal D_q^B}\alpha_{q\ell}},
\qquad s_{qi}^B=(u_q^B)^\top v_i^B.
\end{equation}
If $\mathcal D_q^B$ is empty, the source is unavailable.
Using different center and context vectors
is consistent with the training objective. Replacing the dot product
with cosine similarity or averaging $u_i^{B},v_i^{B}$ defines a separate variant
that requires evaluation. Behavioral and content vectors
are not mixed coordinatewise without a specifically trained mapping.

Retrieval extensions are compared with the same overall budget
$M$ and the same request schedule. Added complexity
is accepted only with a reproducible improvement in the selected
purchase coverage measure and compliance with the serving budget.
Their purpose is to measure the contribution of each retrieval method,
not to increase the number of mandatory system components.

\section{Experiment with real impressions}\label{sec:eval-online}

This protocol is deferred for the public dataset. Before serving real impressions,
it is required, together with checks of historical quality, technical
limits on the target hardware, and outcome recording. The control is
the current policy or no recommendation block; it is fixed
before launch \cite{experiments}.

\subsection{Assignment unit and causal effect}\label{sec:eval-itt}

Before launch, fix the assignment unit, eligible population,
exposure horizon, and corresponding analysis unit.
If the service has a stable permitted account/user/device key,
assign variant $Z_u\in\{0,1\}$ by this key at first eligibility,
with $\Prob(Z_u=1)=1/2$, and retain it throughout the experiment and outcome maturation.
$1$ denotes the new system and $0$ the control. All sessions with the same key
inherit the assignment; the assignment algorithm version, key, and result
are recorded before calling the ranker. Rules for key changes, multiple
devices, and cookie loss are fixed in advance; these cases must not allow
hidden rerandomization. Eligibility is determined before exposure.

The primary session metric retains the unit from~\eqref{eq:eval-itt-estimand}.
The confirmatory cohort includes only the first eligible session
$s(u)$ of each key new to the experiment; this rule is applied
before assignment. The primary outcome $Y_{s(u)}(z)$ is a purchase after assignment
in that same session before $t_{0,s(u)}+H$. Each key contributes exactly one
session to the primary denominator:
\begin{align}\label{eq:eval-unit-itt}
\Delta_{1000,\mathrm{entry}}^{\mathrm{ITT}}
&=1000\bigl(\E[Y_{s(u)}(1)]-\E[Y_{s(u)}(0)]\bigr),\nonumber\\
\widehat\Delta_{1000,\mathrm{entry}}^{\mathrm{ITT}}
&=1000\left(\frac{\sum_{u:Z_u=1}Y_{s(u)}}{N_1}
             -\frac{\sum_{u:Z_u=0}Y_{s(u)}}{N_0}\right),\quad N_0,N_1>0.
\end{align}
Here $N_z$ is both the number of assigned keys and the number of entry-cohort
sessions; the unit is purchasing entry sessions per thousand, not all visits
or people. A device key is not identical to a person. Later sessions
inherit the variant but do not add to the primary denominator. The effect on
any purchase by the key within $H$, and the ratio of purchasing sessions to all
subsequent sessions, are separate secondary estimands only;
the number of subsequent sessions may itself change under treatment.
They do not replace the primary criterion $S_{\mathrm{business}}$.
For the entry cohort, the formulas below are applied with $s=s(u)$ and $n_z=N_z$
only if outcomes are independent and identically distributed across keys
and there is no substantial interference between keys. If these conditions
are not justified, an analysis by assignment units accounting for
calendar dependence is specified in advance; an independent binomial SE is not used.
The threshold, power, and covered traffic fraction are specified for this cohort.

Only when a stable key is unavailable is the limited
short-term session protocol below permissible. At the first eligible request,
assign $Z_s\in\{0,1\}$, with $\Prob(Z_s=1)=1/2$, and retain the variant
through the end of the session and its outcome window. It does not estimate a user-level effect.

Let $Y_s(z)$ be the potential outcome: at least one purchase
after assignment in the same session and before $t_{0,s}+H$, if the session is assigned
variant $z$. For any specific session, the only observed outcome is
\begin{equation}\label{eq:eval-observed-outcome}
 Y_s=Z_sY_s(1)+(1-Z_s)Y_s(0).
\end{equation}
\begin{assumption}[Conditions for the session experiment]
\label{ass:eval-randomization}
Assignment is independent of potential outcomes.
The policy within each variant is unambiguously defined.
Assignment in one session does not substantially affect another session's outcome:
for example, through users remembering recommendations, shared item
scarcity, or popularity changed by the experiment.
Purchase recording is equally complete in both variants.
Independent binomial intervals additionally require
independent and identically distributed outcomes across sessions within
each group; randomization alone does not guarantee this condition.
\end{assumption}

Coveo has no stable person identifier across sessions, and
this does not determine key availability in a future service. For session
assignment, permitted privacy-safe aggregate signals of repeat visits
and a sensitivity plan are defined before launch. The fraction of likely
repeat visits is a diagnostic, not proof of the absence of carryover effects.
If upper bounds $r_z$ are justified for the fractions of outcomes potentially
affected by cross-session interference in group $z$, then for binary outcomes
the difference interval is expanded on both sides by $1000(r_0+r_1)$,
and intersected with $[-1000,1000]$. Data-estimated bounds require separate
accounting for their uncertainty and the overall error level. Without
a justified bound and a plausible no-interference assumption,
no confirmatory session effect is claimed; the result is descriptive.

Under Assumption~\ref{ass:eval-randomization}, for the entry
cohort or the limited session protocol, the effect and estimator are
\begin{align}
 \Delta_{1000}^{\mathrm{ITT}}
 &=1000\bigl(\E[Y_s(1)]-\E[Y_s(0)]\bigr),
 \label{eq:eval-itt-estimand}\\
 \widehat\Delta_{1000}^{\mathrm{ITT}}
 &=1000\left(\frac{k_1}{n_1}-\frac{k_0}{n_0}\right),
 \quad n_z=\sum_s\ind\{Z_s=z\},\quad
 k_z=\sum_s\ind\{Z_s=z\}Y_s,\qquad n_0,n_1>0.
 \label{eq:eval-itt-estimator}
\end{align}
ITT is analysis by original assignment. Sessions with errors,
empty responses, or no actual impression remain in $n_z$; $Y_s$ is determined
by a real purchase, not serving success. All purchases are counted
without attributing an order to an impression. Selecting only those who clicked or saw
the block breaks the original randomized comparison.

\subsection{Interval for the business effect}\label{sec:eval-itt-ci}

The following formulas apply to the stated primary cohort with independent
outcomes and a stopping point fixed in advance. For $n_0,n_1>0$ and
sufficiently large counts of purchasing and nonpurchasing sessions
in each group, use the approximation
\begin{align}
 \widehat p_z&=k_z/n_z,\qquad
 \widehat{\operatorname{se}}_{\mathrm{ITT}}
 =1000\sqrt{\frac{\widehat p_1(1-\widehat p_1)}{n_1}
              +\frac{\widehat p_0(1-\widehat p_0)}{n_0}},
 \label{eq:eval-itt-se}\\
 I_{\mathrm{normal}}
 &=\left[\widehat\Delta_{1000}^{\mathrm{ITT}}
       -z_{1-\alpha/2}\widehat{\operatorname{se}}_{\mathrm{ITT}},
         \widehat\Delta_{1000}^{\mathrm{ITT}}
       +z_{1-\alpha/2}\widehat{\operatorname{se}}_{\mathrm{ITT}}\right].
 \label{eq:eval-itt-normal-ci}
\end{align}
Here $z_u$ is a standard normal quantile.
The confidence level is not the probability that a hypothesis is true.
If a group is empty, the estimate and interval are undefined: status
\texttt{insufficient}, with no substitution of a zero effect.

For rare outcomes, a conservative binomial
interval can be selected in advance. For $k$ purchasing sessions out of $n>0$, $0\le k\le n$,
the bounds invert binomial tails using a validated implementation
of beta-distribution quantiles:
\begin{subequations}\label{eq:eval-exact-beta}
\begin{align}
\ell&=\begin{cases}
0,&k=0,\\
Q_{\mathrm{Beta}(k,n-k+1)}(\alpha/4),&k>0,
\end{cases}\label{eq:eval-exact-lower}\\
u&=\begin{cases}
1,&k=n,\\
Q_{\mathrm{Beta}(k+1,n-k)}(1-\alpha/4),&k<n.
\end{cases}\label{eq:eval-exact-upper}
\end{align}
\end{subequations}
Here $0<\alpha<1$, and $Q$ is a quantile of the specified distribution;
the result is checked for finiteness and $0\le\ell\le u\le1$.
For two groups, the conservative effect interval is
\begin{equation}\label{eq:eval-exact-diff}
 I_{\mathrm{exact}}=1000[\ell_1-u_0,\;u_1-\ell_0].
\end{equation}
Noncoverage of each group's probability is at most $\alpha/2$,
and that of the difference is at most $\alpha$ by the union bound.
The independence conditions are given in~\ref{sec:eval-itt}
and Assumption~\ref{ass:eval-randomization}; multiple sessions from one key
require cluster analysis outside this binomial contract.

\subsection{Business result acceptance rule}\label{sec:eval-acceptance}

Let $[\ell_B,u_B]$ be the prespecified confidence interval
for $\Delta_{1000}^{\mathrm{ITT}}$ after the experiment is completed.
Two claims are distinguished:
\begin{align}
 \ell_B>0
 &\quad\Longrightarrow\quad
 \text{a positive effect is confirmed},
 \label{eq:eval-positive-acceptance}\\
 \ell_B>\delta_{\mathrm{accept}}
 &\quad\Longrightarrow\quad
 \text{the business threshold is confirmed to be exceeded}.
 \label{eq:eval-business-acceptance}
\end{align}
Threshold $\delta_{\mathrm{accept}}\ge0$, in purchasing sessions per thousand,
is set in advance. The claims hold under the assumptions and confidence level
of the selected procedure; a point estimate exceeding the threshold is insufficient.
The other acceptance conditions are specified in Section~\ref{sec:acceptance}.
With key-based assignment, these conditions apply to
$\Delta_{1000,\mathrm{entry}}^{\mathrm{ITT}}$ and a threshold in purchasing
entry sessions per thousand; a secondary user-level outcome
does not replace this session criterion.

\subsection{Experiment size and duration}\label{sec:eval-power}

Before launch, specify a detectable increase of
$\delta_{\mathrm{plan}}>\delta_{\mathrm{accept}}$ purchasing sessions per thousand, error level $\alpha$,
and desired power $1-\beta$: the probability of detecting an effect
of the assumed size. Let $p_0$ be the expected control rate,
$\delta=\delta_{\mathrm{plan}}/1000$, $p_1=p_0+\delta$, and
$\bar p=(p_0+p_1)/2$.
For equal groups testing a zero effect, the approximate number
of independent sessions in each group is
\begin{equation}\label{eq:eval-sample-size}
 n\approx\frac{
 \left[z_{1-\alpha/2}\sqrt{2\bar p(1-\bar p)}
       +z_{1-\beta}\sqrt{p_0(1-p_0)+p_1(1-p_1)}\right]^2}
 {\delta^2}.
\end{equation}
The approximation is checked for applicability to the expected outcome rarity
and analysis method. $p_0$ is obtained from a comparable stream or pilot;
the Coveo rate does not transfer automatically to another store.
With key-based assignment, plan the number of entry sessions, equal to the number
of new keys; $p_0$ also refers to the entry cohort. If
these outcomes are demonstrably independent, the formula applies, but substituting the count
of all subsequent sessions is not allowed.

Formula~\eqref{eq:eval-sample-size} plans a test of no effect,
not a lower bound above $\delta_{\mathrm{accept}}$.
For the business threshold, test
$H_0:\Delta_{1000}^{\mathrm{ITT}}\le\delta_{\mathrm{accept}}$
against $\delta_{\mathrm{plan}}>\delta_{\mathrm{accept}}$;
the required size depends on the gap between these quantities.

The lower-bound rule for a normal interval gives a simple
planning approximation. Denote $d_*=(\delta_{\mathrm{plan}}-
\delta_{\mathrm{accept}})/1000$ and
$v_*=p_0(1-p_0)+p_1(1-p_1)$, where $p_1=p_0+\delta_{\mathrm{plan}}/1000$
must lie in $(0,1)$. With equal groups, the standard error
of the difference in proportions is approximately $\sqrt{v_*/n}$. Therefore,
\begin{equation}\label{eq:eval-threshold-power}
n\approx\frac{v_*(z_{1-\alpha/2}+z_{1-\beta})^2}{d_*^2}.
\end{equation}
Condition $\sqrt{n}\,d_*/\sqrt{v_*}-z_{1-\alpha/2}=z_{1-\beta}$
plans for the lower bound to exceed the threshold when the standard error
is close to its planned value. For rare outcomes, the power of the selected rule is checked
by summing probabilities of pairs $(k_0,k_1)$ or by simulation.

At $v$ newly eligible primary-cohort units per second,
the approximate enrollment time is $2n/v$,
followed by $H+g$ of observation. The duration also covers the specified
calendar cycles; varying traffic and session dependence affect power.

The date or rule for ending enrollment is fixed in advance; analysis is performed
after all outcomes mature, following~\ref{sec:eval-srm-stop}.

\subsection{Assignment check and stopping rule}
\label{sec:eval-srm-stop}

Before the first assignment, save a plan of $L$ calendar checks,
the SRM control level $\alpha_{\mathrm{SRM}}=0.001$, and expected shares of $1/2$.
Each check counts unique assigned units, including
errors and missing impressions. If the plan includes $G$ tested strata,
membership $S_u=\phi(X_u^{\mathrm{pre}})$ is computed before assignment
from features available at that point. The value $S_u$ is recorded
in the assignment log and never changed. The plan specifies fields $X_u^{\mathrm{pre}}$,
the version and hash of rule $\phi$, missing-value handling, and the expected share
of each variant in each stratum. Clicks, impressions, fallback serving,
and post-assignment outcomes do not define SRM strata; such breakdowns
are permissible only for separately justified analysis.

The standard mode requires independent assignments with
$\Prob(Z_u=1\mid X_u^{\mathrm{pre}})=1/2$. Each scheduled check
additionally requires that, conditional on the composition and size $N$ of the overall cohort
or each tested stratum, the number of assignments to variant 1 follows
$\mathrm{Bin}(N,1/2)$. Composition is not adapted to assignment history
or experimental exposure. Under these conditions, apply an exact
two-sided binomial test; $N$ includes all assigned members
of the tested population regardless of serving and outcome.
Without strata, each check uses threshold $\alpha_{\mathrm{SRM}}/L$;
with $G$ strata, the threshold for every test, including the overall test, is
$\alpha_{\mathrm{SRM}}/[L(G+1)]$. Dependent assignment or adaptive
population composition requires a prespecified test of the actual
mechanism with overall error control, not a binomial substitution.
The plan contains the overall cohort composition rule, the test method,
and justification of the conditional distribution and overall error control.

SRM means a mismatch in assignment proportions. When it is detected,
new enrollment is paused, all assignments are retained, and missing
logs, repeated keys, eligibility changes, and routing errors are investigated.
Outcomes are not deleted to restore the proportion. Additional checks cover
balance of prespecified pre-exposure features, integrity of assignment
logs, and agreement between assignment and actual serving.
Fallback does not change the ITT variant. Post-impression balance is not used
to filter participants; the absence of SRM proves neither correct
purchase recording nor the absence of interference.

The safety monitor receives only predeclared measures;
before launch, fix their windows, denominators, minimum sample sizes, thresholds,
response actions, and check schedule. Crossing an emergency
threshold stops exposure; the availability of such a channel does not by itself
make the final interval valid under data-dependent stopping.
The main protocol has one final analysis with no interim
efficacy decisions. Any data-dependent early
stop, including a planned safety stop or a stop
due to SRM, as well as any unplanned interruption, makes the causal result
descriptive: ordinary intervals and $p$-values are not reported as
confirmatory, and rules~\eqref{eq:eval-positive-acceptance}
and~\eqref{eq:eval-business-acceptance} are not applied. After investigation, a new
preregistered experiment is required for a new confirmatory conclusion.
Confirmatory analysis may be retained after planned interim
stops only under a procedure separately specified before launch
with overall error control (for example, alpha-spending or always-valid);
it cannot be chosen retrospectively.

\subsection{Incomplete outcomes and safety measures}
\label{sec:eval-missing}

A mature calendar horizon does not guarantee telemetry completeness.
If group $z$ contains $m_z$ sessions with unknown outcomes,
and $k_z^{\mathrm{obs}}$ purchases are reliably recorded among the remainder,
then for $n_z>0$, without additional assumptions,
\begin{equation}\label{eq:eval-missing-rate}
 \frac{k_z^{\mathrm{obs}}}{n_z}\le p_z^{\mathrm{sample}}
 \le\frac{k_z^{\mathrm{obs}}+m_z}{n_z}.
\end{equation}
The corresponding range of compatible sample effects is
\begin{equation}\label{eq:eval-missing-effect}
 1000\left[
 \frac{k_1^{\mathrm{obs}}}{n_1}
 -\frac{k_0^{\mathrm{obs}}+m_0}{n_0},\quad
 \frac{k_1^{\mathrm{obs}}+m_1}{n_1}
 -\frac{k_0^{\mathrm{obs}}}{n_0}\right].
\end{equation}
This is a range of compatible results, not a confidence interval:
statistical uncertainty is assessed separately. Unknown outcomes
are not replaced by zeros or removed from the ITT denominator without justification.

Before the experiment, specify limits on deterioration in error rate, latency,
empty outputs, purchase recording, and the fraction of ineligible items.
Stopping on these measures is permitted; the status of the causal analysis
is governed by Section~\ref{sec:eval-srm-stop}. Unknown outcomes,
maturity rules, and completeness are checked at the same assignment unit:
with key-based assignment, $n_z,m_z,k_z^{\mathrm{obs}}$ refer to
the entry sessions of those keys, not all their sessions or purchases.

Session length, clicks, and cart additions are reported as secondary
results. Multiple confirmatory effects require a prespecified
multiplicity correction. Breakdowns by post-assignment features
do not admit a causal interpretation without separate justification.

\section{Why historical data do not identify an increase in purchases}
\label{sec:eval-identification}

Nonidentifiability is established in Proposition~\ref{prop:business-nonidentification}.
For example, Hit@20 of an empty output is zero, so a higher value from the new
system does not prove an increase in purchases relative to no recommendations.

\subsection{Conditions for evaluating another policy from an impression log}
\label{sec:eval-ope}

Evaluating a new policy from the old policy's log in a single-decision setting
requires records $(X_s,A_s^{\mathrm{act}},Y_s)$: pre-assignment context,
the list actually displayed, and the outcome. The old policy's probability
$b(a\mid x)$ of choosing the list must be known.

This requires no unobserved causes of list selection that also affect the outcome,
correct impression recording, and the support condition:
\begin{equation}\label{eq:eval-support}
 \pi(a\mid x)>0\ \Longrightarrow\ b(a\mid x)>0.
\end{equation}
Under these conditions, the inverse-propensity estimator is
\begin{equation}\label{eq:eval-ips}
 \widehat V_{\mathrm{IPS}}(\pi)
 =\frac1N\sum_s
  \frac{\pi(A_s^{\mathrm{act}}\mid X_s)}
       {b(A_s^{\mathrm{act}}\mid X_s)}Y_s.
\end{equation}
Rare lists with small probability $b$ receive large weights,
so variance may be unacceptable \cite{ips}.
Weight clipping reduces variance at the cost of bias.

If an estimate
$\widehat m(x,a)\approx\E[Y\mid X=x,A^{\mathrm{act}}=a]$
has also been trained, the outcome model and weights can be combined:
\begin{align}
 \widehat V_{\mathrm{DR}}(\pi)
 =\frac1N\sum_s\Bigg[&\sum_a\pi(a\mid X_s)\widehat m(X_s,a)
 \notag\\[-0.2em]
 &+\frac{\pi(A_s^{\mathrm{act}}\mid X_s)}
       {b(A_s^{\mathrm{act}}\mid X_s)}
       \bigl(Y_s-\widehat m(X_s,A_s^{\mathrm{act}})\bigr)\Bigg].
 \label{eq:eval-dr}
\end{align}
The outcome model is trained without the evaluated record. Double robustness
means consistency when at least one of the two
components is correctly specified, under causal identification and regularity conditions.
It does not resolve a lack of support \cite{dr}.

Formulas~\eqref{eq:eval-ips}--\eqref{eq:eval-dr} apply to a single decision
and a fixed subsequent regime. Sequential recommendations
change future context; substituting individual requests with a shared
session purchase requires a different justification.

Coveo lacks the required assignment probabilities and a complete recommendation
impression log. Search results with clicks do not replace them, so IPS
and DR are excluded from the main calculation. Predicted purchase probabilities
also do not replace assignment probabilities $b$.

\subsection{What event replay can establish}
\label{sec:eval-replay}

Replay checks feature availability, session state,
candidate retrieval, latency, and constraints. A simulated user
tests scenarios under a declared behavior model; neither method
establishes an increase in real purchases.

Test~\ref{test:eval-effect} checks the distinction between coverage and effect.


\section{Run configurations and parameter registry}\label{app:parameters}

\subsection{Fixed configuration for the first run}\label{app:first-config}

The first run is GCS+Rk: three graph relations, descriptions, images,
the search query, and LambdaMART. It uses direct relations and an unchanged
graph under \code{observed_purchases_v1}, with no neural
sources or parameter tuning. The settings below are
initial settings, not measured optima on Coveo.

\paragraph{Data and time.}
Before training, fix the boundaries of the seven parts, session selection,
file versions, and identifier order. The graph, catalog, and statistics
are built on $A_G$; the ranker is trained on $B$. The reserved $A_N$
is not used by the baseline. The first report is produced on $C_{\mathrm{smoke}}$;
$C_{\mathrm{select}},E,D$ remain closed.
\begin{align}
H&=86\,400\ \text{s},&g&=1800\ \text{s},&K&=20,\notag\\
Q_{\max}&=1,&M&=300,&d_G&=100,\notag\\
\Delta_G&=1800\ \text{s},&\tau_G&=600\ \text{s},&n_G&=3,\notag\\
\tau_H&=300\ \text{s},&\Delta_S&=600\ \text{s},&\lambda_G&=10.
\label{eq:app-start}
\end{align}
The graph uses $B_G=10$ preceding item timestamp groups;
its session representation uses the $L_H=20$ latest views
under Section~\ref{sec:model-graph}. Store $L_{\mathrm{hist}}=20$
timestamp groups; popularity smoothing is $\eta_P=1$.
The purchase-rate feature is not used.

\paragraph{Retrieval and features.}
Search over description, image, and query vectors is exact
over the full eligible catalog. Quotas are fixed:
\begin{center}
\begin{tabular}{L{0.63\textwidth}r}
\toprule Source & Quota \\\midrule
Each of the three graph relations & 50 \\
Item description & 40 \\
Image & 30 \\
Search query & 50 \\
Recent views & 20 \\
Historical popularity & 10 \\
\midrule Total & 300 \\
\bottomrule
\end{tabular}
\end{center}
Deduplication, unavailable-source masks, and filling
by popularity follow~\eqref{eq:model-candidate-fill}.
The ranker receives separate graph and content scores;
mixtures with coefficients $\omega_r$ and $\gamma_m$ are not used.
The additional factor $\chi_{ij}^r$ is disabled.
Double-precision tolerances are $\eps_v=10^{-12}$,
$\eps_f=10^{-8}$. The ordering of tied scores and categorical encodings
are fixed together with the feature schema.

\paragraph{Ranker and output.}
The LightGBM version, derivative function, and all settings are specified
in~\ref{sec:model-ranker-implementation}. The first-run budget
is not selected on $C_{\mathrm{select}}$. The similarity penalty is disabled: $\lambda_D=0$.

\paragraph{Validation and experiment completion.}
The first run ends with the checks in~\ref{app:tests-base}
and the report on $C_{\mathrm{smoke}}$ specified in~\ref{app:diagnostic-output}.
Compliance with load requirements under exact search is measured;
approximate search is introduced through a separate comparison before $D$ is opened.

Calibration is disabled. Final evaluation follows
Sections~\ref{sec:eval-observed} and~\ref{sec:eval-primary-comparison}.
The initial $\alpha=0.05$ and
$B_{\mathrm{boot}}=2000$ do not replace the precision plan and conditions
for applying temporal intervals.

\subsection{Neural extension parameters}\label{app:neural-config}

The graph, catalog, and statistics from $A_G$ are retained; representations are trained
on $A_N$. Quotas are redistributed under the overall limit $M=300$:
\begin{longtable}{L{0.38\textwidth}L{0.15\textwidth}L{0.15\textwidth}L{0.17\textwidth}}
\toprule Source & Baseline & Neural & With SID \\\midrule\endhead
Each of the three graph relations & 50 & 30 & 30 \\
Item description & 40 & 30 & 30 \\
Image & 30 & 20 & 20 \\
Search query & 50 & 40 & 40 \\
Recent views & 20 & 20 & 20 \\
Popularity & 10 & 10 & 10 \\
Mixed neural vector & 0 & 30 & 20 \\
Each of the two intents & 0 & 30 & 10 \\
Semantic code & 0 & 0 & 50 \\
\midrule Total & 300 & 300 & 300 \\
\bottomrule
\end{longtable}
An inactive intent returns an empty list. After deduplication,
vacant positions are filled by popularity according to
formula~\eqref{eq:model-candidate-fill}; its initial quota
does not limit additional filling.
In a subsequent evaluation with $J_I=4$, each intent has a quota of 15
for the neural architecture and 5 for the SID variant; the total neural
quotas of 90 and 40 are retained. With $J_I=1$, the single direction
receives 60 and 20 positions, respectively, in addition to the mixed source.


The initial configuration is
\begin{align}
d&=64,&J_I&=2,&H_P&=2,&d_P&=50,\notag\\
B_P&=100,&R_P&=20,&c_G&=0.1,&\eta_N&=1,\notag\\
\tau_P&=1,&\tau_N&=0.1,&\tau_C&=0.1,&\tau_{SN}&=300\ \text{s}.
\label{eq:app-modern-start}
\end{align}
For the three relations, initially $\zeta_r=1/3$; the direct search
mixing limit is $\bar g_N=0.5$. The initial actual weight when search
is available is $g_0=0.1$: initially,
$v_g=w_g=0$, $a_g^t=0$, and the bias is
\begin{equation}\label{eq:app-search-init}
b_g=\log\frac{g_0}{\bar g_N-g_0}=\log\frac14.
\end{equation}
When the search signal is absent, the mask still gives
$g_q^{SN}=0$. These values do not separately limit the influence
of $\beta_S$ on the intent-weight distribution.
Loss weights are
\begin{equation}\label{eq:app-modern-loss-start}
\lambda_{\mathrm{next}}=0.2,\qquad \lambda_C=0.05,\qquad
\lambda_G^N=10^{-3},\qquad \lambda_\theta=10^{-5}.
\end{equation}
In the initial extension, $p_S=0$: random search dropout
is not applied. The diagnostic threshold for intent agreement
is fixed at $\delta_I=0.05$.
For 32-bit floating-point computation, the initial value is
$\eps_N=10^{-6}$. The probability of dropping each available branch
when constructing contrastive representations is $0.2$, with branch
restoration following Section~\ref{sec:modern-contrastive}.
The optimizer, initialization, batches, and three training epochs are specified
in~\ref{sec:modern-training-recipe}. Checkpoints after one and five epochs
are reserved for a separately planned continuation; evaluating them
also counts toward the selection budget on $C_{\mathrm{select}}$.

The cost of an exact epoch is measured using the actual $|\mathcal I_G|$
and number of training requests. If the budget is exceeded, use the separately
identified mode in Appendix~\ref{sec:sampled-training}: $m=512$,
$\delta=0.1$. Values $m=1024,2048$ require a separate plan;
a sampled denominator changes the training procedure.


\subsection{Initial settings for the deferred SID study}\label{app:sid-config}

The deferred study in Appendix~\ref{sec:sid-overview} uses
\begin{align}
L_{\mathrm{SID}}&=4,&V_{\mathrm{SID}}&=32,&d_{\mathrm{SID}}&=16,\notag\\
d_h^{\mathrm{SID}}&=64,&B_{\mathrm{SID}}&=32,&R_{\mathrm{SID}}&=4,
\label{eq:app-sid-start}
\end{align}
and $m_{\mathrm{SID}}=50$. Initial weights are
$\lambda_R^{\mathrm{SID}}=1$,
$\lambda_X^{\mathrm{SID}}=0.1$,
$\lambda_B^{\mathrm{SID}}=0.01$,
$\lambda_2^{\mathrm{SID}}=10^{-5}$; $\tau_R^{\mathrm{SID}}=1$.
For $T_{\mathrm{SID}}$ joint steps specified in advance, the following
schedule is permissible:
\begin{equation}\label{eq:app-sid-schedule}
\sigma_t^{\mathrm{SID}}=\pos{1-t/T_{\mathrm{SID}}},\qquad
\tau_t^{\mathrm{SID}}=0.2+0.8\pos{1-t/T_{\mathrm{SID}}}.
\end{equation}
Training phases are specified in~\ref{sec:modern-training-recipe}
and Section~\ref{sec:sid-decoder}. Hard-code utilization is checked
over the entire catalog: $32^4=1\,048\,576$ possible codes do not rule out
collisions with fewer items. Increasing $L_{\mathrm{SID}}$
or $V_{\mathrm{SID}}$ must be justified by measured retrieval losses.


\subsection{Run manifest and stage checks}
\label{app:run-manifest}

One run manifest stores the fixed inputs,
code, and experiment settings in JSON; stage results are stored separately
and refer to this manifest. A separate manifest snapshot and SHA-256 file
for every stage are not mandatory. Before opening $C_{\mathrm{select}}$,
$D$, or assigning A/B variants, the corresponding plan is saved
immutably; a version or checksum reference links it
to the results. Before computation, verify that the code, settings, and inputs
match the plan; a result cannot change a plan that has already been fixed.
The normative template
\code{documents/math/coveo_session_recommender_run_manifest.template.json}
is not ready to run: unknown dates, hashes, and measurements are
\code{null}. They must not be replaced with invented data.
Template \code{manifest_version=2.10} specifies the mode through
\code{evaluation_protocol}; the run manifest fixes
the first-request rule, observed target, evaluation population,
precision thresholds in proportions, and disabled components
under Section~\ref{sec:eval-observed}. Unknown coverage is not replaced with complete coverage.
Types, states, temporal relationships, input integrity, and conditions of the applicable stage
are validated with ordinary data checks and tests. A separate test
registration system is unnecessary; the validating code version is reproducible.
Check B13 includes valid and invalid run manifests.

\begin{longtable}{L{0.24\textwidth}L{0.66\textwidth}}
\toprule Field group & Required contents \\\midrule\endhead
Data & Versions and SHA-256 of the three tables, schema, vector dimensions,
observation coverage and gaps; UTC and integer milliseconds,
canonicalization and session subsampling rules. \\
Periods & Boundaries of all seven parts as $[a,b)$; nonoverlap,
whole sessions, and eligibility~\eqref{eq:eval-admission}.
Calendar fields have the prefix \code{period_}; matrices have \code{matrix_}. \\
Code and models & Source code hash, pinned dependency file
with package versions and hashes, feature schema, derivative function,
full settings of active modules, hardware, numeric types,
random-number generators and their states. \\
Selection and evaluation & \code{evaluation_protocol=observed_purchases_v1},
variants and selection rule, primary comparison,
metadata sensitivity, precision and temporal dependence plan.
Requirements are determined by the stage table below. \\
Serving & Load and limits, delivery and recovery rules,
trusted sources, retention periods, monitoring, responsible role;
required for service validation and admission to A/B testing. \\
Checks & Standard tests mapped to requirement IDs; their code
defines the input or recipe, an independent reference, and tolerances. Separate stage
results contain the actual outcome, code version, configuration, and seed.
A separate hashed manifest for every test is unnecessary. \\
\bottomrule
\end{longtable}

\paragraph{Stage conditions.}\label{app:run-stages}
Field \code{stage} specifies the stage being validated. The table determines
required fields; data for future and inapplicable stages are not required.

Each result record contains the requirement/test ID, state,
reason, actual report or a link to it, and reproducible
code version, configuration, and seed. \code{null} is allowed before the execution deadline,
but does not replace a required result when the stage is completed.
For a disabled module, record \code{not_used} and a reason;
the same deadlines apply when it is enabled.
$E$ is reserved and unused; calibration is disabled.

\begin{longtable}{L{0.18\textwidth}L{0.72\textwidth}}
\toprule Stage & Completion conditions \\\midrule\endhead
\code{preflight} &
Before training: data, periods, code, active model parameters,
and numerical tolerances are filled in; applicable checks of data, canonicalization,
eligibility, and mathematical preparation properties are completed before computation.
Trained models, planned precision parameters for evaluation on $D$,
and future load results are not required. \\
\code{post_training} &
Trained models and their hashes are saved; applicable
mathematical checks of training and data separation are completed.
The real A/B protocol and B18 are not mandatory at this stage. \\
\code{post_load} &
Separate service acceptance, not a prerequisite for local historical evaluation.
Three load runs, memory, latency, failures,
and replay and recovery checks are saved. The test result
may be negative; this prohibits acceptance but not recording the result. \\
\code{pre_selection} &
Before opening $C_{\mathrm{select}}$, the variants, budget,
and selection rule are fixed; training results, local cost
measurements, and applicable checks of variants participating in selection are available.
The plan includes the configuration of the variant without metadata.
Service load acceptance is not required.
Further development based on $C_{\mathrm{select}}$ is prohibited. \\
\code{pre_D} &
The primary policy or pair and the full evaluation plan are fixed:
half-width in proportions, planned variance, assumptions and dependence multiplier,
provenance of the pilot on $C_{\mathrm{smoke}}$, calendar grid, and blocks;
selection conditions and applicable evaluation checks are satisfied, including
the complete variant without metadata.
Ranker errors and historical protocol violations block computation;
service readiness is not required for \code{ideal_deterministic}.
Results from $D$ do not yet exist. \\
\code{post_D} &
Final estimates, intervals, and states from
Section~\ref{sec:acceptance} are saved, including negative and insufficient
results. The plan is not changed retrospectively. \\
\code{pre_AB} &
Before variant assignment, the experiment protocol,
assignment key, sample size and duration calculation, SRM, emergency conditions,
responsible roles, trusted keys, and production limits are filled in.
The A/B plan is saved immutably; B18, serving, and applicable safety
checks have passed, and readiness~\eqref{eq:acceptance-pilot} is satisfied. \\
\code{post_AB} &
After outcomes mature, results by assigned
variant, integrity checks, and stopping information are saved. \\
\bottomrule
\end{longtable}

Deadlines for applicable tests are specified in~\ref{app:tests-normative}.
After selection, the final model is frozen
and its link to the selection result is checked.
Before each preparation, training, or evaluation computation, check the
relevant data, mathematical properties, and immutability of the applicable plan.
Saved checks are repeated when affected code or inputs change.
For local evaluation, \code{post_load} is not a prerequisite for
\code{pre_selection} or \code{pre_D}.
A missing required report
produces \code{not_run}; a test that runs and detects a violation produces \code{failed}.
Data suitability, evaluation precision, and model quality are tracked
separately under Section~\ref{sec:acceptance}.

\subsection{Preliminary calendar plan}

Preliminary time shares for $A_G,A_N,B,C,E,D$ are
$20\%$, $20\%$, $25\%$, $15\%$, $10\%$, $10\%$.
After applying calendar eligibility, publish the actual sizes
$N_J^{\mathrm{obs}}$ and $N_J^+$; $E$ remains reserved.
Within $C$, equal calendar halves are initially allocated to
$C_{\mathrm{smoke}}$ and $C_{\mathrm{select}}$; eligibility with $H+g$
is applied to each separately. The shares are only an initial planning grid.
If they do not yield eligible sessions or the required number of temporal
blocks, the plan is changed before training using the early stream, or the
research status is retained. Insufficient duration cannot be remedied
by dividing a dependent period into artificially short independent blocks.
A period shorter than $H+g$ is unsuitable. Boundaries and the rule for additional
collection are set before training using preceding data; purchasing sessions
are not selectively added, and results from $D$ are not used for this purpose.


\subsection{Parameter sources}\label{app:parameter-roles}

The first run uses the recipe in~\ref{app:first-config} without tuning.
In subsequent comparisons, hyperparameters are selected only on
$C_{\mathrm{select}}$; model parameters are trained on the assigned period.
The table identifies where each is defined without repeating its values.

\begin{longtable}{L{0.34\textwidth}L{0.57\textwidth}}
\toprule Group & Where defined and fixed \\\midrule\endhead
Problem setup: $H,g,K,Q_{\max}$ &
Request rules~\ref{sec:requests}; fixed before session selection. \\
Graph and history: $\Delta_G,B_G,\tau_G,\tau_H,n_G,d_G,\lambda_G$,
$L_H,L_{\mathrm{hist}}$ &
Section~\ref{sec:model-graph}; statistics from $A_G$ only.
View limits and timestamp-group limits differ. \\
Sources: $M,m_a,\Delta_S,\eta_P$ &
Recipe~\ref{app:first-config}; comparisons preserve the overall budget.
Quotas do not replace filling after deduplication. \\
Ranker &
Section~\ref{sec:model-ranker-implementation}; training on $B$.
Penalty $\lambda_D=0$ in the baseline. \\
Precision and intervals &
Sections~\ref{sec:eval-precision}, \ref{sec:eval-time-block}.
The plan and assumptions are fixed before opening $D$ using the permitted pilot
on $C_{\mathrm{smoke}}$; they are not required for \code{preflight} or training. \\
Serving and numerical tolerances &
Sections~\ref{sec:ops-load-profile}, \ref{sec:ops-numerics}.
Thresholds are fixed before model selection. \\
Neural network: dimensions, paths, temperatures, loss weights &
Appendix~\ref{app:neural-config}; training on $A_N$ with
the graph, catalog, and statistics of $A_G$ unchanged.
Optimization rules are in~\ref{sec:modern-training-recipe}. \\
Search mixing: $\bar g_N,g_0,p_S$ &
$0\le\bar g_N\le1$; with a positive bound,
$0<g_0<\bar g_N$. Zero direct mixing does not disable the influence of search
on intent weights. Initial values are in~\ref{app:neural-config}. \\
Sampled denominator: $m,\delta$ &
Appendix~\ref{sec:sampled-training}; a separate training procedure,
with the proposal built from $A_G$. \\
Sequence and aging &
Appendices~\ref{sec:sequence-control}, \ref{sec:freshness};
the history window and aging do not replace within-session decay. \\
SID: dimensions, codebooks, schedules, losses &
Appendix~\ref{app:sid-config}; training on $A_N$ with
the neural network fixed. In the trainable variant,
$\lambda_R^{\mathrm{SID}}>0$; the other penalties are nonnegative. \\
Simple mixtures and rare-edge penalty &
Section~\ref{sec:model-extensions}; separate comparisons.
The baseline does not use $\omega_r,\gamma_m,\lambda_{\mathrm{edge}}$. \\
Impression experiment &
Section~\ref{sec:eval-online}; $\alpha$, power, the business threshold
$\delta_{\mathrm{accept}}$, and the alternative
$\delta_{\mathrm{plan}}>\delta_{\mathrm{accept}}$
are fixed before assigning variants. \\
\bottomrule
\end{longtable}

\subsection{Period sizes and data sufficiency}

A sample of 200--500 thousand sessions may be sufficient for an engineering
prototype, but the number of sessions with an observed target is unknown in advance.
When a period with targets is opened, record
\begin{equation}\label{eq:app-positive-counts}
N_J^{\mathrm{obs}},\qquad N_J^+=|\mathcal S_J^+|,\qquad
N_J^{\mathrm{pair}}=\sum_{s\in\mathcal S_J^+}\ind\{|\mathcal J_{q_s^{(1)}}|>0\}.
\end{equation}
The first two numbers count sessions; the third counts requests with training
pairs. Candidate rows are not independent purchases.
For $A_N$, also report the number of requests with $Y_q\cap\mathcal I_G\ne\varnothing$,
the number of usable contrastive pairs, and the fraction of targets outside the catalog.
Data sufficiency is checked for every stage.

Coverage precision conditions are specified
in~\ref{sec:eval-precision} and~\ref{sec:eval-information}.
Additional collection is allowed only under a prior plan,
not until a convenient purchase count is reached.

\clearpage
\section{Comparison order and isolating component contributions}\label{app:ablation}

After the run in~\ref{app:first-config}, select the next comparison
under Section~\ref{sec:eval-next-step}. Before starting, fix the question,
data, and budget. The tables below list permissible studies;
running every row is not required. P and R are mandatory descriptive
controls in the first evaluation; G--GCS provide diagnostics on $C_{\mathrm{smoke}}$.
The main control for improvement studies is GCS+Rk.

\begin{longtable}{L{0.12\textwidth}L{0.40\textwidth}L{0.38\textwidth}}
\toprule Variant & Components & Question tested \\\midrule\endhead
P & Historical popularity & Does the system outperform an available
nonpersonalized response? \\
R & Repeat views and popularity & Is the result explained
by a simple reminder of viewed items? \\
G & Graph relations and fallback sources & Are item-to-item transitions
useful beyond repeat views? \\
GC & Graph and item features & Does content similarity improve coverage
under the same candidate limit? \\
GCS & Graph, item features, and search query & Does available search
context provide additional information? \\
GCS+Rk & The same sources and a trained ranker & Does ranking
reduce the gap between candidate coverage and final-output coverage? \\
\bottomrule
\end{longtable}

The second group of comparisons isolates the contributions of parts of the trainable module.
Baseline sources are retained unless a row states otherwise.
\begin{longtable}{L{0.25\textwidth}L{0.65\textwidth}}
\toprule Variant & Change relative to the full trainable model \\\midrule\endhead
Without graph branch & Remove $e_i^G$; retain the other branches and training
on $A_N$. \\
Direct relations only & $H_P=1$; test the contribution of indirect paths
with the same number of retained neighbors. \\
One intent & $J_I=1$; the total quota of intent sources
is assigned to the single direction. \\
Without item content features & Remove branches $e_i^T,e_i^V$ of the new
encoder; retain identifiers, graph, and available search context. \\
Without search in the encoder & Set mask $M_q^S$ to zero throughout the new
module, including the search-only fallback representation;
the baseline search source is unchanged. \\
Without direct search mixing & Only $\bar g_N=0$; retain the influence of search
on intent weights and all other rules. \\
Without contrastive component & $\lambda_C=0$; retain the other targets and available
modalities. \\
Without next view & $\lambda_{\mathrm{next}}=0$; retain
purchase training and the selected regularizers. \\
Full model & All enabled components from Section~\ref{sec:modern-model};
the main architecture's quotas from the initial plan. \\
Sequential control & Replace the session encoder with one
causal attention block while retaining the graph item
encoder; the same inputs, purchase target, and candidate budget. \\
Sampled denominator & The approximate loss from
Appendix~\ref{sec:sampled-training}; the exact version is a reference
at a feasible scale. \\
Three SID variants & Random fixed codes, fixed codes
after reconstruction, and jointly trained codes; quotas of the architecture with SID. \\
\bottomrule
\end{longtable}

In a single-component comparison, retain $A_N$, the other targets,
and the training schedule. Compare the sequential encoder
with the main model using $\lambda_C=0$ in both. When changing $J_I$,
retain the total quota and $M$. The static-vector assumption
is tested by separately removing all supplied vectors,
including baseline content sources and search: this compares
available information, not a single branch.

When comparing retrieval, fix the ranker algorithm and tuning budget,
but retrain it on each variant's candidates in $B$.
Compare the candidate bound and the output.
When comparing only rankers, fix
the same candidates; select simple-mixture coefficients on $C_{\mathrm{select}}$.
Intents are checked under Section~\ref{sec:modern-intent-diagnostics};
high weight entropy does not prove distinct directions,
and the diversity penalty is initially disabled.

Build the graph, catalog, and statistics from the same $A_G$; note additional
training on $A_N$ in the report. Compare the contribution of architecture
using identical training data. Equal $M$ does not mean equal
cost: before starting, limit the number of configurations, time, and memory;
report hardware, completed runs, and reasons for failures.

The next stage, separate from the first run, permits at most
12 complete configurations in total. The initial plan contains nine:
GCS+Rk and neural variants $J_I=1,2$, each with budget
\code{num_boost_round}$\in\{100,200,400\}$; the actual $N_T$
does not exceed the selected budget.
Neural variants are trained for three epochs using the same procedure;
the random seed is 17. Additions are registered before evaluation.
Seeds 29 and 43 are reserved for a separate
stability study; the best run is not selected on $D$.

The report format is specified below. An NDCG improvement does not replace the primary criterion
or operational constraints. Before opening $D$, fix one
primary pair: the model selected on $C_{\mathrm{select}}$ and GCS+Rk. Other comparisons
are exploratory; confirming several superiority claims requires
a prespecified family of tests and a multiplicity correction.
Selecting an improvement based on $D$ requires a new final sample.
The business effect is evaluated in a separate experiment with real impressions.

\subsection{Required metadata sensitivity analysis}
\label{sec:metadata-sensitivity}

The main comparison using supplied vectors is conditional on
their unknown construction date. The required variant
\code{graph_recent_popularity_no_provided_metadata} excludes all
supplied text, visual, and search vectors, category,
price, their masks, and derived features from retrieval and the ranker.
Search clicks/lists and query context are also unused.
The remaining inputs are IDs, available events, the graph, repeat views, and popularity
statistics built exclusively on $A_G$; the neural network and SID
are disabled. Catalog $\mathcal I_G$, sessions, the request schedule,
eligibility rules, $K$, and $M$ are retained so the denominator does not change.

The three graph relations each have a quota of 50, repeat views have
20, and the initial popularity quota is 10; after deduplication,
popularity fills the remaining positions up to $M=300$. A new ranker
is trained on its own candidates from $B$ with the same budget and callback,
using a feature schema without metadata fixed in advance; $\lambda_D=0$.
The configuration and its hash are fixed before $C_{\mathrm{select}}$.
The complete implementation and applicable checks of this variant are required
before its computations and by final evaluation, but do not block the first $C_{\mathrm{smoke}}$.
For $D$, the published results must include the chain $V,U,\operatorname{Hit}$,
the difference from the baseline, temporal sensitivity intervals, and cost.

This is a required check of the result's dependence on information,
not an automatically confirmatory second hypothesis. There remains a single primary
pair; if several superiority claims are made, the family and multiplicity correction
are fixed before $D$. Stability in this
check does not prove the absence of leakage in the unknown snapshot;
instability limits the applicability of the original result.
Replacing the supplied metadata while events remain unchanged must not
change the candidates, features, or response of the variant without metadata.

\subsection{Required report for the first run and comparisons}\label{app:diagnostic-output}

The table defines report fields, not results on Coveo. For the first run,
it is filled in for $C_{\mathrm{smoke}}$; for $D$, only after the protocol is fixed. Unavailable
quantities receive a state with a reason, not a substituted number.
For local \code{ideal_deterministic}, cost means measured
local time and memory; service load metrics remain
unmeasured until \code{post_load} and do not block the historical report.
Protocol \code{observed_purchases_v1} is stated explicitly.
All quality measures refer to $\mathcal S_J^+$; costs
also account for sessions without an observed target. Telemetry completeness
remains unknown; conversion across all sessions is not estimated.

\begin{longtable}{L{0.25\textwidth}L{0.65\textwidth}}
\toprule Block & Required quantities \\\midrule\endhead
Data and version & Period, policy version, $N_J^{\mathrm{obs}}$, $N_J^+$,
number of sessions without an observed target, and exclusion reasons. \\
Coverage chain & $\sum_s V_s$, $\sum_s U_s$, $\sum_s\operatorname{Hit}_s$;
the three losses~\eqref{eq:eval-loss-decomposition},
conditional coverage~\eqref{eq:eval-candidate-conditional}, and $\sum_s V_s/N_J^+$.
A proportion is undefined when its denominator is zero. \\
Final output & $\operatorname{Hit@20}_J$,
$\operatorname{Coverage}_{\mathrm{cand},J}$; paired difference from the baseline
and its confidence interval when comparing models. \\
Sources & Availability, list sizes, overlap, unique coverage;
losses from quota limits and $\Delta_a^{\mathrm{fixed}}$
from~\eqref{eq:eval-source-removal}. Removal with fallback filling
and retraining is performed only as a separate selected experiment
(Section~\ref{sec:eval-source-contribution}). \\
Request conditions & Session counts and coverage in prespecified groups:
search, cart, modality availability, item history, and session history.
Purchases from the observed cart and outside it are reported separately. \\
Cost & Training time, accelerator hours, peak RAM
and GPU memory; hardware, load, P95, P99, error,
deadline-exceedance, and fallback-response rates; agreement of approximate and exact search. \\
Next action & Reason for selecting the next comparison
under Section~\ref{sec:eval-next-step}, data to be fixed, and budget. \\
\bottomrule
\end{longtable}

Save the short source lists before merging, item provenance,
the output, and reasons for fallback responses. Quota losses are estimated
from target ranks or with diagnostic lists of a depth bounded
in advance; unknown ranks are not reconstructed by assumption.
A technical failure does not form a fourth term in the decomposition:
it can affect several stages, while fallback may preserve a hit.

\section{Numerical example}\label{app:example}

Consider three eligible items $a,b,c$. Normalizing a row
of one graph relation gives $P_{ab}=0.8$, $P_{ac}=0.2$;
the total row weight is $d_a=20$, and $\lambda_G=10$.
The history contains only a view of $a$, so
\begin{equation}\label{eq:app-graph-example}
\rho_a=\frac23,\qquad \alpha_{q1}=1,\qquad
g_{qb}=\frac{8}{15},\qquad g_{qc}=\frac{2}{15}.
\end{equation}
These numbers are not purchase probabilities.
Let the normalized text vectors and query be
\begin{equation}\label{eq:app-content-example}
\bar d_a=(1,0),\quad \bar d_b=(0.8,0.6),\quad
\bar d_c=(0,1),\quad\bar z_q=(0,1).
\end{equation}
Text similarity to the viewed item is $0.8$ for $b$
and $0$ for $c$, while search similarity is $0.6$ for $b$ and $1$ for $c$.
The graph and history favor $b$; search favors $c$.
This is a meaningful disagreement between features, resolved by the learned
ranker, not by an unspecified addition of different scales.

If the recorded future purchase is $c$ and the candidates are $\{b,c\}$, then
$y_{qb}=0$, $y_{qc}=1$. With equal initial scores,
$\beta_R=1$, and pair weight $\omega=1$, the derivatives are
\begin{equation}\label{eq:app-gradient-example}
g^R_{qc}=-\frac12,\quad g^R_{qb}=\frac12,
\qquad h^R_{qc}=h^R_{qb}=\frac14.
\end{equation}
A minimization step increases the score of $c$ relative to $b$.
The numbers illustrate the derivatives of one pair and do not describe
a completed model on Coveo.

Now suppose that, of 1000 eligible recorded sessions, 30 have
a recorded future purchase, and the first output covered such
an item in 12 sessions. The first request's candidates contained it in 20 sessions.
Then
\begin{equation}\label{eq:app-proxy-example}
N_J^{\mathrm{obs}}=1000,\quad N_J^+=30,\quad
\operatorname{Hit@20}_J=\frac{12}{30}=0.4,\quad
\operatorname{Coverage}_{\mathrm{cand},J}=\frac{20}{30}.
\end{equation}
If an eligible target existed in 24 sessions, the three losses as proportions are
$(30-24)/30$, $(24-20)/30$, and $(20-12)/30$. Their sum, $0.6$,
complements Hit@20 to one. The other 970 sessions do not become
confirmed negative outcomes.
Conditional retrieval coverage is $20/24$; neither the gap nor this proportion
measures additional purchases caused by recommendations.

In a separate randomized experiment, 3000 purchasing sessions out of 100 thousand in the control
and 3200 out of 100 thousand in the new variant give
\begin{equation}\label{eq:app-business-example}
\widehat\Delta_{1000}^{\mathrm{ITT}}=2,\qquad
\widehat{\operatorname{se}}_{\mathrm{ITT}}\approx0.775.
\end{equation}
The normal 95\% interval is approximately $[0.48,3.52]$.
It supports a positive effect under the experiment's
assumptions, but does not prove that the threshold of one additional
purchasing session per thousand is exceeded: the lower bound is below one.
All numbers in this appendix are invented to check the formulas.

\section{Testable implementation properties}\label{app:tests}

Active modules are checked. Changes to data, features, the model, or
output rules require rechecking the affected properties.

\subsection{Requirement mapping and check order}\label{app:tests-normative}

Each requirement below has a permanent ID mapped to a standard
test through its name or marker. Test code defines the input or the recipe for
constructing it, an independent reference, tolerances, and the initial generator state.
The result retains the actual report, code version, configuration, and seed;
separate preregistration or a hashed manifest for every
test is not mandatory. The deadline follows applicability under~\ref{app:run-stages}.
The absence of a reference or report does not count as a passed check.

Counters, sets, masks, versions, and order are compared exactly.
Numbers must satisfy
$|x-x_{\mathrm{ref}}|\le\mathrm{atol}+\mathrm{rtol}|x_{\mathrm{ref}}|$;
tolerances are specified in~\ref{sec:ops-numerics}, and for tree derivatives in
\ref{sec:model-ranker-implementation}. An unexpected NaN/Inf is
a failure; exceptions, such as latency $+\infty$, are listed explicitly.
Generative checks contain at least 1000 cases per property,
with initial generator seed 17. Save its algorithm, version,
and minimal counterexample. The numerical examples below are mandatory.

Deadlines: B01, B02, B06, B09, B13, and \code{B10.data} are due at \code{preflight};
B04, B05, B07, B08, and \code{B10.ranker} are due before applying the respective
computations, with actual artifacts checked at \code{post_training}.
B14--B17 and \code{B03.service} (delivery and state) belong
to \code{post_load}; \code{B03.data} (schema and canonicalization) belongs to \code{preflight}.
Applicable replay properties for local evaluation are checked
before computation, without mandatory implementation of service delivery and recovery.
B11 checks that calibration is disabled before the first output calculation.
\code{B10.metadata} is mandatory for the complete variant
without metadata before its computations and before $D$; the variant configuration
is fixed before $C_{\mathrm{select}}$.
B12 is performed before computing the reported temporal intervals
and on actual inputs before $D$.
For sampled training, \code{X02.support} is performed at
\code{preflight}, and \code{X02.loss} at \code{post_training}.
Other mathematical checks of enabled research components
are performed before the respective computations are applied and confirmed
at \code{post_training}. Applicable B10--B12 checks are repeated on actual
artifacts before \code{pre_D}. B18 and the real experiment protocol
are required by \code{pre_AB}, but not by \code{post_training}.
Results on $D$ and actual A/B outcomes
are recorded only at \code{post_D} and \code{post_AB}, respectively.
Tests of an enabled ANN are completed before it is admitted to serving.

\subsection{Baseline system}\label{app:tests-base}

\begin{testcase}[B01. Time and canonicalization]
On the millisecond grid, check events exactly at $t_q$ and after it,
a purchase in the first group and at $t_{0,s}+H$, a start at the eligibility boundary,
and the recorded-history boundary $t_{\mathrm{last},s}^{\mathrm{obs}}+g=b_{A_G}$.
Equality excludes the session from $A_G$; the calendar margin does not establish
telemetry completeness.
The reference is the inequalities in~\ref{sec:time} and~\ref{sec:requests}, without rounding.
Permuting simultaneous rows preserves the result.
Two identical rows with distinct IDs have raw multiplicity 2:
\code{set_v1} gives 1, \code{multiset_v1} gives 2; repeating the same ID
does not increase multiplicity. Future events do not change past inputs.
For the last group with nine views of $A$ and one of $B$
(distinct event IDs) and candidate similarities of $1$ and $0$,
the mean~\eqref{eq:model-last-group-mean} is $0.5$ under \code{set_v1}
and $0.9$ under \code{multiset_v1}. Without vector $B$, both means equal $1$;
without either vector, the value and mask are $0$.
Changing multiplicity semantics preserves the set of distinct items
in the group and the maxima, but may change retrieval by averaged scores.
\end{testcase}

\begin{testcase}[B02. Search ambiguity]
Two vectors $(1,0)$ and $(0,1)$ with the same $(s,t)$ remain
ambiguous after canonicalization and any row permutation.
Two identical usable vectors do not create false ambiguity.
An unusable or missing vector in the latest search group
disables context even if that group contains one usable vector.
Replacing lists \code{clicked_skus_hash}, \code{product_skus_hash}
does not change the inputs or response of the main variant.
\end{testcase}

\begin{testcase}[B03. Retries, closed blocks, and limits]
Accept 100 events and repeat the same block after 1 ms:
insufficient tokens do not change counters, completeness, or the previous response.
The same holds for a verifiable retry after $\delta_{\mathrm{late}}$.
With $\delta_{\mathrm{future}}=0$, a new block at $t=a_{\mathrm{block}}$
passes the time check; a lead of 1 ms or one hour produces
\code{timestamp_future} without changing the history or $t_{\mathrm{last},s}$.
Check both calendar-range boundaries, incorrect units,
overflow, and replay of an accepted block after the allowed range changes.
Check a changed payload under the same ID, an ID list of the same length
with different contents, a different owner,
an empty block, an incorrect event count, and an incomplete envelope.
An envelope consisting of two parts of 60 events is not accepted in parts;
new IDs for an already closed time produce \code{closed_block_conflict}.
With 950 accepted events, a new block of 100 is rejected in full.
The accepted history is preserved; a gap in the new trusted portion is reflected
under~\ref{sec:ops-transport-block}. Untrusted input does not corrupt another session.
Two \code{remove} actions do not produce a negative quantity; simultaneous
\code{add}/\code{remove} actions without an order leave the cart unknown.
Only \code{add} yields the corresponding state; only \code{remove}
yields the removed state; the absence of both leaves the state unknown.
\end{testcase}

\begin{testcase}[B04. Graph]
One hundred repetitions of a pair in one session yield support 1 and weight at most 1.
A nonempty normalized row sums to 1; an empty row remains zero.
For fixed $W_{ij}^r>0$, increasing the other row weights
decreases $\rho_i^rP_{ij}^r=W_{ij}^r/(d_i^r+\lambda_G)$.
A separate edge penalty does not replace the reference path matrix.
\end{testcase}

\begin{testcase}[B05. Ranker numerics and derivatives]
A missing, zero, nonfinite, or incorrectly dimensioned external
vector receives a zero mask and a separate reason; valid branches remain operational.
A nonfinite internal feature or loss stops training
and activates fallback during serving. Check opposite
vectors, coordinates on the order of $10^{200}$, empty masks, and scores
$-10^4,-100,0,100,10^4$. The reference uses a scaled norm
and stable softplus/LSE. Original modality spaces are not mixed.
For an enabled neural network, also check entropy and KL at saturation
under N06; finite softmax values without finite derivatives are insufficient.

For pair $(1,0)$ with score differences $-5,-0.1,0,0.1,5$,
compare derivatives with central differences, step $10^{-4}$,
at fixed ranks and weights: absolute tolerance $10^{-7}$
for the gradient and $10^{-6}$ for the Hessian. Check the direct sum of pairs,
FP32 casting, a zero group, and a zero sum of pairwise gradients.
On a small tree with $\gamma_T=0$, check the leaf value and
$\operatorname{Gain}_{\mathrm{LGBM}}=2\operatorname{Gain}_{\mathrm{math}}$;
external Dataset weights are disabled.
A separate example contains at least five categories: one with $G=H=0$,
and the others with informative pairs and a positive Hessian sum.
The check must exercise the grouped categorical split branch;
with \code{cat_smooth=1}, internal keys and the result are finite.
\end{testcase}

\begin{testcase}[B06. Observed targets and eligibility]
Three eligible sessions have
unknown coverage and observed targets $\{a,b\}$, $\{z\}$,
$\varnothing$, where $z\notin\mathcal I_G$.
The first output contains $a$ only for the first session:
$N_J^{\mathrm{obs}}=3$, $N_J^+=2$,
$\operatorname{Hit@20}_J=1/2$.
Purchase $z$ remains among the targets and counts as a miss;
the future catalog is not substituted.
An empty target does not become a confirmed negative outcome;
a positive target does not establish telemetry completeness.
A repeated purchase of the same SKU does not change the set; subsequent outputs
do not change the metric. A purchase in the first-view group makes the session
ineligible. For an eligible session, a purchase at $t_{0,s}+H$
is included in the target, while a later purchase is excluded. At $t_{0,s}+H+g=b_J$, the session is excluded
regardless of purchase presence. Skipping published rows of a selected
session is a read error, not unknown telemetry.
\end{testcase}

\begin{testcase}[B07. Candidates, delivery, and loss decomposition]
Baseline quotas sum to 300; after filtering and filling,
$L_q\subseteq\mathcal C_q^{\mathrm{served}}\subseteq\mathcal I_q$,
$|\mathcal C_q^{\mathrm{served}}|\le M$.
A small catalog is not padded with ineligible items; the shared limit $K$
does not require equal actual response lengths.
A purchase outside the candidates is not inserted during training or evaluation.

Check the normal branch, fallback, failure before selection, and delivery failure
after selection. In the last case, the list is empty and the selected candidates
are preserved; before selection, both sets are empty. Branches are not merged retrospectively.
Four purchasing sessions with one loss at each of the three stages and one
hit yield $(N_J^+,\sum V_s,\sum U_s,\sum\operatorname{Hit}_s)
=(4,3,2,1)$, with three losses of $1/4$ each.
A repeated hit does not add one. Always
$\operatorname{Hit}_s\le U_s\le V_s\le1$;
if $\sum V_s=0$, conditional coverage is undefined.
\end{testcase}

\begin{testcase}[B08. Sources and fallback]
A target with provenance $\{a,b\}$ remains after removing $a$:
$\Delta_a^{\mathrm{fixed}}=0$. Removing its only source
loses the target; fallback filling may restore it.
These two calculations are recorded separately. A failure does not remove the session
from the denominator; the response actually delivered is used.
\end{testcase}

\begin{testcase}[B09. Unavailable purchase rate]
The historical purchase-probability feature is absent from the feature schema;
unknown outcomes are not replaced with zeros to calculate it.
\end{testcase}

\begin{testcase}[B10. Data separation and selection]
Parts \code{B10.data}, \code{B10.ranker}, and \code{B10.metadata}
check the data, ranker, and variant without metadata by the applicable deadlines.
Changing $D$ does not change the model, variant list, block length,
or retrieval rules; changing reserved $E$ does not change the computations.
Check that
the seven parts share no sessions, further development on $C_{\mathrm{select}}$ is prohibited,
and tie-breaking order is respected. Arbitrary replacement of supplied
metadata does not change the variant in~\ref{sec:metadata-sensitivity}.
\end{testcase}

\begin{testcase}[B11. Disabled calibration]
The run manifest and computations require no calibrator; data from $E$ are not read,
and output scores are not described as purchase probabilities.
Penalty $\lambda_D=0$; the probabilistic diversity procedure is disabled.
\end{testcase}

\begin{testcase}[B12. Temporal intervals]
For $J=4,m=2$, starts $(1,3)$ produce indices $(1,2,3,4)$;
both policies are moved together, and the statistic is a ratio of sums.
Empty calendar cells are retained; the $H+g$ calendar margin is not added.
A zero denominator and too few blocks produce \code{insufficient}.
The same state is required when all $d_s=0$, all Hit=0, or all Hit=1,
even if the formal interval has zero width.
Check the boundaries of 99/100 sessions and 19/20 blocks
under~\eqref{eq:eval-information-gate}; sufficient counts do not permit
a degenerate interval, for example with identically balanced cells.
The single-coverage formula matches replacement of $d_s$ by Hit.
The resample denominator equals
the number of sessions with an observed target; there is no factor of 1000;
cells without such sessions remain on the calendar grid.

For $B_0$ independent equiprobable signs $\pm1$,
each repeated $m$ times,
the variance of the mean is $1/B_0$, while an independence-based calculation incorrectly gives
$1/(mB_0)$. The coverage check fixes the number of repetitions
and the binomial tolerance for simulation error in advance. It is not required that
every random interval be wider than the interval for independent observations.
\end{testcase}

\begin{testcase}[B13. Run manifest stages]
When prerequisites are filled in, missing future
training/load results do not block \code{preflight}.
Local \code{ideal_deterministic} permits selection and evaluation on $D$
without \code{post_load}; the absence of a service does not establish pilot readiness.
The first $C_{\mathrm{smoke}}$ does not require B18 or a completed variant
without metadata. Omission of these checks by their
required deadlines is rejected. Results are saved separately;
changes to an already fixed plan or a mismatch between its version
and the code, inputs, and settings in use are detected before computation.
At the deadline for a test, a missing result produces \code{not_run}.
A negative report does not become \code{passed} merely because it exists.
Check conditional requirements for fields of enabled modules.
The evaluation protocol is explicit; an unknown mode is rejected.
Disabled calibration does not block \code{pre_D}, but a successful
historical report does not establish complete outcome recording for A/B testing.
An incorrect type, unknown state, inconsistent times,
violated input integrity, or missed applicable check deadlines are rejected.
\end{testcase}

\begin{testcase}[B14. Versions and replay]
Batch and sequential processing of the same visible prefix
produce identical masks, categories, candidates, and versions; numbers agree
within the declared tolerance. A late event does not change a past response.
For $t_q=10{:}00$, $a_q=10{:}05$, a package published at 10:03 is prohibited.
A trace of a disabled image source reproduces the same
mask and $\mathcal C_q^{\mathrm{served}}$; without a trace, a healthy replay
is not declared to be the measured $U_s$. Deleted inputs are not replaced by their hash.
\end{testcase}

\begin{testcase}[B15. Load, memory, and recovery]
Failed, overdue, and undelivered requests remain in the load
and denominators. Input acceptance and queueing count toward the deadline; stage P99
values are not added. A timely fallback response remains a fallback.
Exact search peaking at 30 GiB fails the overall 28 GiB limit;
both packages are counted during an update.
Check queue overflow, a Sybil/SKU burst, signatures, checkpoint
loss, concurrent writes, and catalog failure.
Insert more than $L_{\mathrm{hist}}$ search blocks between views:
the queue of $L_H$ views is preserved. Recovery limits,
feature agreement, and RTO/RPO are checked by measurement.
\end{testcase}

\begin{testcase}[B16. Stopping and alert rates]
At one request per second, 20 consecutive technical failures
of valid requests trigger protection, even though five minutes contain fewer than 1000 requests.
Three failed internal checks also trigger protection.
Invalid requests and an unauthentic new package do not roll back a healthy version.
For $A=(2,0.5,2,0.5)$ and $h=1$, the exceedance fraction is $1/2$,
but the rule requiring two adjacent windows never triggers.
An empty window does not yield a zero error rate; a missing status counts as a failure.
\end{testcase}

\begin{testcase}[B17. Approximate search]
When ANN is enabled, compare it with exact search over the same eligible
catalog. Dot products of $0.05$ and $0.045$ yield
a loss of $0.005$ regardless of the ranking-score temperature.
An empty approximate list with a nonempty exact list fails acceptance.
Duplicates, ineligible items, and incomplete repeatability are detected separately.
\end{testcase}

\begin{testcase}[B18. Causal outcome accounting]\label{test:eval-effect}
Two sessions of one stable key receive one variant and contribute one
entry session to the primary denominator. A service error does not remove it
or a reliably recorded purchase. Check SRM and descriptive status after
a safety stop; an empty group produces \code{insufficient}.
A plan with a stratum based on post-assignment clicks is rejected before launch;
a plan based on features available before assignment checks the missing-value rule,
immutable membership, and retention of all assigned units in the denominator.
With independent $Z_1,Z_2\sim\mathrm{Bernoulli}(1/2)$, the standard plan
$S_1=1$, $S_2=Z_1$ is also rejected: when the stratum size is 1, the number of
variant-1 assignments is identically zero, rather than following $\mathrm{Bin}(1,1/2)$.
For $k=0$ and $k=n$, check the boundary beta intervals.
With all other conditions satisfied, $\ell_B=\delta_{\mathrm{accept}}<u_B$
produces \code{insufficient}, $\ell_B\le u_B=\delta_{\mathrm{accept}}$ produces
\code{failed}, and $\ell_B>\delta_{\mathrm{accept}}$ produces \code{passed}.
With predefined purchases independent of recommendations,
an ideal output may have high Hit@20 while retaining
$Y_s(1)=Y_s(0)$ and $\Delta_{1000}^{\mathrm{ITT}}=0$.
Describing this result as conversion growth is an error
in metric interpretation.
\end{testcase}

\subsection{Neural extension}\label{app:tests-neural}

\begin{testcase}[N01. Graph reference and paths]
Changing $A_N$ with $A_G$ fixed preserves the catalog,
$\mathcal N_i^P,D_{ij}^G$, and frequencies. With $\eta_N=0$,
nonempty rows have $A_{ij}^N=\nu_{ij}^P$, and $\mathcal L_G^N=0$.
On a small graph, exhaustive enumeration of up to $H_P$ transitions matches search
over states ``vertex, number of transitions'' with sufficient $B_P$.
Truncated search returns a valid path and a truncation flag,
but does not claim an exact minimum.
\end{testcase}

\begin{testcase}[N02. Masks and zero projections]
When branches are disabled, the remaining weights are normalized; without inputs, the source is empty.
For a usable input, $Y_q^G\ne\varnothing$, and $u_q^N=0$, the example is retained:
$p_N(i\given q)=1/|\mathcal I_q|$, with loss $\log|\mathcal I_q|$.
For $M_q^S=1,m_q^{SN}=0$, search mixing is zero and history is not attenuated.
Zeroing a learned projection does not remove the example.
\end{testcase}

\begin{testcase}[N03. Intents and search]
For $J_q=J_I$, jointly permuting intents, parameters, and quotas
preserves representations, scores, and candidates.
For identical $h_{q1}=h_{q2}=(1,0)$ with equal weights, entropy is
$\log2$, even though the directions are identical. With one direction,
pairwise diagnostics have a zero mask.
For $h_{q1}=(1,0)$, $z_q^N=(0,1)$, and $g_q^{SN}=1/4$,
the result is $u_q^N=(3/4,1/4)$ and the score of item $(0,1)$ is $1/(4\tau_N)$.
With $\bar g_N=0$, search may still change the weights of distinct intents;
disabling all its paths eliminates this influence. Mask $p_S$
is not generated during serving.
\end{testcase}

\begin{testcase}[N04. Purchase and next-view targets]
Direct and blockwise computations of the full denominator agree.
Two purchases have masses of $1/2$; score derivatives sum to zero;
a nontarget score affects the denominator.
For two views at 0 and 60 without a purchase, the last view
enters the next-view target; also check $t_{0,s}+H$ and a view
simultaneous with the first purchase under~\eqref{eq:modern-next-time}.
\end{testcase}

\begin{testcase}[N05. Contrastive comparisons]
A batch contains no two requests from the same session. Identical signatures
and a shared purchase are excluded by the specified rule.
Without eligible comparisons, the loss is zero and inputs and candidates are unchanged.
Batch order and masks are preserved; variant \code{label_free_v1}
removes only the shared-purchase filter, without reading $C,E,D$.
\end{testcase}

\begin{testcase}[N06. Initialization and training]
Repeating the specified profile and generator states reproduces
initial tensors and batch order. The shared matrix is initialized
once; with $U_N=0$, initially $A^N=\nu^P$.
Compare two Adam steps on a scalar with an independent calculation:
norm clipping precedes moment updates, and L2 is counted once.
At fixed parameters, the average purchase component of the loss
over all batches, including the short final batch, matches the full
loss under~\eqref{eq:modern-batch-loss}.
For logits $(0,-200)$ in FP32, entropy is finite and close to zero,
and KL relative to $(1/2,1/2)$ is close to $\log2$. Values and gradients with respect to logits
are compared with a stable FP64 reference; a single position yields zero
entropy and gradient, while an empty mask yields the specified zero branch.
Rounding a probability to zero does not cause NaN.
Check stream independence: a disabled optional consumer
does not change batch order or other random transformations.
\end{testcase}

\begin{testcase}[N07. Quotas and baseline restoration]
The overall limit $M$ is respected for any number of intents.
Disabling trainable sources while restoring the original quotas,
features, and ranker reproduces the baseline; merely zeroing the new score is insufficient.
With an unrecorded neural-source failure, the result matches
the common fallback: the branch, $\mathcal C_q^{\mathrm{served}}$, and output order
are unambiguous; fallback-output scores receive no probabilistic interpretation.
\end{testcase}

\subsection{Additional variants}\label{app:tests-extra}

\begin{testcase}[X01. Aging, publication, and graph traversal]
Shifting an event pair into the past while keeping its interval unchanged preserves
the within-session factor and decreases the aging factor;
outside the window, the relation is excluded. The graph, paths, vectors, and index
are published as a consistent package.
With fixed $c_k$ and eligible whole sessions, events outside the window or after
$c_k$ do not change the package. Pair repetitions do not increase a session's contribution beyond 1;
on retained edges, $\rho_i^rP_{ij}^r=W_{ij}^r/(d_i^r+\lambda_G)\le P_{ij}^r$.
A package with $r_k\ge t_q$ is unavailable; an update does not change an old response or its version.
For enabled PPR on a graph without edges, check
$z^{(1)}=\eta_Rh_q+(1-\eta_R)\pi^{\mathrm{pop}}$
and measure the cost of a full step as the catalog grows.
\end{testcase}

\begin{testcase}[X02. Sampled denominator]
\emph{Support (before training in sampled mode).}
In order $a,b$, with frequencies $(1,0)$ and $\eta_P=\delta=10^{-20}$,
the computed $Q_G\approx(1,1.5\cdot10^{-20})$, but the saved $F=(1,1)$.
For target $a$, the directly computed allowed mass is positive, while $\rho_b=0$;
before training or consuming the stream, failure
\code{negative_sampling_precision_unsupported} is required, not removal of $b$
or waiting for a budget failure. Also check strictly increasing
$F=(2^{-54},3\cdot2^{-55},1)$: its second interval is positive,
but contains no generator-grid points and must be rejected.
Check the initial parameters on the actual catalog against all conditions in
\ref{sec:sampled-proposal}; the check does not consume the random stream.

\emph{Loss (after training).}
In catalog $\{a,b,c\}$, with $T_q^a=\{a\}$, zero scores,
$m=1$, $Q_q^a(b)=1/4$, $Q_q^a(c)=3/4$, we have $Z_q=3$.
The estimate takes values $5$ and $7/3$ with these probabilities:
$\E[\widehat Z_q^a]=3$, but $\E[\log\widehat Z_q^a]<\log3$.
The mean loss gradient with respect to the positive score is $-22/35$,
versus the exact $-2/3$. With a fixed sample, finite differences
reproduce~\eqref{eq:sampled-gradient}.
Positive items enter once; negative items retain
their sampling multiplicity. An empty complement yields exact mode.
Check positive support of $Q_q^a$, exclusion of all targets from
the negative sample, multiple targets, zero vectors,
and a purchase outside the catalog: it does not expand the catalog and remains a possible miss.
Compare the mean $\widehat Z$ with exact $Z$, accounting for random
sampling error and the allowed distribution discretization error;
exhaustive enumeration of the complement with unit weights reproduces the exact
loss and gradients. A finite random sample is not required to do so.
Repeating stream $j=4$ preserves sampled items and the number of rejected attempts;
a zero task and an empty complement do not change the generator state.
\end{testcase}

\begin{testcase}[X03. Group sequence]
Permutation within a group does not change the representation; the causal mask
prevents later groups from influencing earlier ones. Moving items between groups
may change the response. The control retains action types and the purchase target;
an action-enriched variant is identified separately.
For $\gamma=1,\beta=0$, input $(0.001,-0.001)$, and $\eps_{\mathrm{LN}}=10^{-5}$,
the first LN coordinate is $0.001/\sqrt{0.000011}$, approximately $0.3015113$.
\end{testcase}

\begin{testcase}[X04. SID: derivatives and stages]
With fixed noise, finite differences reproduce
\eqref{eq:sid-soft-gradient}. Check the path from $\mathcal L_R^{\mathrm{SID}}$
to $A_\ell^{\mathrm{SID}}$ through soft assignments (a nonzero gradient
on a nondegenerate example), and from
$\mathcal L_G^{\mathrm{SID}}$ to the decoder's input code vectors.
With target codes fixed, $\mathcal L_G^{\mathrm{SID}}$ gives no
derivative with respect to $A_\ell^{\mathrm{SID}}$; $e_i,u_q^N$ are fixed.
For each stage, check the set of updated parameters, the reset of Adam
state, and noise repeatability. Only the decoder changes in the final stage.
Saturated assignments retain finite balancing loss and gradients;
finite-difference checks do not apply to the discrete $\argmax$.
\end{testcase}

\begin{testcase}[X05. SID collisions]
For $L_{\mathrm{SID}}=2$, two items with code $(2,1)$
and $R_{\mathrm{SID}}=1$ yield the first item
in the specified order. A purchase of the second does not change the mapping
and remains a miss for the SID source; another source may retrieve the item.
Without noise, codes and candidates repeat; all returned items belong to $\mathcal I_q$.
The random control with the same stream $m=2,j=6$ reproduces unique codes.
Inconsistent codebooks, decoder, and tree are not admitted to serving.
\end{testcase}

\section{Contents of the final research deliverable}\label{app:deliverables}

Deliver the run manifest, checks, and report~\ref{app:diagnostic-output},
including data sources and evaluation denominators, as well as the A/B plan.
Without the corresponding observations, causal and monetary results
have state \code{not_observed}.

\begingroup
\small
\setlength{\parskip}{0pt}
\begin{thebibliography}{99}
\setlength{\itemsep}{0.3\baselineskip}
\addcontentsline{toc}{section}{References and documentation}

\bibitem{coveo-data}
Coveo. \emph{SIGIR eCom Data Challenge: dataset documentation}.
\url{https://github.com/coveooss/SIGIR-ecom-data-challenge}.
Description of the schema, anonymization, features, and search data.

\bibitem{coveo-paper}
Tagliabue J., Greco C., Roy J.-F., Bianchi F., Cassani G., Yu B., Chia P. J.
\emph{SIGIR 2021 E-Commerce Workshop Data Challenge}.
2021. \url{https://arxiv.org/abs/2104.09423}.

\bibitem{spare2025}
Peintner A., Mohammadi A. R., Zangerle E.
\emph{Efficient Session-based Recommendation with Contrastive
Graph-based Shortest Path Search}.
ACM Transactions on Recommender Systems, 2025.
\url{https://doi.org/10.1145/3701764}.
Author version:
\url{https://evazangerle.at/publication/peintner-tors-2025/}.

\bibitem{miasrec2024}
Choi M., Kim H.-Y., Cho H., Lee J.
\emph{Multi-intent-aware Session-based Recommendation}.
SIGIR, 2024. \url{https://doi.org/10.1145/3626772.3657928}.
\url{https://arxiv.org/abs/2405.00986}.

\bibitem{macl2025}
Zhang X., Xu B., Ma F., Wang Z., Yang L., Lin H.
\emph{Rethinking Contrastive Learning in Session-based Recommendation}.
Pattern Recognition. Vol.~169, 2026. Article~111924.
\url{https://doi.org/10.1016/j.patcog.2025.111924}.
\url{https://arxiv.org/abs/2506.05044}.

\bibitem{grid2025}
Ju C. M., Collins L., Neves L., Kumar B., Wang L. Y., Zhao T., Shah N.
\emph{Generative Recommendation with Semantic IDs: A Practitioner's Handbook}.
2025. \url{https://arxiv.org/abs/2507.22224}.
Authors' repository: \url{https://github.com/snap-research/GRID}.

\bibitem{diger2026}
Fu J., Ge X., Karatzoglou A., Arapakis I., Verberne S., Jose J. M., Ren Z.
\emph{Differentiable Semantic ID for Generative Recommendation}.
SIGIR, 2026. \url{https://doi.org/10.1145/3805712.3809641}.
\url{https://arxiv.org/abs/2601.19711}.

\bibitem{sasrec2018}
Kang W.-C., McAuley J.
\emph{Self-Attentive Sequential Recommendation}.
ICDM, 2018. \url{https://arxiv.org/abs/1808.09781}.

\bibitem{sampledsoftmax2025}
Chen J., Zhang J., Huang X., Yang Y., Lian D., Chen E.
\emph{Adaptive Sampled Softmax with Inverted Multi-Index:
Methods, Theory and Applications}.
2025. \url{https://arxiv.org/abs/2501.08563}.

\bibitem{logq2025}
Khrylchenko K., Baikalov V., Makeev S., Matveev A., Liamaev S.
\emph{Correcting the LogQ Correction: Revisiting Sampled Softmax
for Large-Scale Retrieval}.
RecSys, 2025. \url{https://doi.org/10.1145/3705328.3748033}.
\url{https://arxiv.org/abs/2507.09331}.

\bibitem{candidatesampling}
TensorFlow. \emph{Candidate Sampling Algorithms}.
\url{https://www.tensorflow.org/extras/candidate_sampling.pdf}.
Comparison of training methods with sampled classes and noise examples.

\bibitem{lambdamart}
Burges C. J. C.
\emph{From RankNet to LambdaRank to LambdaMART: An Overview}.
Microsoft Research Technical Report MSR-TR-2010-82, 2010.
\url{https://www.microsoft.com/en-us/research/publication/from-ranknet-to-lambdarank-to-lambdamart-an-overview/}.

\bibitem{lightgbm}
LightGBM 4.6.0. \emph{Parameters; Python training API;
custom derivatives and split gain implementation}.
\url{https://lightgbm.readthedocs.io/en/v4.6.0/Parameters.html};
\url{https://lightgbm.readthedocs.io/en/v4.6.0/pythonapi/lightgbm.train.html};
\url{https://github.com/microsoft/LightGBM/blob/v4.6.0/python-package/lightgbm/basic.py};
\url{https://github.com/microsoft/LightGBM/blob/v4.6.0/src/treelearner/feature_histogram.hpp}.
Version \code{lightgbm==4.6.0}, binary package hashes, and the callback
are saved in the run lockfile.

\bibitem{item2vec}
Barkan O., Koenigstein N.
\emph{Item2Vec: Neural Item Embedding for Collaborative Filtering}.
2016. \url{https://arxiv.org/abs/1603.04259}.

\bibitem{ips}
Schnabel T., Swaminathan A., Singh A., Chandak N., Joachims T.
\emph{Recommendations as Treatments: Debiasing Learning and Evaluation}.
2016. \url{https://arxiv.org/abs/1602.05352}.

\bibitem{dr}
Dud\'ik M., Langford J., Li L.
\emph{Doubly Robust Policy Evaluation and Learning}.
2011. \url{https://arxiv.org/abs/1103.4601}.

\bibitem{experiments}
Kohavi R., Tang D., Xu Y.
\emph{Trustworthy Online Controlled Experiments:
A Practical Guide to A/B Testing}.
Cambridge University Press, 2020.
\url{https://doi.org/10.1017/9781108653985}.

\bibitem{holm1979}
Holm S. \emph{A Simple Sequentially Rejective Multiple Test Procedure}.
Scandinavian Journal of Statistics. 1979. Vol.~6, no.~2. P.~65--70.
\url{https://www.jstor.org/stable/4615733}.

\bibitem{bootstrap}
Efron B.
\emph{Bootstrap Methods: Another Look at the Jackknife}.
The Annals of Statistics. 1979. Vol.~7, no.~1. P.~1--26.
\url{https://doi.org/10.1214/aos/1176344552}.

\bibitem{kunsch1989}
Künsch H. R.
\emph{The Jackknife and the Bootstrap for General Stationary Observations}.
The Annals of Statistics. 1989. Vol.~17, no.~3. P.~1217--1241.
\url{https://doi.org/10.1214/aos/1176347265}.

\end{thebibliography}
\endgroup
\end{document}
